%%%%%%%%%%%%%%%%%%%%%%%%%%%%%%%%%%%%%%%%%%%%%%%%%%%%%%%%%%%%%%%%%%%%%%%%%%%%%%
%%%%%%%%%%%%%%%%%%%%%%%%%%%%%%%%%%%%%%%%%%%%%%%%%%%%%%%%%%%%%%%%%%%%%%%%%%%%%%
%%%%%%%%%%%%%%%%%%%%%%%%%%%%%%%%%%%%%%%%%%%%%%%%%%%%%%%%%%%%%%%%%%%%%%%%%%%%%%

\chapter{Fungsi sebagai Limit} \label{approx:chapter}

%%%%%%%%%%%%%%%%%%%%%%%%%%%%%%%%%%%%%%%%%%%%%%%%%%%%%%%%%%%%%%%%%%%%%%%%%%%%%%

\section{Bilangan kompleks}
\label{sec:complexnums}

%mbxINTROSUBSECTION

\sectionnotes{setengah kuliah}

\subsection{Bidang kompleks}

Dalam bab ini, kita membahas aproksimasi fungsi, atau dengan kata lain,
fungsi sebagai limit barisan dan deret.
Kita akan memperluas beberapa hasil yang telah kita lihat ke suatu
latar yang sedikit lebih umum, dan akan melihat beberapa hasil yang sama sekali baru.
Secara khusus, kita membahas fungsi bernilai kompleks.
Sebelumnya kita telah memberikan bilangan kompleks sebagai contoh, tetapi
mari kita mulai dari awal dan mendefinisikan lapangan bilangan kompleks secara tepat.

Bilangan kompleks hanyalah pasangan $(x,y) \in \R^2$ yang padanya kita
mendefinisikan perkalian (lihat di bawah).
Kita menyebut himpunan ini \emph{bilangan kompleks}\index{bilangan kompleks}
dan menotasikannya dengan $\C$.
Kita mengidentifikasi $x \in \R$ dengan $(x,0) \in \C$.
Sumbu-$x$ kemudian disebut \emph{\myindex{sumbu real}} dan sumbu-$y$ disebut
\emph{\myindex{sumbu imajiner}}.
Karena $\C$ hanyalah bidang, kita juga menyebut himpunan $\C$ sebagai
\emph{\myindex{bidang kompleks}}.

Definisikan:
\begin{equation*}
(x,y) + (s,t) \coloneqq (x+s,y+t) , \qquad
(x,y) (s,t) \coloneqq (xs-yt,xt+ys) .
\end{equation*}
Dengan identifikasi di atas, kita memiliki $0 = (0,0)$ dan $1 = (1,0)$.
Kedua operasi ini menjadikan bidang tersebut sebuah lapangan (latihan).
Kita menulis bilangan kompleks $(x,y)$ sebagai $x+iy$, dengan
mendefinisikan\footnote{Perhatikan bahwa para insinyur menggunakan $j$ sebagai pengganti $i$.}
\glsadd{not:imaginary}
\begin{equation*}
i \coloneqq (0,1) .
\end{equation*}
Perhatikan bahwa $i^2 = (0,1)(0,1) = (0-1,0+0) = -1$.
Artinya, $i$ adalah solusi persamaan polinomial
\begin{equation*}
z^2+1=0 .
\end{equation*}
Mulai sekarang, kita tidak akan menggunakan notasi $(x,y)$ dan hanya menggunakan $x+iy$.
Lihat \figureref{fig:complexplane}.
\begin{myfigureht}
\myincludegraphics{complexplane}{%
Diagram sebuah bilangan pada bidang kompleks. Bilangan 1 ditandai pada
sumbu horizontal di sebelah kanan titik asal,
dan bilangan i ditandai pada sumbu vertikal di atas titik asal.
Bilangan x ditandai pada sumbu horizontal dan bilangan i y ditandai
pada sumbu vertikal. Kombinasi x ditambah i y ditandai pada titik
(x,y).}
\caption{Titik-titik $1$, $i$, $x$, $iy$, dan $x+iy$ pada bidang
kompleks.\label{fig:complexplane}}
\end{myfigureht}

Kita umumnya menggunakan $x,y,r,s,t$ untuk nilai real dan $z,w,\xi,\zeta$
untuk nilai kompleks, meskipun ini bukan aturan yang kaku. Khususnya,
$z$ sering digunakan sebagai variabel real ketiga dalam $\R^3$.

\begin{defn}
Misalkan $z= x+iy$.
Kita menyebut $x$ 
sebagai \emph{\myindex{bagian real}} dari $z$, dan 
menyebut $y$
sebagai \emph{\myindex{bagian imajiner}} dari $z$. Kita menulis
\glsadd{not:realpart}\glsadd{not:imagpart}
\begin{equation*}
\Re\, z \coloneqq x , \qquad
\Im\, z \coloneqq y .
\end{equation*}
Definisikan
\glsadd{not:conj}
\emph{\myindex{konjugat kompleks}} sebagai
\begin{equation*}
\bar{z} \coloneqq x-iy ,
\end{equation*}
dan definisikan \emph{\myindex{modulus}} sebagai
\glsadd{not:modulus}
\begin{equation*}
\sabs{z} \coloneqq \sqrt{x^2+y^2} .
\end{equation*}
\end{defn}

Modulus merupakan analog kompleks dari nilai mutlak dan
memiliki sifat-sifat serupa.
Sebagai contoh,
$\sabs{zw} = \sabs{z} \, \sabs{w}$ (latihan).
Konjugat kompleks merupakan pencerminan bidang terhadap sumbu real.
Bilangan real tepatnya adalah bilangan-bilangan yang bagian imajinernya
$y=0$. Secara khusus, bilangan-bilangan tersebut tepatnya adalah
bilangan yang memenuhi persamaan
\begin{equation*}
z = \bar{z} .
\end{equation*}

Karena $\C$ sebenarnya adalah $\R^2$, kita gunakan pada $\C$ metrik
Euklides standar pada $\R^2$.
Secara khusus,
\begin{equation*}
\sabs{z} = d(z,0) , \qquad 
\text{dan juga} \qquad 
\sabs{z-w} = d(z,w) .
\end{equation*}
Jadi, topologi pada $\C$ sama persis dengan topologi standar pada $\R^2$
dengan metrik Euklides,
dan $\sabs{z}$ sama dengan norma Euklides pada $\R^2$.
Yang penting, karena $\R^2$ merupakan ruang metrik lengkap, maka
$\C$ juga demikian.
Karena $\sabs{z}$ merupakan norma Euklides pada $\R^2$, kita memiliki
\emph{ketaksamaan segitiga}\index{ketaksamaan segitiga!bilangan kompleks}
dalam kedua bentuk:
\begin{equation*}
\sabs{z+w} \leq \sabs{z}+\sabs{w} \qquad \text{dan} \qquad
\big\lvert \sabs{z}-\sabs{w} \big\rvert \leq \sabs{z-w} .
\end{equation*}

Konjugat kompleks dan modulus ternyata berkaitan lebih erat:
\begin{equation*}
\sabs{z}^2 =
x^2+y^2 =
(x+iy)(x-iy) =
z \bar{z} .
\end{equation*}

\begin{remark}
Tidak ada urutan alami pada bilangan kompleks.
Secara khusus,
tidak ada urutan yang menjadikan bilangan kompleks sebagai lapangan terurut.
Pengurutan adalah salah satu hal yang hilang ketika kita beralih dari bilangan real
ke bilangan kompleks.
\end{remark}

\subsection{Bilangan kompleks dan limit}

Operasi aljabar pada bilangan kompleks bersifat kontinu karena konvergensi di
$\R^2$ sama dengan konvergensi pada setiap komponen, dan kita telah mengetahui
bahwa operasi aljabar real bersifat kontinu. Sebagai contoh,
tuliskan $z_n = x_n + i\,y_n$ dan
$w_n = s_n + i\,t_n$, dan misalkan
$\lim_{n\to\infty} z_n = z = x+i\,y$ serta $\lim_{n\to\infty} w_n = w = s+i\,t$.
Mari kita tunjukkan
\begin{equation*}
\lim_{n\to\infty} z_n w_n = zw .
\end{equation*}
Pertama,
\begin{equation*}
z_n w_n = (x_ns_n-y_nt_n) + i(x_nt_n+y_ns_n) .
\end{equation*}
Topologi pada $\C$ sama dengan topologi pada $\R^2$, sehingga
$x_n \to x$, $y_n \to y$, $s_n \to s$, dan $t_n \to t$.
Dengan demikian,
\begin{equation*}
\lim_{n\to\infty} (x_ns_n-y_nt_n) = xs-yt \qquad \text{dan} \qquad
\lim_{n\to\infty} (x_nt_n+y_ns_n) = xt+ys .
\end{equation*}
Karena
$(xs-yt)+i(xt+ys) = zw$,
\begin{equation*}
\lim_{n\to\infty} z_n w_n = zw .
\end{equation*}

Demikian pula, modulus dan konjugat kompleks merupakan fungsi kontinu. Kita
serahkan sisa pembuktian proposisi berikut sebagai latihan.

\begin{prop} \label{prop:continuityofcomplex}
Misalkan $\{ z_n \}_{n=1}^\infty$, $\{ w_n \}_{n=1}^\infty$ merupakan barisan bilangan kompleks yang masing-masing konvergen
ke $z$ dan $w$. Maka
\begin{enumerate}[(i)]
\item
$\displaystyle \lim_{n\to \infty} (z_n + w_n) = z + w$.
\item
$\displaystyle \lim_{n\to \infty} z_n w_n = z w$.
\item
Andaikan $w_n \neq 0$ untuk semua $n$ dan $w \neq 0$,
$\displaystyle \lim_{n\to \infty} \frac{z_n}{w_n} = \frac{z}{w}$.
\item
$\displaystyle \lim_{n\to \infty} \sabs{z_n} = \sabs{z}$.
\item
$\displaystyle \lim_{n\to \infty} \bar{z}_n = \bar{z}$.
\end{enumerate}
\end{prop}

Seperti telah kita lihat di atas, konvergensi di $\C$ sama dengan konvergensi di
$\R^2$. Secara khusus, suatu barisan di $\C$ konvergen jika dan hanya jika
bagian real dan imajinernya konvergen. Karena itu, silakan terapkan
semua yang telah Anda pelajari tentang konvergensi di $\R^2$, serta
terapkan hasil-hasil tentang bilangan real pada bagian real dan imajinernya.

Kita juga memerlukan konvergensi deret kompleks. Misalkan $\{ z_n \}_{n=1}^\infty$ merupakan
barisan bilangan kompleks. Deret
\begin{equation*}
\sum_{n=1}^\infty z_n
\end{equation*}
\emph{konvergen}\index{konvergen!deret kompleks} jika barisan jumlah parsialnya konvergen, yaitu jika
\begin{equation*}
\lim_{k\to\infty} \sum_{n=1}^k z_n \qquad \text{ada.}
\avoidbreak
\end{equation*}
Suatu deret \emph{konvergen absolut}\index{konvergen absolut!deret kompleks}
jika $\sum_{n=1}^\infty \sabs{z_n}$ konvergen.

Kita mengatakan suatu deret
bersifat \emph{Cauchy}\index{Cauchy!deret kompleks}
jika barisan jumlah parsialnya bersifat Cauchy. Dua
proposisi berikut pada dasarnya memiliki pembuktian yang sama
dengan padanannya untuk deret real, dan
kita serahkan sebagai latihan.

\begin{prop} \label{prop:cachysercomplex}
Deret kompleks $\sum_{n=1}^\infty z_n$ bersifat Cauchy jika untuk setiap $\epsilon > 0$
terdapat $M \in \N$ sedemikian sehingga untuk setiap $n \geq M$
dan setiap $k > n$, berlaku
\begin{equation*}
\abs{ \sum_{j={n+1}}^k z_j }
< \epsilon .
\end{equation*}
\end{prop}

\begin{prop} \label{prop:absconvmeansconv}
Jika suatu deret kompleks $\sum_{n=1}^\infty z_n$ konvergen absolut, maka deret tersebut konvergen.
\end{prop}

Deret $\sum_{n=1}^\infty \sabs{z_n}$ merupakan deret real. Semua
uji konvergensi (uji rasio, uji akar, dan sebagainya)\ yang membahas
konvergensi absolut bekerja pada bilangan $\sabs{z_n}$; artinya, uji-uji tersebut
sebenarnya membahas konvergensi deret bilangan real nonnegatif.
Anda dapat langsung menerapkan uji-uji ini
tanpa perlu membuktikan ulang apa pun untuk deret kompleks.

\subsection{Fungsi bernilai kompleks}

Ketika membahas fungsi bernilai kompleks
$f \colon X \to \C$, kita sering menulis
$f = u+i\,v$ untuk fungsi bernilai real $u \colon X \to \R$ dan $v \colon X \to
\R$.

Misalkan kita ingin mengintegralkan
$f \colon [a,b] \to \C$. Kita menulis
$f = u+i\,v$ untuk $u$ dan~$v$ yang bernilai real.
Kita mengatakan bahwa $f$ \emph{terintegralkan secara Riemann}\index{terintegralkan secara Riemann!
fungsi bernilai kompleks}
jika $u$ dan $v$ terintegralkan secara Riemann,
dan dalam hal ini kita mendefinisikan
\begin{equation*}
\int_a^b f \coloneqq \int_a^b u + i \int_a^b v .
\end{equation*}
Kita menggunakan definisi yang sama untuk setiap jenis integral lainnya (tak wajar,
multivariabel, dan sebagainya).

Demikian pula, ketika melakukan diferensiasi, tuliskan $f \colon [a,b] \to \C$ sebagai
$f = u+i\,v$. Dengan memandang $\C$ sebagai $\R^2$, kita mengatakan bahwa $f$ diferensiabel
jika $u$ dan $v$ diferensiabel. Untuk fungsi yang bernilai di $\R^2$, turunannya
direpresentasikan oleh suatu vektor di $\R^2$. Vektor di $\R^2$ tersebut sekarang merupakan bilangan
kompleks. Dengan kata lain,
kita menulis
\emph{turunan}\index{turunan!fungsi bernilai kompleks}
sebagai
\glsadd{not:mvder}
\begin{equation*}
f'(t) \coloneqq u'(t) + i \, v'(t) .
\end{equation*}
Operator linear yang merepresentasikan turunan adalah perkalian dengan
bilangan kompleks $f'(t)$, sehingga tidak ada yang hilang dalam identifikasi ini.


\subsection{Latihan}

\begin{exercise}
Periksa bahwa $\C$ merupakan lapangan.
\end{exercise}

\begin{exercise}
Buktikan bahwa untuk $z,w \in \C$ berlaku
$\sabs{zw} = \sabs{z} \, \sabs{w}$.
\end{exercise}

\begin{exercise}
Selesaikan pembuktian \propref{prop:continuityofcomplex}.
\end{exercise}

\begin{exercise}
Buktikan \propref{prop:cachysercomplex}.
\end{exercise}

\begin{exercise}
Buktikan \propref{prop:absconvmeansconv}.
\end{exercise}

\begin{samepage}
\begin{exercise}
Untuk $x +iy$, definisikan matriks
$\left[ \begin{smallmatrix} x & -y \\ y & x \end{smallmatrix} \right]$.
Buktikan:
\begin{enumerate}[a)]
\item
Aksi matriks ini pada vektor $(s,t)$ sama dengan mengalikan
$(x+iy)(s+it)$.
\item
Mengalikan dua matriks semacam ini sama dengan mengalikan bilangan kompleks
yang mendasarinya, lalu mencari matriks yang bersesuaian dengan hasil kali tersebut.
Dengan kata lain, lapangan $\C$ dapat diidentifikasi dengan suatu subhimpunan matriks 2-kali-2.
\item
Matriks
$\left[ \begin{smallmatrix} x & -y \\ y & x \end{smallmatrix} \right]$
memiliki nilai eigen $x+iy$ dan $x-iy$. Ingat bahwa $\lambda$ adalah
nilai eigen matriks $A$ jika $A-\lambda I$ (matriks kompleks dalam kasus kita)
tidak invertibel, yaitu,
jika baris-barisnya bergantung secara linear: satu baris merupakan kelipatan
(kompleks) dari baris lainnya.
\end{enumerate}
\end{exercise}
\end{samepage}

\begin{exercise}
Buktikan teorema Bolzano--Weierstrass untuk barisan kompleks.
Misalkan $\{ z_n \}_{n=1}^\infty$ merupakan barisan bilangan kompleks yang terbatas.
Artinya, terdapat $M$ sedemikian sehingga $\sabs{z_n} \leq M$ untuk semua $n$.
Buktikan bahwa terdapat subbarisan $\{ z_{n_k} \}_{k=1}^\infty$
yang konvergen ke suatu $z \in \C$.
\end{exercise}

\begin{exercise}
\leavevmode
\begin{enumerate}[a)]
\item
Buktikan bahwa tidak ada teorema nilai rata-rata sederhana untuk fungsi bernilai kompleks:
Temukan fungsi diferensiabel $f \colon [0,1] \to \C$ sedemikian sehingga
$f(0) = f(1) = 0$, tetapi $f'(t) \neq 0$ untuk semua $t \in [0,1]$.
\item
Namun, ada bentuk teorema nilai rata-rata yang lebih lemah sebagaimana untuk fungsi bernilai vektor.
Buktikan: Jika $f \colon [a,b] \to \C$ kontinu dan diferensiabel pada
$(a,b)$, dan untuk suatu $M$ berlaku $\babs{f'(x)} \leq M$ untuk semua $x \in (a,b)$, maka
$\babs{f(b)-f(a)} \leq M \sabs{b-a}$.
\end{enumerate}
\end{exercise}

\begin{exercise}
Buktikan bahwa tidak ada teorema nilai rata-rata sederhana untuk integral
fungsi bernilai kompleks:
Temukan fungsi kontinu $f \colon [0,1] \to \C$ sedemikian sehingga
$\int_0^1 f = 0$ tetapi $f(t) \neq 0$ untuk semua $t \in [0,1]$.
\end{exercise}


%%%%%%%%%%%%%%%%%%%%%%%%%%%%%%%%%%%%%%%%%%%%%%%%%%%%%%%%%%%%%%%%%%%%%%%%%%%%%%

\sectionnewpage
\section{Pertukaran limit}
\label{sec:swaplim}

%mbxINTROSUBSECTION

\sectionnotes{2 kuliah}

\subsection{Kekontinuan}

Mari kita kembali ke pertukaran limit dan memperluas pembahasan
\volIref{\chapterref*{vI-fs:chapter} dari jilid I}{\chapterref{fs:chapter}}.
Misalkan $\{ f_n \}_{n=1}^\infty$ suatu barisan
fungsi $f_n \colon X \to Y$ untuk suatu himpunan $X$ dan ruang metrik $Y$.
Misalkan $f \colon X \to Y$ suatu
fungsi dan andaikan untuk setiap $x \in X$ berlaku
\begin{equation*}
f(x) = \lim_{n\to \infty} f_n(x) .
\end{equation*}
Kita mengatakan barisan $\{ f_n \}_{n=1}^\infty$
\emph{\myindex{konvergen titik demi titik}}\index{konvergensi titik demi titik} ke $f$.

Untuk $Y=\C$, suatu deret fungsi
\emph{konvergen titik demi titik}\index{konvergen titik demi titik!deret kompleks}\index{konvergensi titik demi titik!deret kompleks} ke $f$ jika
untuk setiap $x \in X$, berlaku
\begin{equation*}
f(x) = \lim_{n\to \infty} \sum_{k=1}^n f_k(x) =
\sum_{k=1}^\infty f_k(x) .
\end{equation*}

\medskip

Pertanyaannya:
Jika semua $f_n$ kontinu, apakah $f$ kontinu?  Diferensiabel?
Dapat diintegralkan?  Apa turunan atau integral dari $f$?
Sebagai contoh, untuk kekontinuan limit titik demi titik dari suatu barisan
fungsi
$\{ f_n \}_{n=1}^\infty$, kita menanyakan apakah
\begin{equation*}
\lim_{x\to x_0} \lim_{n\to\infty} f_n(x)
\overset{?}{=}
\lim_{n\to\infty} \lim_{x\to x_0} f_n(x) .
\end{equation*}
Secara apriori, kita bahkan tidak tahu apakah kedua ruas tersebut ada,
apalagi apakah keduanya sama.

\begin{example}
Fungsi $f_n \colon \R \to \R$,
\begin{equation*}
f_n(x) \coloneqq \frac{1}{1+nx^2},
\end{equation*}
kontinu dan konvergen titik demi titik ke fungsi yang tidak kontinu
\begin{equation*}
f(x) \coloneqq 
\begin{cases}
1 & \text{jika } x=0, \\
0 & \text{selainnya.}
\end{cases}
\end{equation*}
\end{example}

Jadi, konvergensi titik demi titik tidak cukup untuk mempertahankan kekontinuan
(bahkan keterbatasan). Untuk itu, kita memerlukan konvergensi seragam.
Misalkan $f_n \colon X \to Y$ fungsi-fungsi. Maka
$\{f_n\}_{n=1}^\infty$ \emph{\myindex{konvergen seragam}}\index{konvergensi seragam}
ke $f$ jika
untuk setiap $\epsilon > 0$, terdapat suatu $M$ sedemikian sehingga
untuk semua $n \geq M$ dan semua $x \in X$, berlaku
\begin{equation*}
d\bigl(f_n(x),f(x)\bigr) < \epsilon .
\end{equation*}

Deret $\sum_{n=1}^\infty f_n$ dari fungsi-fungsi bernilai kompleks konvergen seragam
ke $f$ jika barisan
jumlah parsialnya konvergen seragam, yaitu, jika untuk setiap $\epsilon > 0$,
terdapat suatu $M$ sedemikian sehingga
untuk semua $n \geq M$ dan semua $x \in X$,
\begin{equation*}
\abs{\left(\sum_{k=1}^n f_k(x)\right)-f(x)} < \epsilon .
\end{equation*}

Sifat paling sederhana yang dipertahankan oleh konvergensi seragam adalah
keterbatasan. Pembuktian proposisi berikut kami serahkan sebagai
latihan. Pembuktiannya hampir identik dengan pembuktian untuk fungsi bernilai real.

\begin{prop} \label{prop:uniformconvbounded}
Misalkan $X$ suatu himpunan dan $(Y,d)$ suatu ruang metrik.
Jika $f_n \colon X \to Y$ adalah fungsi-fungsi terbatas dan konvergen seragam ke $f
\colon X \to Y$, maka $f$ terbatas.
\end{prop}

Jika $X$ suatu himpunan dan $(Y,d)$ suatu ruang metrik, maka suatu barisan $f_n \colon X
\to Y$ disebut \emph{\myindex{Cauchy seragam}} jika untuk setiap $\epsilon > 0$, terdapat suatu $M$ sedemikian sehingga
untuk semua $n, m \geq M$ dan semua $x \in X$, berlaku
\begin{equation*}
d\bigl(f_n(x),f_m(x)\bigr) < \epsilon .
\end{equation*}
Gagasan ini sama seperti untuk fungsi bernilai real.
Pembuktian proposisi berikut
pada dasarnya juga sama seperti pada kasus tersebut dan
diserahkan sebagai latihan.

\begin{prop} \label{prop:unifcauchymetric}
Misalkan $X$ suatu himpunan, $(Y,d)$ suatu ruang metrik, dan
$f_n \colon X \to Y$ fungsi-fungsi.
Jika $\{ f_n \}_{n=1}^\infty$ konvergen seragam,
maka $\{f_n\}_{n=1}^\infty$ bersifat Cauchy seragam. Sebaliknya, jika
$\{f_n\}_{n=1}^\infty$ bersifat Cauchy seragam dan $(Y,d)$ lengkap secara Cauchy,
maka $\{f_n\}_{n=1}^\infty$ konvergen seragam.
\end{prop}

Untuk $f \colon X \to \C$, kita tuliskan\glsadd{not:uniformnorm}
\begin{equation*}
\snorm{f}_X \coloneqq \sup_{x \in X} \babs{f(x)} .
\end{equation*}
Kita menyebut $\snorm{\cdot}_X$
sebagai \emph{\myindex{norma supremum}} atau \emph{\myindex{norma seragam}},
dan subskripnya menunjukkan himpunan tempat supremum diambil.
Maka suatu barisan fungsi
$f_n \colon X \to \C$ konvergen seragam ke $f \colon X \to \C$ jika dan hanya jika
\begin{equation*}
\lim_{n\to \infty} \snorm{f_n-f}_X = 0 .
\end{equation*}

Norma supremum memenuhi ketaksamaan segitiga: Untuk setiap $x \in X$,
\begin{equation*}
\babs{f(x)+g(x)} \leq
\babs{f(x)}+\babs{g(x)} \leq
\snorm{f}_X+\snorm{g}_X .
\end{equation*}
Ambil supremum pada ruas kiri untuk memperoleh
\begin{equation*}
\snorm{f+g}_X \leq
\snorm{f}_X+\snorm{g}_X .
\end{equation*}

Untuk ruang metrik kompak $X$,
norma seragam merupakan norma pada ruang vektor $C(X,\C)$.
Kami menyerahkan pembuktiannya sebagai latihan.
Meskipun kita tidak akan memerlukannya, $C(X,\C)$ sebenarnya merupakan ruang
vektor kompleks, yaitu, dalam definisi ruang vektor kita dapat mengganti
$\R$ dengan $\C$.
Konvergensi dalam ruang metrik $C(X,\C)$ adalah
konvergensi seragam.

Kita akan mempelajari beberapa jenis deret fungsi, dan
suatu uji yang berguna untuk konvergensi seragam suatu deret adalah
\emph{\myindex{uji-$M$ Weierstrass}}.

\begin{thm}[\myindex{uji-$M$ Weierstrass}] \label{thm:weiermtest}
Misalkan $X$ suatu himpunan.
Andaikan $f_n \colon X \to \C$ adalah fungsi-fungsi dan $M_n > 0$ adalah bilangan-bilangan sedemikian
sehingga
\begin{equation*}
\babs{f_n(x)}\leq M_n \quad \text{untuk semua } x \in X,
\qquad \text{dan} \qquad
\sum_{n=1}^\infty M_n
\quad \text{konvergen}.
\end{equation*}
Maka
\begin{equation*}
\sum_{n=1}^\infty f_n(x)
\quad \text{konvergen seragam}.
\end{equation*}
\end{thm}

Cara lain untuk menyatakan teorema tersebut adalah dengan mengatakan bahwa jika
$\sum_{n=1}^\infty \snorm{f_n}_X$ konvergen, maka $\sum_{n=1}^\infty f_n$ konvergen seragam.
Perhatikan bahwa kebalikan teorema ini tidak benar.
Dengan menerapkan teorema tersebut pada $\sum_{n=1}^\infty \babs{f_n(x)}$, kita
melihat bahwa deret ini juga konvergen seragam. Jadi, deret tersebut konvergen absolut
sekaligus seragam.

\begin{proof}
Andaikan $\sum_{n=1}^\infty M_n$ konvergen. Ambil $\epsilon > 0$.
Jumlah-jumlah parsial deret $\sum_{n=1}^\infty M_n$ bersifat Cauchy, sehingga
terdapat suatu $N$ sedemikian sehingga untuk semua $m, n \geq N$ dengan $m > n$, berlaku
\begin{equation*}
\sum_{k=n+1}^m M_k < \epsilon .
\end{equation*}
Kita menaksir selisih Cauchy antara jumlah-jumlah parsial
fungsi-fungsi tersebut:
\begin{equation*}
\abs{\sum_{k=n+1}^m f_k(x)} \leq
\sum_{k=n+1}^m \babs{f_k(x)} \leq
\sum_{k=n+1}^m M_k < \epsilon .
\end{equation*}
Deret tersebut konvergen berdasarkan \propref{prop:cachysercomplex}.
Konvergensinya seragam karena $N$ tidak bergantung pada $x$.
Sesungguhnya, untuk semua $n \geq N$,
\begin{equation*}
\abs{\sum_{k=1}^\infty f_k(x) - \sum_{k=1}^n f_k(x)} \leq
\abs{\sum_{k=n+1}^\infty f_k(x)} \leq \epsilon .
\qedhere
\end{equation*}
\end{proof}

\begin{example} \label{example:sinnsqfourier}
Deret
\begin{equation*}
\sum_{n=1}^\infty \frac{\sin(nx)}{n^2}
\end{equation*}
konvergen seragam pada $\R$. Lihat \figureref{fig:fouriersern2}.
Deret ini merupakan deret Fourier, dan
kita akan melihat lebih banyak deret semacam ini pada bagian selanjutnya. Pembuktian:
Deret tersebut konvergen seragam karena
$\sum_{n=1}^\infty \frac{1}{n^2}$
konvergen dan
\begin{equation*}
\abs{\frac{\sin(nx)}{n^2}} \leq 
\frac{1}{n^2} .
\end{equation*}
\end{example}

\begin{myfigureht}
\myincludegraphics{fouriersern2}{%
Grafik suatu fungsi yang digambarkan dengan garis tebal, berbentuk
gelombang dengan lereng turun yang landai dan lereng naik yang curam.
Beberapa aproksimasi diperlihatkan dalam berbagai nuansa abu-abu;
aproksimasi yang lebih gelap (lebih awal) lebih menyerupai gelombang sinus biasa,
sedangkan yang lebih terang lebih menyerupai garis tebal gelap.}
\caption{Grafik
$\sum_{n=1}^\infty \frac{\sin(n x)}{n^2}$ yang mencakup
8 jumlah parsial pertamanya dalam berbagai nuansa abu-abu.\label{fig:fouriersern2}}
\end{myfigureht}

\begin{example}
Deret
\begin{equation*}
\sum_{n=0}^\infty \frac{x^n}{n!} 
\end{equation*}
konvergen seragam pada setiap interval terbatas.
Deret ini merupakan deret pangkat yang akan segera kita pelajari.
Pembuktian: Ambil interval $[-r,r] \subset \R$ (setiap interval terbatas
termuat dalam suatu $[-r,r]$).
Deret $\sum_{n=0}^\infty \frac{r^n}{n!}$ konvergen berdasarkan uji rasio,
sehingga $\sum_{n=0}^\infty \frac{x^n}{n!}$ konvergen seragam pada $[-r,r]$ karena
\begin{equation*}
\abs{\frac{x^n}{n!} } \leq 
\frac{r^n}{n!} .
\end{equation*}
\end{example}

Sekarang kita ingin mengetahui sesuatu tentang limit tersebut. Sebagai contoh,
apakah limit itu kontinu?


\begin{prop} \label{prop:uniformswitch}
Misalkan $(X,d_X)$ dan $(Y,d_Y)$ ruang-ruang metrik, dan andaikan $(Y,d_Y)$
lengkap secara Cauchy.
Andaikan $f_n \colon X \to Y$ konvergen seragam ke
 $f \colon X \to Y$.
Misalkan $\{ x_k \}_{k=1}^\infty$ suatu barisan dalam $X$ dan $x \coloneqq \lim_{k \to \infty} x_k$. Andaikan
\begin{equation*}
a_n \coloneqq \lim_{k \to \infty} f_n(x_k)
\end{equation*}
ada untuk semua $n$. Maka
$\{a_n\}_{n=1}^\infty$ konvergen dan
\begin{equation*}
\lim_{k \to \infty} f(x_k) = \lim_{n\to\infty} a_n .
\end{equation*}
\end{prop}

Dengan kata lain,
\begin{equation*}
\lim_{k \to \infty} \lim_{n\to\infty} f_n(x_k) =
\lim_{n \to \infty} \lim_{k\to\infty} f_n(x_k) .
\end{equation*}

\begin{proof}
Pertama, kita tunjukkan bahwa $\{ a_n \}_{n=1}^\infty$ konvergen. Karena
$\{ f_n \}_{n=1}^\infty$ konvergen seragam, barisan ini bersifat Cauchy seragam.
Ambil $\epsilon > 0$. Terdapat
suatu $M$ sedemikian sehingga untuk semua $m,n \geq M$, berlaku
\begin{equation*}
d_Y\bigl(f_n(x_k),f_m(x_k)\bigr) < \epsilon \qquad \text{untuk semua } k .
\end{equation*}
Perhatikan bahwa
$d_Y(a_n,a_m) \leq
d_Y\bigl(a_n,f_n(x_k)\bigr) +
d_Y\bigl(f_n(x_k),f_m(x_k)\bigr) +
d_Y\bigl(f_m(x_k),a_m\bigr)$. Dengan mengambil limit ketika $k \to \infty$, kita memperoleh
\begin{equation*}
d_Y(a_n,a_m) \leq \epsilon .
\end{equation*}
Karena itu, $\{a_n\}_{n=1}^\infty$ bersifat Cauchy dan konvergen karena $Y$ lengkap. Tuliskan
$a \coloneqq \lim_{n \to \infty} a_n$.

Pilih suatu $\ell \in \N$ sedemikian sehingga
\begin{equation*}
d_Y\bigl(f_\ell(p),f(p)\bigr) < \nicefrac{\epsilon}{3}
\end{equation*}
untuk semua $p \in X$. Pastikan pula $\ell$ cukup besar
sehingga
\begin{equation*}
d_Y(a_\ell,a) < \nicefrac{\epsilon}{3}  .
\end{equation*}
Pilih suatu $N \in \N$ sedemikian sehingga untuk $m \geq N$,
\begin{equation*}
d_Y\bigl(f_\ell(x_m),a_\ell\bigr) < \nicefrac{\epsilon}{3}  .
\end{equation*}
Maka untuk
$m \geq N$,
\begin{equation*}
d_Y\bigl(f(x_m),a\bigr)
\leq
d_Y\bigl(f(x_m),f_\ell(x_m)\bigr)
+
d_Y\bigl(f_\ell(x_m),a_\ell\bigr)
+
d_Y\bigl(a_\ell,a\bigr)
<
\nicefrac{\epsilon}{3} +
\nicefrac{\epsilon}{3} +
\nicefrac{\epsilon}{3} = \epsilon . \qedhere
\end{equation*}
\end{proof}

Kita memperoleh korolari langsung mengenai kekontinuan.
Jika semua $f_n$ kontinu, maka $a_n = f_n(x)$.
Karena itu,
$\{ a_n \}_{n=1}^\infty$ secara otomatis konvergen ke
$f(x)$ berdasarkan hipotesis, sehingga kita tidak memerlukan kelengkapan $Y$.

\begin{cor} \label{cor:metricuniformcontinuous}
Misalkan $X$ dan $Y$ ruang-ruang metrik.
Jika $f_n \colon X \to Y$ fungsi-fungsi kontinu
sedemikian sehingga
$\{ f_n \}_{n=1}^\infty$ konvergen seragam ke $f \colon X \to Y$,
maka $f$ kontinu.
\end{cor}

Kebalikannya tidak benar. Fakta bahwa limitnya kontinu tidak berarti
bahwa konvergensinya seragam. Sebagai contoh:
$f_n \colon (0,1) \to \R$ yang didefinisikan oleh $f_n(x) \coloneqq x^n$ konvergen ke
fungsi nol, tetapi tidak secara seragam. Namun, jika kita menambahkan syarat-syarat tambahan
pada barisan tersebut, kita dapat memperoleh kebalikan parsial seperti teorema Dini,
\volIref{lihat \exerciseref*{vI-exercise:dinisthm} dari jilid I}{lihat \exerciseref{exercise:dinisthm}}.

Dalam \exerciseref{exercise:CXCnormedspace}, pembaca diminta untuk membuktikan
bahwa untuk $X$ kompak, $C(X,\C)$ merupakan ruang vektor bernorma
dengan norma seragam dan dengan demikian merupakan ruang metrik. Kita baru saja menunjukkan bahwa
$C(X,\C)$ lengkap secara Cauchy: \propref{prop:unifcauchymetric} mengatakan bahwa suatu barisan Cauchy
dalam $C(X,\C)$ konvergen seragam ke suatu fungsi,
dan \corref{cor:metricuniformcontinuous} menunjukkan bahwa limitnya
kontinu dan dengan demikian merupakan anggota $C(X,\C)$.

\begin{cor}
Misalkan $(X,d)$ suatu ruang metrik kompak.
Maka $C(X,\C)$ merupakan ruang metrik yang lengkap secara Cauchy.
\end{cor}

\begin{example}
Berdasarkan \exampleref{example:sinnsqfourier},
deret Fourier
\begin{equation*}
\sum_{n=1}^\infty \frac{\sin(nx)}{n^2}
\end{equation*}
konvergen seragam dan karenanya kontinu berdasarkan \corref{cor:metricuniformcontinuous} (seperti yang tampak
pada \figureref{fig:fouriersern2}).
\end{example}

\subsection{Integrasi}

\begin{prop} \label{prop:complexlimitswapintegral}
Misalkan $f_n \colon [a,b] \to \C$
terintegralkan secara Riemann dan andaikan $\{ f_n \}_{n=1}^\infty$ konvergen
seragam ke $f \colon [a,b] \to \C$. Maka $f$ terintegralkan secara Riemann
dan
\begin{equation*}
\int_a^b f = \lim_{n\to \infty} \int_a^b f_n .
\end{equation*}
\end{prop}

Karena integral suatu fungsi bernilai kompleks hanyalah integral
bagian real dan imajinernya secara terpisah,
pembuktiannya langsung mengikuti hasil-hasil dalam
\volIref{\chapterref*{vI-fs:chapter} dari jilid~I}{\chapterref{fs:chapter}}.
Rinciannya kita serahkan sebagai latihan.

\begin{cor}
\pagebreak[2]
Misalkan $f_n \colon [a,b] \to \C$
terintegralkan secara Riemann dan andaikan bahwa
\begin{equation*}
\sum_{n=1}^\infty f_n(x)
\end{equation*}
konvergen seragam. Maka deret tersebut terintegralkan secara Riemann pada $[a,b]$
dan
\begin{equation*}
\int_a^b \sum_{n=1}^\infty f_n(x) \,dx
=
\sum_{n=1}^\infty \int_a^b f_n(x) \,dx
\end{equation*}
\end{cor}

\begin{example}
Mari kita tunjukkan cara mengintegralkan suatu deret Fourier.
\begin{equation*}
\int_{0}^x \sum_{n=1}^\infty \frac{\cos(nt)}{n^2} \,dt
=
\sum_{n=1}^\infty \int_{0}^x \frac{\cos(nt)}{n^2}\,dt
=
\sum_{n=1}^\infty \frac{\sin(nx)}{n^3}
\end{equation*}
Pertukaran integral dan jumlah dimungkinkan oleh konvergensi seragam,
yang sebelumnya telah kita buktikan menggunakan uji-$M$ Weierstrass
(\thmref{thm:weiermtest}).
\end{example}

Perhatikan bahwa kita dapat mempertukarkan integral dan limit di bawah hipotesis yang jauh lebih lemah,
tetapi untuk itu kita memerlukan integral yang lebih kuat daripada integral Riemann,
seperti integral Lebesgue.

\subsection{Diferensiasi}

Ingat bahwa suatu fungsi bernilai kompleks
$f \colon [a,b] \to \C$, dengan $f(x) = u(x)+i\,v(x)$,
diferensiabel jika $u$ dan $v$ diferensiabel
dan turunannya adalah
\begin{equation*}
f'(x) = u'(x)+i\,v'(x) .
\end{equation*}

Pembuktian teorema berikut dilakukan dengan menerapkan teorema yang bersesuaian untuk
fungsi real pada $u$ dan $v$, dan diserahkan sebagai latihan.

\begin{thm} \label{thm:dersconvergecomplex}
Misalkan $I \subset \R$ suatu interval terbatas dan misalkan
$f_n \colon I \to \C$ fungsi-fungsi yang diferensiabel secara kontinu.
Andaikan $\{ f_n' \}_{n=1}^\infty$ konvergen seragam ke $g \colon I \to \C$,
dan andaikan $\{ f_n(c) \}_{n=1}^\infty$ merupakan suatu
barisan konvergen untuk suatu $c \in I$.  Maka $\{ f_n \}_{n=1}^\infty$ konvergen seragam ke 
suatu fungsi $f \colon I \to \C$ yang diferensiabel secara kontinu, dan $f' = g$.
\end{thm}

Konvergensi seragam dari fungsi-fungsi itu sendiri tidaklah cukup.
Turunannya mungkin tidak konvergen dan fungsi limitnya bahkan tidak harus
diferensiabel.
Dalam \sectionref{sec:stoneweier}, kita akan membuktikan bahwa
fungsi-fungsi kontinu merupakan limit seragam dari polinom, tetapi seperti yang ditunjukkan oleh
contoh berikut, suatu fungsi kontinu tidak harus diferensiabel
di titik mana pun.

\begin{example}
Terdapat fungsi kontinu yang tidak diferensiabel di mana pun.
Fungsi semacam itu sering disebut
\emph{fungsi Weierstrass}\index{fungsi Weierstrass},
meskipun fungsi
khusus ini, yang pada hakikatnya berasal dari
Takagi\footnote{\href{https://en.wikipedia.org/wiki/Teiji_Takagi}{Teiji
Takagi} (1875--1960) adalah seorang matematikawan Jepang.}, merupakan contoh yang berbeda dari yang diberikan Weierstrass.
Definisikan
\begin{equation*}
\varphi(x) \coloneqq \sabs{x} \qquad \text{untuk } x \in [-1,1] .
\end{equation*}
Perluas $\varphi$ ke seluruh $\R$ dengan menjadikannya
2-periodik:
Tetapkan bahwa
$\varphi(x) = \varphi(x+2)$.  Fungsi $\varphi \colon \R \to \R$
kontinu; bahkan, $\babs{\varphi(x)-\varphi(y)} \leq \sabs{x-y}$ (mengapa?).
Lihat \figureref{fig:triangwave}.
\begin{myfigureht}
\myincludegraphics{triangwave}{%
Suatu fungsi yang grafiknya berupa deretan segitiga; yakni,
fungsi tersebut bernilai 0 pada semua bilangan bulat genap, dan bernilai 1 pada semua bilangan bulat
ganjil.  Di antaranya, grafik tersebut berupa garis lurus.}
\caption{Fungsi 2-periodik $\varphi$.\label{fig:triangwave}}
\end{myfigureht}

Karena $\sum_{n=0}^\infty {\left(\frac{3}{4}\right)}^n$ konvergen dan $\babs{\varphi(x)} \leq
1$ untuk semua $x$, menurut uji-$M$
(\thmref{thm:weiermtest}),
\begin{equation*}
f(x) \coloneqq \sum_{n=0}^\infty 
{\left(\frac{3}{4}\right)}^n \varphi(4^n x)
\end{equation*}
konvergen seragam sehingga kontinu.
Lihat \figureref{fig:nowherediff}.

\begin{myfigureht}
\myincludegraphics{nowherediff}{%
Grafik suatu fungsi yang sangat liar dengan banyak 
puncak tinggi yang runcing dan titik minimum yang rendah, seakan-akan digambar dengan menggerakkan pena secara sangat
cepat naik turun.}
\caption{Grafik fungsi $f$ yang tidak diferensiabel di mana pun.\label{fig:nowherediff}}
\end{myfigureht}

Kita mengklaim bahwa $f \colon
\R \to \R$ tidak diferensiabel di mana pun.
Ambil sembarang $x$, dan akan kita tunjukkan bahwa $f$ tidak diferensiabel di $x$.
Definisikan
\begin{equation*}
\delta_m \coloneqq \pm \frac{1}{2} 4^{-m} ,
\avoidbreak
\end{equation*}
dengan tanda dipilih sedemikian sehingga tidak ada bilangan bulat
di antara $4^m x$ dan $4^m(x+\delta_m) = 4^m x \pm \frac{1}{2}$.

Kita akan meninjau hasil bagi selisih
\begin{equation*}
\frac{f(x+\delta_m)-f(x)}{\delta_m}
=
\sum_{n=0}^\infty 
{\left(\frac{3}{4}\right)}^n
\frac{\varphi\bigl(4^n(x+\delta_m)\bigr)-\varphi(4^nx)}{\delta_m} .
\end{equation*}
Untuk sementara, tetapkan $m$.  Tinjau ekspresi di dalam deret:
\begin{equation*}
\gamma_{n} \coloneqq
\frac{\varphi\bigl(4^n(x+\delta_m)\bigr)-\varphi(4^nx)}{\delta_m} .
\end{equation*}
Jika $n > m$, maka $4^n\delta_m$ merupakan bilangan bulat genap.  Karena $\varphi$
2-periodik, diperoleh $\gamma_n = 0$.

Karena tidak ada bilangan bulat di antara 
$4^m(x+\delta_m) = 4^m x\pm\nicefrac{1}{2}$ dan $4^m x$, maka pada interval ini
$\varphi(t) = \pm t + \ell$ untuk suatu bilangan bulat $\ell$.
Secara khusus,
$\abs{\varphi\bigl(4^m(x+ \delta_m)\bigr)-\varphi(4^mx)} =
\abs{4^mx\pm\nicefrac{1}{2}-4^mx} = \nicefrac{1}{2}$.  Oleh karena itu,
\begin{equation*}
\sabs{\gamma_m} =
\abs{
\frac{\varphi\bigl(4^m(x+\delta_m)\bigr)-\varphi(4^mx)}{\pm (\nicefrac{1}{2}) 4^{-m}}
}
= 4^m .
\end{equation*}
Demikian pula, andaikan $n < m$.  Karena $\babs{\varphi(s) -\varphi(t)} \leq
\sabs{s-t}$,
\begin{equation*}
\sabs{\gamma_n} =
\abs{\frac{\varphi\bigl(4^nx\pm(\nicefrac{1}{2})4^{n-m}\bigr)-\varphi(4^nx)}{\pm
(\nicefrac{1}{2}) 4^{-m}}}
\leq
\abs{\frac{\pm(\nicefrac{1}{2})4^{n-m}}{\pm (\nicefrac{1}{2}) 4^{-m}}} = 4^n
.
\end{equation*}

Dengan demikian,
\begin{equation*}
\begin{split}
\abs{
\frac{f(x+\delta_m)-f(x)}{\delta_m}
}
%& =
%\abs{
%\sum_{n=0}^\infty 
%{\left(\frac{3}{4}\right)}^n
%\frac{\varphi\bigl(4^n(x+\delta_m)\bigr)-\varphi(4^nx)}{\delta_m}
%}
=
\abs{
\sum_{n=0}^\infty 
{\left(\frac{3}{4}\right)}^n
\gamma_n
}
& =
\abs{
\sum_{n=0}^m 
{\left(\frac{3}{4}\right)}^n
\gamma_n
}
\\
& \geq
\abs{
{\left(\frac{3}{4}\right)}^m
\gamma_m}
-
\abs{
\sum_{n=0}^{m-1} 
{\left(\frac{3}{4}\right)}^n
\gamma_n
}
\\
& \geq
3^m
-
\sum_{n=0}^{m-1} 
3^n
=
3^m
-
\frac{3^{m}-1}{3-1}
=
\frac{3^m +1}{2} .
\end{split}
\end{equation*}
Ketika $m \to \infty$, kita memperoleh
$\delta_m \to 0$, tetapi $\frac{3^m+1}{2}$
menuju tak hingga.  Jadi $f$ tidak mungkin diferensiabel di~$x$.
\end{example}

\subsection{Latihan}

\begin{exercise}
Buktikan \propref{prop:uniformconvbounded}.
\end{exercise}

\begin{exercise}
Buktikan \propref{prop:unifcauchymetric}.
\end{exercise}

\begin{exercise} \label{exercise:CXCnormedspace}
Misalkan $(X,d)$ adalah ruang metrik kompak.
Buktikan bahwa norma seragam $\snorm{\cdot}_X$ merupakan norma pada ruang vektor
fungsi-fungsi kontinu bernilai kompleks $C(X,\C)$.
\end{exercise}

\begin{exercise}
\pagebreak[2]
\leavevmode
\begin{enumerate}[a)]
\item
Buktikan bahwa
$f_n(x) \coloneqq 2^{-n} \sin(2^n x)$
konvergen seragam ke nol, tetapi terdapat himpunan padat $D \subset \R$
sedemikian sehingga $\lim_{n\to\infty} f_n'(x) = 1$ untuk setiap $x \in D$.
\item
Buktikan bahwa
$\sum_{n=1}^\infty 2^{-n} \sin(2^n x)$
konvergen seragam ke suatu fungsi kontinu,
dan terdapat himpunan padat $D \subset \R$
sehingga turunan jumlah-jumlah parsial tersebut tidak konvergen di setiap titik himpunan itu.
\end{enumerate}
\end{exercise}

\begin{exercise}
Buktikan bahwa $\snorm{f}_{C^1} \coloneqq \snorm{f}_{[a,b]}+\snorm{f'}_{[a,b]}$
merupakan norma pada ruang vektor fungsi-fungsi bernilai kompleks yang
diferensiabel secara kontinu $C^1\bigl([a,b],\C\bigr)$.
\end{exercise}

\begin{exercise}
Buktikan \thmref{thm:dersconvergecomplex}.
\end{exercise}

\begin{exercise}
Buktikan \propref{prop:complexlimitswapintegral} dengan mereduksinya ke
hasil untuk kasus real.
\end{exercise}

\begin{exercise}
Kerjakan rincian contoh penyangkal berikut untuk kebalikan
uji-$M$ Weierstrass (\thmref{thm:weiermtest}).  Definisikan
$f_n \colon [0,1] \to \R$ melalui
\begin{equation*}
f_n(x) \coloneqq
\begin{cases}
\frac{1}{n} & \text{jika } \frac{1}{n+1} < x < \frac{1}{n},\\
0	    & \text{selainnya.}
\end{cases}
\end{equation*}
Buktikan bahwa $\sum_{n=1}^\infty f_n$ konvergen seragam, tetapi
$\sum_{n=1}^\infty \snorm{f_n}_{[0,1]}$
tidak konvergen.
\end{exercise}

\begin{exercise}
Misalkan $f_n \colon [0,1] \to \R$ adalah fungsi-fungsi yang monoton naik dan
misalkan $\sum_{n=1}^\infty f_n$ konvergen titik demi titik.  Buktikan bahwa
$\sum_{n=1}^\infty f_n$
konvergen seragam.
\end{exercise}

\begin{exercise}
Buktikan bahwa
\begin{equation*}
\sum_{n=1}^\infty e^{-nx}
\end{equation*}
konvergen untuk setiap $x > 0$ ke suatu fungsi yang diferensiabel.
\end{exercise}

%%%%%%%%%%%%%%%%%%%%%%%%%%%%%%%%%%%%%%%%%%%%%%%%%%%%%%%%%%%%%%%%%%%%%%%%%%%%%%

\sectionnewpage
\section{Deret pangkat dan fungsi analitik}
\label{sec:analfuncs}

%mbxINTROSUBSECTION

\sectionnotes{2--3 kuliah}

\subsection{Fungsi analitik}

Suatu deret pangkat (kompleks) adalah deret berbentuk
\begin{equation*}
\sum_{n=0}^\infty c_n {(z-a)}^n
\end{equation*}
untuk $c_n, z, a \in \C$.  Kita mengatakan deret tersebut
\emph{konvergen}\index{konvergen!deret pangkat} jika deret itu konvergen untuk
suatu $z \neq a$.

Misalkan $U \subset \C$ merupakan himpunan terbuka dan
$f \colon U \to \C$ suatu fungsi.
Andaikan untuk setiap $a \in U$ terdapat $\rho > 0$ dan suatu deret
pangkat yang konvergen ke fungsi tersebut,
\begin{equation*}
f(z) = \sum_{n=0}^\infty c_n {(z-a)}^n
\end{equation*}
untuk semua $z \in B(a,\rho)$.
Maka kita mengatakan bahwa $f$ adalah fungsi \emph{\myindex{analitik}}.
Serupa dengan itu, untuk interval $(a,b) \subset \R$, kita mengatakan bahwa $f \colon (a,b) \to \C$
bersifat analitik atau dapat pula disebut \emph{\myindex{analitik real}}
jika untuk setiap titik $c \in
(a,b)$ terdapat deret pangkat di sekitar $c$ yang konvergen pada suatu
$(c-\rho,c+\rho)$
untuk suatu $\rho > 0$.
Karena kita kadang-kadang akan membahas deret pangkat real dan kadang-kadang
deret pangkat kompleks, kita akan memakai $z$ untuk menyatakan bilangan kompleks dan $x$ bilangan real.
Kita akan selalu menyebutkan kasus mana yang sedang dibahas.

Suatu fungsi analitik memiliki pengembangan yang berbeda di sekitar titik-titik yang berbeda.
Selain itu, konvergensi tidak otomatis terjadi pada seluruh
domain fungsi.  Sebagai contoh, jika $\sabs{z} < 1$, maka
\begin{equation*}
\frac{1}{1-z} = \sum_{n=0}^\infty z^n .
\end{equation*}
Walaupun ruas kiri terdefinisi untuk setiap $z \neq 1$, ruas kanan
ternyata hanya konvergen jika $\sabs{z} < 1$.  Lihat grafik
sebagian kecil dari $\frac{1}{1-z}$ pada \figureref{fig:1over1mz}.
Kita tidak dapat menggambar grafik
fungsi itu sendiri; kita hanya dapat menggambar bagian real atau bagian imajinernya karena
dimensi alam semesta kita tidak mencukupi.

\begin{myfigureht}
\myincludepdft{real_imag_1over1mz}{%
Dua grafik permukaan 3D di atas bidang xy.  Pada kedua grafik, suatu garis vertikal
ditandai di atas titik (1,0) pada bidang xy.  Grafik-grafik tersebut ditampilkan pada
petak kecil di sekitar titik asal dan berakhir sebelum mencapai titik (1,0).
Grafik kiri lebih datar jauh dari garis vertikal dan menanjak
curam ketika mendekati garis vertikal.  Grafik kanan serupa,
tetapi menanjak curam ketika mendekati garis vertikal dari
daerah tempat y bernilai negatif dan menurun curam ketika mendekati
garis vertikal dari daerah tempat y bernilai positif.}
\caption{Grafik bagian real dan bagian imajiner dari $z=x+iy \mapsto \frac{1}{1-z}$
pada persegi $[-0.8,0.8]^2$.  Singularitas di $z=1$ ditandai dengan
garis vertikal putus-putus.\label{fig:1over1mz}}
\end{myfigureht}

\subsection{Konvergensi deret pangkat}

Kita telah membuktikan beberapa hasil untuk deret pangkat dengan peubah real dalam
\volIref{\sectionref*{vI-sec:moreonseries} pada Jilid I}{\sectionref{sec:moreonseries}}.   Sebagian besar sifat
konvergensi deret pangkat berkaitan dengan deret
$\sum_{k=0}^\infty \sabs{c_k} \, \sabs{z-a}^k$,
sehingga kita sebenarnya telah membuktikan banyak hasil mengenai deret pangkat
kompleks.  Khususnya, kita telah menghitung apa yang disebut
jari-jari konvergensi suatu deret pangkat.

\begin{prop}
\pagebreak[2]
Misalkan $\sum_{n=0}^\infty c_n {(z-a)}^n$ suatu deret pangkat.
Terdapat $\rho \in [0,\infty]$ sedemikian sehingga
\begin{enumerate}[(i)]
\item Jika $\rho = 0$, maka deret tersebut divergen untuk setiap $z \neq a$.
\item Jika $\rho = \infty$, maka deret tersebut konvergen mutlak untuk semua $z \in \C$.
\item Jika $0 < \rho < \infty$, maka deret tersebut
konvergen mutlak pada $B(a,\rho)$,
dan divergen ketika $\sabs{z-a} > \rho$.
\end{enumerate}
Lebih lanjut, jika $0 < r < \rho$, maka
deret tersebut konvergen secara seragam pada bola tertutup $C(a,r)$.
\end{prop}

Bilangan $\rho$ disebut \emph{\myindex{jari-jari konvergensi}}.
Lihat \figureref{fig:radiusconvcomplex}.
Jari-jari konvergensi menentukan sebuah cakram berpusat di $a$ tempat deret tersebut konvergen.  Suatu deret pangkat
dikatakan konvergen jika $\rho > 0$.
\begin{myfigureht}
\myincludepdft{radiusconvcomplex}{%
Sebuah lingkaran bertitik-titik berpusat di a dengan jari-jari rho.
Bagian dalam lingkaran diarsir dan diberi label "deret konvergen".
Bagian luar lingkaran diberi label "deret tidak konvergen".}
\caption{Jari-jari konvergensi.\label{fig:radiusconvcomplex}}
\end{myfigureht}

\begin{proof}
Kita menggunakan versi real dari proposisi ini,
\volIref{\propref*{vI-prop:powerserrealradius} pada Jilid I}{\propref{prop:powerserrealradius}}.
Misalkan
\begin{equation*}
R \coloneqq \limsup_{n\to\infty} \sqrt[n]{\sabs{c_n}} .
\end{equation*}
Jika $R = 0$, maka
$\sum_{n=0}^\infty \sabs{c_n} \, \sabs{z-a}^n$ konvergen untuk semua $z$.
Jika $R = \infty$, maka
$\sum_{n=0}^\infty \sabs{c_n} \, \sabs{z-a}^n$ hanya konvergen di $z=a$.
Jika tidak, misalkan $\rho \coloneqq \nicefrac{1}{R}$.
Maka
$\sum_{n=0}^\infty \sabs{c_n} \, \sabs{z-a}^n$ konvergen ketika
$\sabs{z-a} < \rho$, dan divergen (bahkan suku-suku deretnya
tidak menuju nol) ketika $\sabs{z-a} > \rho$.

Untuk membuktikan pernyataan \myquote{Lebih lanjut,} andaikan
$0 < r < \rho$ dan $z \in C(a,r)$.  Perhatikan
jumlah parsial
\begin{equation*}
\abs{\sum_{n=0}^k c_n {(z-a)}^n}
\leq
\sum_{n=0}^k \sabs{c_n} \sabs{z-a}^n
\leq
\sum_{n=0}^k \sabs{c_n} r^n . \qedhere
\end{equation*}
\end{proof}

Jika $\sum_{n=0}^\infty c_n {(z-a)}^n$ konvergen untuk suatu $z$, maka
\begin{equation*}
\sum_{n=0}^\infty c_n {(w-a)}^n
\end{equation*}
konvergen mutlak setiap kali $\sabs{w-a} < \sabs{z-a}$.
Sebaliknya, jika deret tersebut divergen di $z$, maka deret itu juga harus divergen di $w$
setiap kali $\sabs{w-a} > \sabs{z-a}$.
Jadi, untuk menunjukkan
bahwa jari-jari konvergensinya sekurang-kurangnya sebesar suatu bilangan tertentu, kita hanya perlu
menunjukkan konvergensi di suatu titik dengan metode apa pun yang kita ketahui.

\begin{example}
Kita daftarkan beberapa deret yang telah kita kenal:
\begin{align*}
& &
& \sum_{n=0}^\infty z^n
& & \text{memiliki jari-jari konvergensi } 1.
& &
\\
& &
& \sum_{n=0}^\infty \frac{1}{n!} z^n
& & \text{memiliki jari-jari konvergensi } \infty.
& &
\\
& &
& \sum_{n=0}^\infty n^n z^n
& & \text{memiliki jari-jari konvergensi } 0.
& &
\end{align*}
\end{example}

\begin{example}
Perhatikan perbedaan antara $\frac{1}{1-z}$ dan deret pangkatnya.  Mari kita
kembangkan $\frac{1}{1-z}$ sebagai deret pangkat di sekitar titik $a \neq 1$.
Misalkan $c \coloneqq \frac{1}{1-a}$, maka
\begin{equation*}
\frac{1}{1-z} =
\frac{c}{1-c(z-a)}
=
c
\sum_{n=0}^\infty c^{n} {(z-a)}^n
=
\sum_{n=0}^\infty \left( \frac{1}{{(1-a)}^{n+1}} \right) {(z-a)}^n .
\end{equation*}
Deret $\sum_{n=0}^\infty c^n {(z-a)}^n$ konvergen jika dan hanya jika
deret di ruas kanan konvergen.  Kita hitung,
\begin{equation*}
\limsup_{n\to\infty}
\sqrt[n]{\sabs{c^n}} = \sabs{c}
= \frac{1}{\sabs{1-a}} .
\end{equation*}
Jari-jari konvergensi deret pangkat tersebut adalah $\sabs{1-a}$, yaitu
jarak dari $1$ ke $a$.  Fungsi $\frac{1}{1-z}$
memiliki representasi deret pangkat
di sekitar setiap $a \neq 1$ sehingga bersifat analitik pada $\C \setminus
\{ 1 \}$.
Domain fungsi tersebut lebih luas daripada daerah
konvergensi deret pangkat yang merepresentasikan fungsi itu di sekitar titik mana pun.
\end{example}

Ternyata,
jika suatu fungsi memiliki representasi deret pangkat yang konvergen ke fungsi tersebut
pada suatu bola,
maka fungsi itu memiliki representasi deret pangkat di sekitar setiap titik dalam bola tersebut.  Kita akan membuktikan
hasil ini nanti.

\subsection{Sifat-sifat fungsi analitik}

\begin{prop}
Jika
\begin{equation*}
f(z) \coloneqq \sum_{n=0}^\infty c_n {(z-a)}^n
\end{equation*}
konvergen pada $B(a,\rho)$ untuk suatu $\rho > 0$, maka
$f \colon B(a,\rho) \to \C$ kontinu.
Khususnya, fungsi analitik bersifat kontinu.
\end{prop}

\begin{proof}
Untuk $z_0 \in B(a,\rho)$, pilih $r < \rho$ sedemikian sehingga $z_0 \in B(a,r)$.
Pada $B(a,r)$,
jumlah parsial (yang kontinu) konvergen secara seragam,
sehingga limitnya, $f|_{B(a,r)}$, kontinu.
Setiap barisan yang konvergen ke
$z_0$ memiliki suatu ekor yang seluruhnya berada dalam bola terbuka $B(a,r)$.
Jadi, $f$ kontinu di $z_0$.
\end{proof}

Dalam \volIref{\corref*{vI-cor:differentiatepowerser} pada Jilid
I}{\corref{cor:differentiatepowerser}}, kita telah membuktikan bahwa kita dapat
mendiferensialkan deret pangkat real suku demi suku.  Dengan kata lain,
kita membuktikan bahwa jika
\begin{equation*}
f(x) \coloneqq \sum_{n=0}^\infty c_n {(x-a)}^n
\end{equation*}
konvergen untuk $x$ real pada suatu interval di sekitar $a \in \R$, maka kita dapat mendiferensialkannya
suku demi suku dan memperoleh deret
\begin{equation*}
f'(x) =
\sum_{n=1}^\infty n c_n {(x-a)}^{n-1}
=
\sum_{n=0}^\infty (n+1)c_{n+1} {(x-a)}^{n} 
\end{equation*}
dengan jari-jari konvergensi yang sama.
Kita hanya membuktikan teorema ini ketika $c_n$ bernilai real.
Namun, untuk $c_n$ kompleks,
kita tulis
$c_n = s_n + i t_n$, dan karena $x$ serta $a$ real,
\begin{equation*}
\sum_{n=0}^\infty c_n {(x-a)}^n
=
\sum_{n=0}^\infty s_n {(x-a)}^n
+
i
\sum_{n=0}^\infty t_n {(x-a)}^n .
\end{equation*}
Kita terapkan teorema tersebut pada bagian real dan
bagian imajiner.

Dengan menerapkan teorema ini berulang kali, kita peroleh bahwa suatu
fungsi analitik diferensiabel tak berhingga kali:
\begin{equation*}
%mbxlatex \begin{aligned}
f^{(\ell)}(x)
%mbxlatex &
=
\sum_{n=\ell}^\infty n(n-1)\cdots(n-\ell+1)c_n {(x-a)}^{n-\ell}
%mbxlatex \\
%mbxlatex &
=
\sum_{n=0}^\infty (n+\ell)(n+\ell-1)\cdots (n+1) c_{n+\ell} {(x-a)}^{n} .
%mbxlatex \end{aligned}
\end{equation*}
Khususnya,
\begin{equation} \label{eq:formulaforpscoeffs}
f^{(\ell)}(a) = \ell! \, c_\ell .
\avoidbreak
\end{equation}
Koefisien-koefisien tersebut ditentukan secara unik oleh turunan
fungsi tersebut, dan sebaliknya.

Di sisi lain, adanya fungsi yang diferensiabel tak berhingga kali
tidak berarti bahwa bilangan $c_n$ yang diperoleh dari
$c_n = \frac{f^{(n)}(0)}{n!}$ menghasilkan deret pangkat yang konvergen.
Terdapat sebuah teorema, yang tidak akan kita buktikan,
yang menyatakan bahwa untuk sebarang barisan $\{ c_n \}_{n=1}^\infty$, terdapat
fungsi $f$ yang diferensiabel tak berhingga kali sedemikian sehingga
$c_n = \frac{f^{(n)}(0)}{n!}$.  Selain itu, sekalipun deret yang diperoleh
konvergen, deret itu belum tentu konvergen ke fungsi semula.
Sebagai contoh,
lihat \volIref{\exerciseref*{vI-exercise:nonanalytic} pada Jilid I}{\exerciseref{exercise:nonanalytic}}:  Fungsi
berikut
\begin{equation*}
f(x) \coloneqq
\begin{cases}
e^{-1/x} & \text{jika } x > 0,\\
0        & \text{jika } x \leq 0,
\end{cases}
\end{equation*}
diferensiabel tak berhingga kali, dan semua turunannya di titik asal bernilai nol.
Jadi,
deretnya di titik asal hanyalah deret nol; walaupun
deret tersebut konvergen, deret itu tidak konvergen ke~$f$ untuk $x > 0$.

\medskip

Kita dapat menerapkan transformasi afin $z \mapsto z+a$ yang
mengubah deret pangkat di sekitar $a$ menjadi deret di sekitar titik asal.
Artinya, jika
\begin{equation*}
f(z) = \sum_{n=0}^\infty c_n {(z-a)}^n,
\qquad \text{kita tinjau} \qquad
f(z+a) = \sum_{n=0}^\infty c_n {z}^n.
\end{equation*}
Karena itu, biasanya
cukup untuk membuktikan hasil mengenai deret pangkat di sekitar titik asal.
Mulai sekarang, demi kesederhanaan, kita sering mengandaikan $a=0$.

\subsection{Deret pangkat sebagai fungsi analitik}

Kita memerlukan teorema tentang pertukaran limit deret, yakni
teorema Fubini untuk penjumlahan.  Untuk deret real, hasil ini merupakan
\volIref{\exerciseref*{vI-exercise:tonellifubiniforsums} pada Jilid
I}{\exerciseref{exercise:tonellifubiniforsums}}, tetapi sekarang kita memiliki
argumen yang lebih elegan.

\begin{thm}[\myindex{Fubini untuk penjumlahan}] \label{thm:fubiniforsums}
Misalkan $\{ a_{k,m} \}_{k=1,m=1}^\infty$ suatu barisan ganda
bilangan kompleks dan andaikan bahwa untuk setiap $k$, deret
\begin{equation*}
\sum_{m=1}^\infty \sabs{a_{k,m}} \qquad \text{konvergen}
\end{equation*}
dan, lebih lanjut, bahwa
\begin{equation*}
\sum_{k=1}^\infty \left( \sum_{m=1}^\infty \sabs{a_{k,m}} \right)
\qquad \text{konvergen}.
\end{equation*}
Maka
\begin{equation*}
\sum_{k=1}^\infty \left( \sum_{m=1}^\infty a_{k,m} \right)
=
\sum_{m=1}^\infty \left( \sum_{k=1}^\infty a_{k,m} \right) ,
\end{equation*}
dan semua deret yang terlibat konvergen.
\end{thm}

\begin{proof}
Misalkan $E$ himpunan $\{ \nicefrac{1}{n} : n \in \N \} \cup \{ 0 \}$,
dan pandang himpunan tersebut sebagai ruang metrik dengan metrik yang diwarisi dari $\R$.
Definisikan barisan fungsi $f_k \colon E \to \C$
dengan
\begin{equation*}
f_k(\nicefrac{1}{n}) \coloneqq \sum_{m=1}^n a_{k,m}
\qquad
\text{dan}
\qquad
f_k(0) \coloneqq \sum_{m=1}^\infty a_{k,m} .
\end{equation*}
Karena deret tersebut konvergen, setiap $f_k$ kontinu di $0$
(karena 0 adalah satu-satunya titik akumulasi, fungsi-fungsi tersebut kontinu di setiap titik
dalam $E$, tetapi kita tidak memerlukan fakta itu).
Untuk setiap $x \in E$, kita memiliki
\begin{equation*}
\babs{f_k(x)} \leq \sum_{m=1}^\infty \sabs{a_{k,m}} .
\end{equation*}
Karena $\sum_k \sum_m \sabs{a_{k,m}}$ konvergen (dan tidak bergantung pada
$x$), kita tahu bahwa
\begin{equation*}
\sum_{k=1}^n f_k(x)
\end{equation*}
konvergen secara seragam pada $E$.  Definisikan
\begin{equation*}
g(x) \coloneqq \sum_{k=1}^\infty f_k(x) ,
\end{equation*}
yang karena itu merupakan fungsi kontinu di $0$.
Dengan demikian,
\begin{equation*}
\begin{split}
\sum_{k=1}^\infty \left( \sum_{m=1}^\infty a_{k,m} \right)
& =
\sum_{k=1}^\infty f_k(0)
= g(0)
= \lim_{n\to\infty} g(\nicefrac{1}{n}) \\
&= 
\lim_{n\to\infty}\sum_{k=1}^\infty f_k(\nicefrac{1}{n})
= 
\lim_{n\to\infty}\sum_{k=1}^\infty \sum_{m=1}^n a_{k,m} \\
&= 
\lim_{n\to\infty}\sum_{m=1}^n \sum_{k=1}^\infty a_{k,m}
= 
\sum_{m=1}^\infty \left( \sum_{k=1}^\infty a_{k,m} \right) . \qedhere
\end{split}
\end{equation*}
\end{proof}

Sekarang kita buktikan bahwa jika suatu deret konvergen ke sebuah fungsi
pada suatu interval, kita dapat mengembangkan fungsi tersebut di sekitar setiap titik.

\begin{thm}[Teorema Taylor untuk fungsi analitik real]
\index{Teorema Taylor!analitik real}%
\label{thm:tayloranal}
Misalkan
\begin{equation*}
f(x) \coloneqq \sum_{k=0}^\infty a_k x^k
\end{equation*}
merupakan suatu deret pangkat yang konvergen pada $(-\rho,\rho)$ untuk suatu $\rho > 0$.  Untuk setiap $a \in
(-\rho,\rho)$
dan setiap $x$ yang memenuhi $\sabs{x-a} < \rho-\sabs{a}$, berlaku
\begin{equation*}
f(x) =
\sum_{k=0}^\infty \frac{f^{(k)}(a)}{k!} {(x-a)}^{k} .
\end{equation*}
\end{thm}

Deret pangkat yang berpusat di $a$ tentu saja dapat konvergen pada interval yang lebih besar, tetapi
konvergensi pada interval di atas dijamin.  Interval tersebut adalah interval simetris terbesar yang berpusat di
$a$ dan termuat dalam $(-\rho,\rho)$.

\begin{proof}
Untuk $a$ dan $x$ sebagaimana dalam teorema,
tuliskan
\begin{equation*}
\begin{split}
f(x) &= \sum_{k=0}^\infty a_k {\bigl((x-a)+a\bigr)}^k \\
&= \sum_{k=0}^\infty a_k \sum_{m=0}^k \binom{k}{m} a^{k-m} {(x-a)}^m .
\end{split}
\end{equation*}
Definisikan $c_{k,m} \coloneqq a_k \binom{k}{m} a^{k-m}$ jika $m \leq k$ dan $0$ jika $m >
k$.  Maka
\begin{equation} \label{eq:tsproof}
f(x) = \sum_{k=0}^\infty \, \sum_{m=0}^\infty c_{k,m} {(x-a)}^m .
\end{equation}
Mari kita tunjukkan bahwa jumlah ganda tersebut konvergen mutlak.
\begin{equation*}
\begin{split}
\sum_{k=0}^\infty \, \sum_{m=0}^\infty \babs{ c_{k,m} {(x-a)}^m}
& = \sum_{k=0}^\infty \, \sum_{m=0}^k \abs{ a_k \binom{k}{m} a^{k-m} {(x-a)}^m }
\\
& = \sum_{k=0}^\infty \sabs{a_k} \sum_{m=0}^k \binom{k}{m} \sabs{a}^{k-m}
{\sabs{x-a}}^m  \\
& = \sum_{k=0}^\infty \sabs{a_k} {\bigl(\sabs{x-a}+\sabs{a}\bigr)}^k ,
\end{split}
\end{equation*}
dan deret ini konvergen selama
$(\sabs{x-a}+\sabs{a}) < \rho$, atau dengan kata lain jika
$\sabs{x-a} < \rho-\sabs{a}$.

Dengan menggunakan \thmref{thm:fubiniforsums},
tukar urutan penjumlahan dalam \eqref{eq:tsproof}; maka
deret berikut konvergen ketika $\sabs{x-a} < \rho-\sabs{a}$:
\begin{equation*}
f(x) =
\sum_{k=0}^\infty \, \sum_{m=0}^\infty c_{k,m} {(x-a)}^m
=
\sum_{m=0}^\infty
\left( \sum_{k=0}^\infty
c_{k,m} \right) {(x-a)}^m .
\end{equation*}
Rumus yang dinyatakan melalui turunan di $a$ diperoleh dengan
mendiferensialkan deret tersebut untuk memperoleh \eqref{eq:formulaforpscoeffs}.
\end{proof}

Perhatikan bahwa jika suatu deret konvergen untuk peubah real $x \in (a-\rho,a+\rho)$, deret itu juga konvergen
untuk semua bilangan kompleks dalam $B(a,\rho)$.
Kita memperoleh korolari berikut, yang menyatakan bahwa fungsi yang didefinisikan
oleh deret pangkat bersifat analitik.

\begin{cor} \label{cor:powerseranalytic}
Untuk setiap $a \in \C$,
jika $\sum_{k=0}^\infty c_k {(z-a)}^k$ konvergen ke $f(z)$ dalam $B(a,\rho)$ dan $b \in
B(a,\rho)$,
maka terdapat deret pangkat
$\sum_{k=0}^\infty d_k {(z-b)}^k$ yang konvergen ke $f(z)$ dalam $B(b,\rho-\sabs{b-a})$.
\end{cor}

\begin{proof}
Tanpa mengurangi keumuman, andaikan $a=0$ dan $b \neq 0$.  Kita dapat melakukan rotasi sehingga $b$ bernilai real, tetapi
karena hal itu lebih sulit dibayangkan, mari kita lakukan secara eksplisit.
Misalkan $\alpha \coloneqq \frac{\bar{b}}{\sabs{b}}$.
Perhatikan bahwa
\begin{equation*}
\abs{\nicefrac{1}{\alpha}} = \sabs{\alpha}
= 1 .
\end{equation*}
Karena itu, deret
$\sum_{k=0}^\infty c_k {(\nicefrac{z}{\alpha})}^k = 
\sum_{k=0}^\infty c_k \alpha^{-k} {z}^k$
konvergen ke $f(\nicefrac{z}{\alpha})$ dalam $B(0,\rho)$.
Ketika $z=x$ bernilai real,
kita menerapkan \thmref{thm:tayloranal} pada titik $\sabs{b}$ dan memperoleh
sebuah deret yang konvergen
ke $f(\nicefrac{z}{\alpha})$ pada $B(\sabs{b},\rho-\sabs{b})$.
Artinya, terdapat deret konvergen
\begin{equation*}
f(\nicefrac{z}{\alpha}) =
\sum_{k=0}^\infty a_k {\bigl(z - \sabs{b}\bigr)}^k .
\end{equation*}
Dengan menggunakan $\alpha b = \sabs{b}$, kita peroleh
\begin{equation*}
f(z) = f(\nicefrac{\alpha z}{\alpha}) =
\sum_{k=0}^\infty a_k {(\alpha z - \sabs{b})}^k 
=
\sum_{k=0}^\infty a_k\alpha^k {\bigl(z - \nicefrac{\sabs{b}}{\alpha}\bigr)}^k
=
\sum_{k=0}^\infty a_k\alpha^k {(z - b)}^k ,
\end{equation*}
dan deret ini konvergen untuk semua $z$ sedemikian sehingga
$\bigl\lvert \alpha z-\sabs{b}\bigr\rvert < \rho-\sabs{b}$
atau $\sabs{z - b} < \rho-\sabs{b}$.
\end{proof}

Di atas telah kita buktikan bahwa deret pangkat yang konvergen mendefinisikan
fungsi analitik pada daerah konvergensinya.  Sebelumnya juga telah kita tunjukkan bahwa
$\frac{1}{1-z}$ bersifat analitik di setiap titik selain $z=1$.

Perhatikan bahwa meskipun suatu fungsi analitik real bersifat analitik pada
seluruh garis real, hal itu belum tentu berarti fungsi tersebut memiliki representasi deret pangkat
yang konvergen di semua titik.  Sebagai contoh, fungsi
\begin{equation*}
f(x) = \frac{1}{1+x^2}
\end{equation*}
merupakan fungsi analitik real pada $\R$ (latihan).  Deret
pangkat di sekitar titik asal yang konvergen ke $f$
memiliki jari-jari konvergensi tepat $1$.
Dapatkah Anda melihat alasannya? (latihan)

\subsection{Teorema identitas untuk fungsi analitik}

\begin{lemma}
Misalkan $f(z) = \sum_{k=0}^\infty a_k z^k$ adalah deret pangkat konvergen dan
$\{ z_n \}_{n=1}^\infty$ adalah barisan bilangan kompleks tak nol yang konvergen ke 0,
sedemikian sehingga $f(z_n) = 0$ untuk semua $n$.  Maka $a_k = 0$ untuk setiap~$k$.
\end{lemma}

\begin{proof}
Berdasarkan kekontinuan, kita tahu bahwa $f(0) = 0$, sehingga $a_0 = 0$.
Misalkan terdapat suatu $a_k$ yang tak nol.
Ambil $m$ sebagai indeks terkecil sedemikian sehingga $a_m \neq 0$.  Maka
\begin{equation*}
f(z) = \sum_{k=m}^\infty a_k z^k = 
z^m \sum_{k=m}^\infty a_k z^{k-m} =
z^m \sum_{k=0}^\infty a_{k+m} z^{k} .
\end{equation*}
Tuliskan $g(z) = \sum_{k=0}^\infty a_{k+m} z^{k}$ (deret ini konvergen
pada himpunan yang sama seperti deret untuk $f$).  Fungsi $g$ kontinu dan
$g(0) = a_m \neq 0$.  Jadi, terdapat suatu $\delta > 0$ sedemikian sehingga
$g(z) \neq 0$ untuk semua $z \in B(0,\delta)$.  Karena $f(z) = z^m g(z)$,
satu-satunya titik di $B(0,\delta)$ tempat $f(z) = 0$ adalah $z=0$, tetapi hal
ini bertentangan dengan asumsi bahwa $f(z_n) = 0$ untuk semua $n$.
\end{proof}

Ingatlah bahwa dalam ruang metrik $X$, sebuah \emph{titik akumulasi}
(atau terkadang \emph{titik limit}) dari suatu himpunan
$E$ adalah titik $p \in X$ sedemikian sehingga
$B(p,\epsilon) \setminus \{ p \}$ memuat titik-titik dari $E$
untuk setiap $\epsilon > 0$.

\begin{thm}[Teorema identitas]
\label{thm:identityanalytic}%
\index{teorema identitas}%
Misalkan $U \subset \C$ terbuka dan terhubung.  Jika $f \colon U \to \C$
dan $g \colon U \to \C$ merupakan fungsi-fungsi analitik yang bernilai
sama pada suatu himpunan $E \subset U$, dan $E$ memiliki titik akumulasi
di $U$, maka $f(z) = g(z)$ untuk semua $z \in U$.
\end{thm}

Dalam penerapan yang paling umum dari teorema ini, $E$ adalah himpunan terbuka
atau mungkin sebuah kurva.

\begin{proof}
Dengan mengganti $E$, jika perlu, dengan himpunan semua titik $z \in U$ yang
memenuhi $g(z)=f(z)$, kita boleh menganggap $E$ adalah himpunan tersebut.
Perhatikan bahwa $E$ harus tertutup relatif terhadap $U$ karena $f$ dan $g$
kontinu.

Misalkan $p \in U$ adalah suatu titik akumulasi dari $E$.  Dengan mengganti
peubah $w=z-p$ serta mentranslasikan domain dan kedua fungsi, kita boleh
mengambil $p=0$.  Di sekitar $0$, kita memiliki ekspansi
\begin{equation*}
f(z) = \sum_{k=0}^\infty a_k {z}^k 
\qquad
\text{dan}
\qquad
g(z) = \sum_{k=0}^\infty b_k {z}^k ,
\end{equation*}
yang konvergen dalam suatu bola $B(0,\rho)$.  Oleh karena itu, deret
\begin{equation*}
0 = f(z)-g(z) = 
\sum_{k=0}^\infty (a_k-b_k) z^k
\end{equation*}
konvergen dalam $B(0,\rho)$.  Karena $0$ adalah titik akumulasi dari $E$,
terdapat barisan titik tak nol $\{ z_n \}_{n=1}^\infty$ yang konvergen ke $0$
sedemikian sehingga $f(z_n) -g(z_n) = 0$ untuk semua $n$.  Jadi, berdasarkan lema di atas,
$a_k = b_k$ untuk semua $k$.  Oleh karena itu, $B(0,\rho) \subset E$.

Dengan demikian, himpunan semua titik akumulasi dari $E$ terbuka relatif
terhadap $U$.  Himpunan itu juga tertutup relatif terhadap $U$: limit di $U$
dari titik-titik akumulasi $E$ berada di $E$ karena $E$ tertutup relatif
terhadap $U$, dan jelas merupakan titik akumulasi dari $E$.
Karena $U$ terhubung, himpunan semua titik akumulasi dari $E$
sama dengan $U$, atau dengan kata lain $E = U$.
\end{proof}

Dengan membatasi perhatian kita pada $x$ real, kita memperoleh teorema yang sama
untuk himpunan bagian terbuka dan terhubung dari $\R$, yaitu interval terbuka.

\subsection{Latihan}

\begin{exercise}
Misalkan
\begin{equation*}
a_{k,m} \coloneqq
\begin{cases}
1        & \text{jika } k=m,\\
-2^{k-m} & \text{jika } k<m,\\
0        & \text{jika } k>m.
\end{cases}
\end{equation*}
Hitung (atau tunjukkan bahwa limitnya tidak ada):
\\
a)~$\displaystyle \sum_{m=1}^\infty \sabs{a_{k,m}}~$ untuk setiap $k$,
%mbxSTARTIGNORE
\hspace{\fill}
%mbxENDIGNORE
%mbxlatex \qquad
b)~$\displaystyle \sum_{k=1}^\infty \sabs{a_{k,m}}~$ untuk setiap $m$,
%mbxSTARTIGNORE
\hspace{\fill}
%mbxENDIGNORE
%mbxlatex \qquad
c)~$\displaystyle \sum_{k=1}^\infty \sum_{m=1}^\infty \sabs{a_{k,m}}$,
%mbxSTARTIGNORE
\hspace{\fill}
%mbxENDIGNORE
%mbx <rabr/>
d)~$\displaystyle \sum_{k=1}^\infty \sum_{m=1}^\infty a_{k,m}$,
%mbxSTARTIGNORE
\hspace{\fill}
%mbxENDIGNORE
%mbxlatex \qquad
e)~$\displaystyle \sum_{m=1}^\infty \sum_{k=1}^\infty a_{k,m}$.
\\
Petunjuk: Teorema Fubini untuk deret tidak berlaku.
Bahkan, jawaban untuk d) dan e) berbeda.
\end{exercise}

\begin{exercise}
\pagebreak[2]
Misalkan $f(x) \coloneqq \frac{1}{1+x^2}$.  Buktikan bahwa
\begin{enumerate}[a)]
\item
$f$ merupakan fungsi analitik pada seluruh $\R$
dengan menemukan deret pangkat untuk $f$ di setiap $a \in \R$,
\item
jari-jari konvergensi dari
deret pangkat untuk $f$ di titik asal adalah 1.
\end{enumerate}
\end{exercise}


\begin{exercise}
\pagebreak[2]
Misalkan $f \colon \C \to \C$ analitik dan bukan fungsi nol.  Tunjukkan bahwa untuk setiap
$n$, hanya terdapat berhingga banyak titik nol dari $f$ di $B(0,n)$, yaitu,
$f^{-1}(0) \cap B(0,n)$ berhingga untuk setiap $n$.
\end{exercise}

\begin{exercise}
Misalkan $U \subset \C$ terbuka dan terhubung, $0 \in U$, dan $f \colon U
\to \C$ analitik.  Dengan memandang $f$ sebagai fungsi dari peubah real $x$
di titik asal, misalkan $f^{(n)}(0) = 0$ untuk semua $n$.  Tunjukkan bahwa $f(z) = 0$
untuk semua $z \in U$.
\end{exercise}

\begin{exercise}
Misalkan $U \subset \C$ terbuka dan terhubung, $0 \in U$, dan $f \colon U
\to \C$ analitik.  Untuk $x$ dan $y$ real yang cukup dekat dengan $0$ sehingga
$x \in U$ dan $iy \in U$, misalkan $h(x) \coloneqq f(x)$ dan
$g(y) \coloneqq -i \, f(iy)$.
Tunjukkan bahwa $h$ dan $g$ diferensiabel tak berhingga kali di titik asal dan
$h'(0) = g'(0)$.
\end{exercise}

\begin{exercise}
Misalkan suatu fungsi $f$ analitik pada suatu lingkungan dari titik asal,
dan terdapat suatu $M$
sedemikian sehingga $\sabs{f^{(n)}(0)} \leq M$ untuk semua $n$.
Buktikan bahwa deret pangkat $f$ di titik asal konvergen untuk semua $z \in \C$.
\end{exercise}

\begin{exercise}
Misalkan $f(z) \coloneqq \sum_{n=0}^\infty c_n z^n$ dengan jari-jari konvergensi 1.  Misalkan $f(0)
= 0$, tetapi $f$ bukan fungsi nol.
Tunjukkan bahwa terdapat suatu $k \in \N$ dan suatu
deret pangkat konvergen $g(z) \coloneqq \sum_{n=0}^\infty d_n z^n$ dengan jari-jari konvergensi 1
sedemikian sehingga $f(z) = z^k g(z)$ untuk semua $z \in B(0,1)$, dan $g(0) \neq 0$.
\end{exercise}

\begin{exercise}
Misalkan $U \subset \C$ terbuka dan terhubung.  Misalkan
$f \colon U \to \C$ analitik, $U \cap \R \neq \emptyset$, dan
$f(x) = 0$ untuk semua $x \in U \cap \R$.  Tunjukkan bahwa $f(z) = 0$ untuk semua $z \in
U$.
\end{exercise}

\begin{exercise}
Untuk $\alpha \in \C$ dan $k=0,1,2,3\ldots$, definisikan
\begin{equation*}
\binom{\alpha}{k} \coloneqq \frac{\alpha(\alpha-1)\cdots(\alpha-k+1)}{k!} .
\end{equation*}
\begin{enumerate}[a)]
\item
Tunjukkan bahwa deret
\begin{equation*}
f(z) \coloneqq \sum_{k=0}^\infty \binom{\alpha}{k} z^k
\end{equation*}
konvergen setiap kali $\sabs{z} < 1$.
Bahkan, buktikan bahwa untuk $\alpha = 0,1,2,3,\ldots$ jari-jari
konvergensinya adalah $\infty$, dan untuk semua $\alpha$ lainnya
jari-jari konvergensinya adalah 1.
\item
Tunjukkan bahwa untuk $x \in \R$, $\sabs{x} < 1$, kita mempunyai
\begin{equation*}
(1+x) f'(x) = \alpha f(x) ,
\end{equation*}
yang berarti bahwa $f(x) = (1+x)^\alpha$.
\end{enumerate}
\end{exercise}

\begin{exercise}
Misalkan $f \colon \C \to \C$ analitik dan misalkan untuk suatu
interval terbuka $(a,b) \subset \R$, $f$ bernilai real pada $(a,b)$.  Tunjukkan
bahwa $f$ bernilai real pada $\R$.
\end{exercise}

\begin{exercise}
Misalkan $\D \coloneqq B(0,1)$ adalah cakram satuan.  Misalkan
$f \colon \D \to \C$ analitik dengan deret pangkat
$\sum_{n=0}^\infty c_n z^n$.  Misalkan $\sabs{c_n} \leq 1$ untuk semua $n$.  Buktikan bahwa untuk
semua $z \in \D$, kita mempunyai
$\babs{f(z)} \leq \frac{1}{1-\sabs{z}}$.
\end{exercise}

%%%%%%%%%%%%%%%%%%%%%%%%%%%%%%%%%%%%%%%%%%%%%%%%%%%%%%%%%%%%%%%%%%%%%%%%%%%%%%

\sectionnewpage
\section{Eksponensial kompleks dan fungsi trigonometri}
\label{sec:complexexp}

%mbxINTROSUBSECTION

\sectionnotes{1 kuliah}

\subsection{Eksponensial kompleks}

Definisikan
\begin{equation*}
E(z) \coloneqq \sum_{k=0}^\infty \frac{1}{k!} z^k .
\end{equation*}
Deret ini konvergen untuk setiap $z \in \C$.
Jadi, menurut
\corref{cor:powerseranalytic}, $E$ analitik pada $\C$. Kita perhatikan bahwa $E(0) = 1$
dan, untuk $z=x \in \R$, berlaku $E(x) \in \R$. Dengan tetap mengambil $x$ real,
perhitungan langsung menunjukkan
\begin{equation*}
\frac{d}{dx} \bigl( E(x) \bigr) = E(x) .
\end{equation*}
Dalam \volIref{\sectionref*{vI-sec:logandexp} pada volume I}{\sectionref{sec:logandexp}} (atau dengan Teorema Picard), kita membuktikan bahwa
satu-satunya fungsi pada garis real yang memenuhi $E' = E$ dan
$E(0) = 1$ adalah fungsi eksponensial. Dengan kata lain, untuk $x \in \R$, $e^x = E(x)$.

Untuk bilangan kompleks $z$, kita definisikan\glsadd{not:complexexp}
\begin{equation*}
e^z \coloneqq E(z) = 
\sum_{k=0}^\infty \frac{1}{k!} z^k .
\end{equation*}
Pada garis real, definisi baru ini sama dengan definisi kita sebelumnya.
Lihat \figureref{fig:complexexpgraphs}. Perhatikan bahwa pada arah $x$
(arah real),
grafik berperilaku seperti eksponensial real, sedangkan pada arah $y$
(arah imajiner) grafik berosilasi.

\begin{myfigureht}
\myincludepdft{real_imag_exp}{%
Grafik permukaan.
Untuk x tetap, kedua grafik berbentuk sinusoidal pada arah y,
dengan amplitudo osilasi yang semakin besar ketika x semakin besar.
Untuk x negatif, grafik bahkan begitu kecil pada skala ini sehingga
osilasinya tidak dapat terlihat.
Pada grafik kanan, gelombang sinusoidal digeser. Pada grafik kiri,
irisan tepat di atas sumbu x naik secara eksponensial, sedangkan pada grafik
kanan irisan tepat di atas sumbu x selalu bernilai nol; dengan kata lain,
sumbu x termuat dalam grafik kanan. Irisan grafik di atas sumbu x ditandai
dengan garis tebal.}
\caption{Grafik bagian real (kiri) dan bagian imajiner (kanan)
dari eksponensial kompleks $e^z = e^{x+iy}$. Sumbu $x$ membentang dari $-4$ hingga
$4$, sumbu $y$ membentang dari $-6$ hingga $6$, dan sumbu vertikal membentang dari
$-e^{4} \approx -54.6$
hingga
$e^{4} \approx 54.6$. Grafik eksponensial real ($y=0$)
ditandai dengan garis tebal.\label{fig:complexexpgraphs}}
\end{myfigureht}

\begin{prop}[Hukum eksponen\index{hukum eksponen}]
Misalkan $z,w \in \C$ adalah bilangan kompleks. Maka
\begin{equation*}
e^{z+w} = e^z e^w.
\end{equation*}
\end{prop}

\begin{proof}
Kita sudah mengetahui bahwa kesamaan
$e^{x+y} = e^x e^y$ berlaku untuk semua
bilangan real $x$ dan $y$.
Untuk setiap $y \in \R$ yang tetap, pandang kedua ruas sebagai
fungsi dari $x$ dan terapkan teorema identitas
(\thmref{thm:identityanalytic}) untuk memperoleh
$e^{z+y} = e^ze^y$ bagi setiap $z \in \C$. Dengan menetapkan sembarang $z \in \C$,
kita memperoleh
$e^{z+y} = e^ze^y$ bagi setiap $y \in \R$. Sekali lagi menurut teorema identitas,
$e^{z+w} = e^z e^w$
bagi setiap $w \in \C$.
\end{proof}

Konsekuensi langsung dari proposisi ini adalah $e^z \neq 0$ untuk setiap $z \in \C$,
karena $e^z e^{-z} = e^{z-z} = 1$.
Perhitungan ini berarti bahwa ${(e^z)}^{-1} = e^{-z}$.
Dengan menggabungkan fakta tersebut dan hukum eksponen, kita memperoleh
\begin{equation*}
{(e^z)}^{n} = e^{nz} \qquad \text{untuk setiap } n \in \Z.
\end{equation*}
Konsekuensi lain yang sedikit lebih rumit adalah bahwa kita
dapat menghitung deret pangkat untuk fungsi eksponensial di sembarang titik $a \in
\C$:
\begin{equation*}
e^z = e^a e^{z-a} = \sum_{k=0}^\infty \frac{e^a}{k!} {(z-a)}^k .
\end{equation*}

\subsection{Fungsi trigonometri dan $\pi$}

Kini akhirnya kita dapat mendefinisikan \emph{\myindex{fungsi sinus}} dan \emph{\myindex{fungsi kosinus}}
melalui persamaan
\begin{equation*}
e^{x+iy} = e^x \bigl( \cos(y) + i \sin(y) \bigr) .
\end{equation*}
Lebih umum, kita mendefinisikan fungsi sinus dan fungsi kosinus untuk setiap bilangan kompleks $z$:
\glsadd{not:sin}\glsadd{not:cos}
\begin{equation*}
\cos(z) \coloneqq \frac{e^{iz} + e^{-iz}}{2}
\qquad\text{dan}\qquad
\sin(z) \coloneqq \frac{e^{iz} - e^{-iz}}{2i} .
\end{equation*}

Mari kita gunakan definisi ini untuk membuktikan sifat-sifat umum
fungsi sinus dan fungsi kosinus. Dalam prosesnya, kita juga mendefinisikan
bilangan $\pi$.

\begin{prop}
Fungsi sinus dan fungsi kosinus memiliki sifat-sifat berikut:
\begin{enumerate}[(i)]
\item Untuk setiap $z \in \C$,\index{rumus Euler}
\begin{equation*}
e^{iz} = \cos(z) + i\sin(z) \qquad
\text{(rumus Euler)}.
\end{equation*}
\item $\cos(0) = 1$, $\sin(0) = 0$.
\item Untuk setiap $z \in \C$,
\begin{equation*}
\cos(-z) = \cos(z), \qquad
\sin(-z) = -\sin(z).
\end{equation*}
\item Untuk setiap $z \in \C$,
\begin{equation*}
\cos(z) = \sum_{k=0}^\infty \frac{{(-1)}^k}{(2k)!} z^{2k} ,
\qquad
\sin(z) = \sum_{k=0}^\infty \frac{{(-1)}^k}{(2k+1)!} z^{2k+1} .
\end{equation*}
\item Untuk setiap $x \in \R$
\begin{equation*}
\cos(x) = \Re (e^{ix})
\qquad\text{dan}\qquad
\sin(x) = \Im (e^{ix}) .
\end{equation*}
\item Untuk setiap $z \in \C$,
\begin{equation*}
{\bigl( \cos(z) \bigr)}^2 + {\bigl( \sin(z) \bigr)}^2 = 1 .
\end{equation*}
\item Untuk setiap $x \in \R$,
\begin{equation*}
\babs{\sin(x)} \leq 1, \qquad \babs{\cos(x)} \leq 1 .
\end{equation*}
\item Untuk setiap $x \in \R$,
\begin{equation*}
\frac{d}{dx} \bigl[ \cos(x) \bigr] = -\sin(x)
\qquad \text{dan} \qquad
\frac{d}{dx} \bigl[ \sin(x) \bigr] = \cos(x) .
\end{equation*}
\item Untuk setiap $x \geq 0$,
\begin{equation*}
\sin(x) \leq x .
\end{equation*}
\item
Terdapat $x > 0$ sedemikian sehingga $\cos(x) = 0$. Kita definisikan
\glsadd{not:pi}
\begin{equation*}
\pi \coloneqq 2 \, \inf \{ x > 0 : \cos(x) = 0 \} .
\end{equation*}
\item
Untuk setiap $z \in \C$,
\begin{equation*}
e^{2\pi i} = 1 \qquad \text{dan} \qquad e^{z + i 2\pi} = e^z.
\end{equation*}
\item
Fungsi sinus dan fungsi kosinus periodik dengan periode $2\pi$ dan tidak memiliki periode positif
yang lebih kecil. Artinya, $2\pi$ adalah bilangan positif terkecil sedemikian sehingga untuk setiap $z \in \C$,
\begin{equation*}
\sin(z+2\pi) = \sin(z)
\qquad \text{dan} \qquad
\cos(z+2\pi) = \cos(z) .
\end{equation*}
\item
Fungsi $x \mapsto e^{ix}$ merupakan pemetaan bijektif dari $[0,2\pi)$
ke himpunan semua $z \in \C$ sedemikian sehingga $\sabs{z} = 1$.
\end{enumerate}
\end{prop}

Proposisi ini langsung menyiratkan bahwa $\sin(x)$ dan $\cos(x)$ bernilai
real apabila $x$ real.

\begin{proof}
Tiga butir pertama langsung mengikuti definisi.
Perhitungan deret pangkat untuk kedua fungsi tersebut diserahkan sebagai latihan.
Karena pemetaan konjugasi kompleks bersifat kontinu, definisi
$e^z$ menyiratkan
$\overline{e^z} = e^{\bar{z}}$. Jika
$x$ real,
\begin{equation*}
\overline{e^{ix}} = e^{-ix} .
\end{equation*}
Jadi, untuk $x$ real,
$\cos(x) =
\frac{e^{ix}+e^{-ix}}{2} =
\frac{e^{ix}+\overline{e^{ix}}}{2} =
\Re (e^{ix})$
dan, secara serupa, $\sin(x) = \Im (e^{ix})$.

Untuk $x$ real, kita hitung
\begin{equation*}
1 =  e^{ix} e^{-ix}
= e^{ix} \, \overline{e^{ix}}
= \sabs{e^{ix}}^2
= \babs{\cos(x) + i \sin(x)}^2
= {\bigl( \cos(x) \bigr)}^2 + {\bigl( \sin(x) \bigr)}^2 .
\end{equation*}
Perhitungan yang sedikit lebih rumit menunjukkan fakta ini untuk
bilangan kompleks; lihat \exerciseref{exercise:cossinidentity}.
Khususnya, $e^{ix}$ bersifat \emph{\myindex{bermodulus satu}} untuk $x$ real;
nilai-nilainya terletak pada lingkaran satuan.
Kuadrat suatu bilangan real selalu nonnegatif:
\begin{equation*}
{\bigl(\sin(x)\bigr)}^2 = 1-{\bigl(\cos(x)\bigr)}^2 \leq 1 .
\end{equation*}
Jadi, $\babs{\sin(x)} \leq 1$ dan, secara serupa,
$\babs{\cos(x)} \leq 1$.

Perhitungan kedua turunan tersebut kita serahkan kepada pembaca sebagai latihan.
Mari kita buktikan bahwa $\sin(x) \leq x$ untuk $x \geq 0$.
Pandang
$f(x) \coloneqq x-\sin(x)$ dan turunkan:
\begin{equation*}
f'(x) = \frac{d}{dx} \bigl[ x - \sin(x) \bigr]
=
1 -\cos(x) \geq 0 ,
\end{equation*}
untuk setiap $x \in \R$ karena $\babs{\cos(x)} \leq 1$.
Dengan kata lain, $f$ monoton naik dan $f(0) = 0$.
Jadi, $f$ nonnegatif ketika $x \geq 0$; oleh karena itu, $\sin(x) \leq x$.

Selanjutnya, kita klaim bahwa terdapat $x$ positif sedemikian sehingga $\cos(x) = 0$.
Karena $\cos(0) = 1 > 0$, berlaku $\cos(x) > 0$
untuk $x$ yang cukup dekat dengan $0$. Definisikan
$S \coloneqq \{ t > 0 : \cos(x) > 0 \text{ untuk setiap } x \in [0,t) \}$.
Himpunan $S$ tak kosong. Pilih $y_0 \in S$ dan $a \in (0,y_0)$.
Karena $\cos(x) > 0$ pada $[0,y_0)$, fungsi $\sin(x)$ naik tegas
pada $[0,y_0)$. Karena $\sin(0) = 0$, kita memperoleh $\sin(a) > 0$.
Sekarang, ambil sembarang $y \in S$ dengan $y > a$. Pada $[0,y)$,
fungsi $\sin$ naik tegas. Menurut teorema nilai rata-rata,
terdapat $c \in (a,y)$ sedemikian sehingga
\begin{equation*}
2 \geq \cos(a)-\cos(y) = \sin(c)(y-a) \geq \sin(a)(y-a) .
\end{equation*}
Karena $\sin(a) > 0$, kita memperoleh
\begin{equation*}
y \leq \frac{2}{\sin(a)} + a .
\end{equation*}
Batas ini berlaku bagi setiap $y \in S$ dengan $y > a$, sedangkan anggota $S$
yang memenuhi $y \leq a$ sudah terbatas. Jadi, $S$ terbatas di atas. Definisikan
$y \coloneqq \sup S$. Dari sifat supremum, $\cos(x) > 0$ untuk setiap $x \in [0,y)$,
sehingga kekontinuan memberikan $\cos(y) \geq 0$. Jika $\cos(y) > 0$, kekontinuan
akan memberikan suatu $t > y$ dengan $t \in S$, bertentangan dengan definisi supremum.
Oleh karena itu, $\cos(y) = 0$. Karena $\cos(x) > 0$ untuk setiap $x \in [0,y)$,
$y$ adalah nilai positif terkecil sedemikian sehingga $\cos(y) = 0$. Seperti telah disebutkan,
$\pi$ didefinisikan sebagai $2y$.

Karena $\cos(\nicefrac{\pi}{2}) = 0$, kita peroleh
${\bigl(\sin(\nicefrac{\pi}{2})\bigr)}^2 = 1$.
Karena $\sin$ positif pada $(0,\nicefrac{\pi}{2})$, kita mempunyai
$\sin(\nicefrac{\pi}{2}) = 1$.
Dengan demikian,
\begin{equation*}
e^{i \pi /2} = i ,
\end{equation*}
dan, menurut hukum eksponen,
\begin{equation*}
e^{i \pi} = -1 ,
\qquad 
e^{i 2\pi} = 1 .
\end{equation*}
Jadi, $e^{i2\pi} = 1 = e^0$. Hukum eksponen juga menyatakan
\begin{equation*}
e^{z+i2\pi} = e^z e^{i2\pi} = e^z
\end{equation*}
untuk setiap $z \in \C$. Kita juga langsung memperoleh
$\cos(z+2\pi) = \cos(z)$ dan $\sin(z+2\pi) = \sin(z)$.
Jadi, $\sin$ dan $\cos$ periodik dengan periode $2\pi$.

Kita klaim bahwa $\sin$ dan $\cos$ tidak periodik dengan periode positif yang lebih kecil.
Untuk itu, ambil sembarang $x$ yang memenuhi $0 < x \leq 2\pi$ dan
$e^{ix} = 1$. Kita akan menunjukkan bahwa $x = 2\pi$.
Menurut hukum eksponen,
\begin{equation*}
{\bigl(e^{ix/4}\bigr)}^4 = 1 .
\end{equation*}
Jika $e^{ix/4} = a+ib$, maka
\begin{equation*}
{(a+ib)}^4
=a^4-6a^2b^2+b^4 + i\bigl(4ab(a^2-b^2)\bigr)
=1 .
\end{equation*}
Karena $0 < \nicefrac{x}{4} \leq \nicefrac{\pi}{2}$, kita mempunyai
$b = \sin(\nicefrac{x}{4}) > 0$.
Selain itu, $a = \cos(\nicefrac{x}{4}) \geq 0$. Bagian imajiner pada persamaan di atas
harus nol. Karena $b > 0$, maka $a = 0$ atau $a^2 = b^2$.
Jika $a^2=b^2$, maka
$a^4-6a^2b^2+b^4 = -4a^4 < 0$, dan khususnya tidak sama dengan 1.
Oleh karena itu, $a=0$; dalam hal ini $\nicefrac{x}{4} = \nicefrac{\pi}{2}$.
Jadi, $x = 2\pi$.

Sekarang, misalkan $T > 0$ merupakan periode fungsi sinus. Maka, untuk setiap $x \in \R$,
$\sin(x+T)=\sin(x)$. Dengan menurunkan identitas ini terhadap $x$, kita memperoleh
$\cos(x+T)=\cos(x)$. Jika $T$ merupakan periode fungsi kosinus, penalaran yang sama,
dengan menurunkan identitas $\cos(x+T)=\cos(x)$ terhadap $x$, menghasilkan
$\sin(x+T)=\sin(x)$. Dengan demikian, setiap periode positif dari salah satu fungsi
memberikan, pada $x=0$, $\sin(T)=\sin(0)$ dan $\cos(T)=\cos(0)$.
Menurut rumus Euler, periode tersebut memenuhi
$e^{iT}=\cos(T)+i\sin(T)=\cos(0)+i\sin(0)=1$.
Hasil di atas menunjukkan bahwa tidak ada periode positif yang lebih kecil daripada $2\pi$.

Terakhir, kita ingin menunjukkan bahwa $e^{ix}$ merupakan pemetaan bijektif
dari himpunan $[0,2\pi)$ ke himpunan semua $z \in \C$ sedemikian sehingga
$\sabs{z} = 1$. Misalkan $e^{ix} = e^{iy}$ dan
$x > y$. Maka,
$0 < x-y < 2\pi$ karena $x,y \in [0,2\pi)$. Akan tetapi, $e^{i(x-y)} = 1$,
bertentangan dengan hasil di atas. Jadi, pemetaan tersebut injektif.
Untuk menunjukkan bahwa pemetaan $x \mapsto e^{ix}$ surjektif, ambil $(a,b) \in \R^2$ sedemikian sehingga $a^2+b^2 = 1$.
Pertama, andaikan $a,b \geq 0$. Menurut teorema nilai antara,
terdapat $x \in [0,\nicefrac{\pi}{2}]$ sedemikian sehingga
$\cos(x) = a$, dan akibatnya $b^2 = \bigl(\sin(x)\bigr)^2$. Karena
$b$ dan $\sin(x)$ nonnegatif, berlaku $b = \sin(x)$.
Karena $-\sin(x)$ merupakan turunan $\cos(x)$
dan $\cos(-x) = \cos(x)$,
kita memperoleh $\sin(x) < 0$ untuk $x \in [\nicefrac{-\pi}{2},0)$.
Dengan penalaran yang sama, kita memperoleh bahwa
jika $a > 0$ dan $b < 0$, kita dapat menemukan $x_0$ dalam $(\nicefrac{-\pi}{2},0)$
sedemikian sehingga $\cos(x_0) = a$ dan $\sin(x_0)=b$. Menurut periodisitas,
$x \coloneqq x_0+2\pi \in (\nicefrac{3\pi}{2},2\pi)$ juga memenuhi
$\cos(x) = a$ dan $\sin(x)=b$. Kasus $(a,b)=(0,-1)$ diperoleh dengan $x=\nicefrac{3\pi}{2}$.
Mengalikan dengan $-1$ sama dengan mengalikan dengan $e^{i\pi}$ atau
$e^{-i\pi}$. Jadi, kita selalu dapat mengasumsikan bahwa $a \geq 0$ (rincian diserahkan
sebagai latihan).
\end{proof}

\subsection{Lingkaran satuan dan koordinat polar}

Jika $\gamma \colon [a,b] \to \C$ merupakan kurva mulus sepotong-sepotong,
panjang busurnya diberikan oleh
\begin{equation*}
\int_a^b \babs{\gamma'(t)} \, dt .
\end{equation*}
Pemetaan $t \mapsto e^{it}$, untuk $t \in [0,2\pi)$, merupakan parametrisasi
bijektif lingkaran satuan. Untuk menghitung keliling, kita gunakan parametrisasi
satu putaran penuh pada interval tertutup $[0,2\pi]$. Karena
$\frac{d}{dt} \bigl( e^{it} \bigr) = ie^{it}$, keliling
lingkaran (panjang busurnya) adalah
\begin{equation*}
\int_0^{2\pi} \sabs{i e^{it}}  \,  dt
=
\int_0^{2\pi} 1  \,  dt  = 2\pi .
\end{equation*}

Lebih umum, pemetaan $t \mapsto e^{it}$ merupakan parametrisasi panjang busur bagi lingkaran.
Artinya, $t$, sebagai ukuran sudut dalam radian, mengukur panjang busur pada
lingkaran berjari-jari 1. Jadi, definisi $\sin$ dan $\cos$ yang diberikan
di atas sesuai dengan definisi geometris standar.

Setiap titik pada lingkaran satuan dapat dituliskan sebagai $e^{it}$ untuk suatu $t$.
Oleh karena itu, setiap bilangan kompleks $z \in \C$ dapat dituliskan dalam
\emph{\myindex{koordinat polar}} sebagai
\begin{equation*}
z = r e^{i\theta}
\end{equation*}
untuk suatu $r \geq 0$ dan $\theta \in \R$. Tentu saja, $\theta$
tidak tunggal karena $\theta$ dan $\theta+2\pi$ menghasilkan bilangan yang sama.
Hukum eksponen $e^{a+b} = e^a e^b$ menghasilkan rumus yang berguna untuk perpangkatan
dan hasil kali bilangan kompleks dalam koordinat polar. Rumus perpangkatan berlaku untuk
$n \in \N$, dan juga untuk $n \in \Z$ apabila $r > 0$:
\begin{equation*}
{(r e^{i\theta})}^n
= r^n e^{i n \theta} ,
\qquad
(r e^{i\theta})
(s e^{i\gamma})
=
rs e^{i(\theta+\gamma)} .
\end{equation*}

\subsection{Latihan}

\begin{exercise}
Turunkan deret pangkat untuk $\sin(z)$ dan $\cos(z)$ di titik asal.
\end{exercise}

\begin{exercise}
Dengan menggunakan deret pangkat, tunjukkan bahwa untuk $x$ real berlaku
$\frac{d}{dx} \bigl[ \sin(x)\bigr] = \cos(x)$ dan
$\frac{d}{dx} \bigl[ \cos(x)\bigr] = -\sin(x)$.
\end{exercise}

\begin{exercise}
Lengkapi bukti bahwa pemetaan $x \mapsto e^{ix}$ dari
$[0,2\pi)$ ke lingkaran satuan bersifat surjektif. Secara khusus, andaikan bahwa
kita telah memperoleh semua titik berbentuk $(a,b)$ dengan $a^2+b^2=1$ dan $a \geq 0$.
Dengan mengalikan dengan $e^{i\pi}$ atau $e^{-i\pi}$, tunjukkan bahwa kita memperoleh
semua titik, yakni juga titik-titik dengan $a < 0$.
\end{exercise}

\begin{exercise}
Buktikan bahwa fungsi eksponensial kompleks $z \mapsto e^z$ surjektif dari $\C$
ke $\C \setminus \{ 0 \}$, dan bahkan
bahwa untuk setiap $w$ tak nol, terdapat tak berhingga banyak $z \in \C$ sedemikian sehingga $e^z=w$.
\end{exercise}

\begin{exercise}
Buktikan bahwa untuk setiap $w \neq 0$ dan setiap $\epsilon > 0$,
terdapat $z \in \C$ dengan $\sabs{z} < \epsilon$ sedemikian sehingga $e^{1/z} = w$.
\end{exercise}

\begin{exercise}\label{exercise:cossinidentity}
Kita telah menunjukkan bahwa
${\bigl( \cos(x) \bigr)}^2 + {\bigl( \sin(x) \bigr)}^2 = 1$
untuk setiap $x \in \R$.
Buktikan bahwa
${\bigl( \cos(z) \bigr)}^2 + {\bigl( \sin(z) \bigr)}^2 = 1$
untuk setiap $z \in \C$.
\end{exercise}

\begin{exercise}
Buktikan identitas trigonometri
$\sin(z + w) = \sin(z) \cos(w) + \cos(z) \sin(w)$ dan
$\cos(z + w) = \cos(z) \cos(w) - \sin(z) \sin(w)$ untuk setiap $z,w \in \C$.
\end{exercise}

\begin{exercise}
Definisikan $\operatorname{sinc}(z) \coloneqq \frac{\sin(z)}{z}$ untuk $z \neq 0$ dan
$\operatorname{sinc}(0) \coloneqq 1$.
Tunjukkan bahwa fungsi sinc bersifat analitik dan hitung deret pangkatnya di titik asal.
\end{exercise}

\begin{exnote}
\pagebreak[2]
Definisikan \emph{\myindex{sinus hiperbolik}} dan
\emph{\myindex{kosinus hiperbolik}} melalui
\begin{equation*}
\sinh(z) \coloneqq \frac{e^z-e^{-z}}{2}, \qquad
\cosh(z) \coloneqq \frac{e^z+e^{-z}}{2}.
\end{equation*}
\end{exnote}

\begin{exercise}
Turunkan deret pangkat fungsi sinus hiperbolik dan fungsi kosinus hiperbolik di titik asal.
\end{exercise}

\begin{exercise}
Tunjukkan bahwa
\begin{enumerate}[a)]
\item
$\sinh(0) = 0$, $\cosh(0) = 1$.
\item
Untuk $x \in \R$,
$\frac{d}{dx} \bigl[ \sinh(x) \bigr] = \cosh(x)$ dan
$\frac{d}{dx} \bigl[ \cosh(x) \bigr] = \sinh(x)$.
\item
$\cosh(x) > 0$ untuk setiap $x \in \R$ dan
pemetaan $x \mapsto \sinh(x)$ naik tegas dan bijektif dari $\R$ ke $\R$.
\item
${\bigl(\cosh(x)\bigr)}^2 = 1 + {\bigl(\sinh(x)\bigr)}^2$ untuk setiap
$x$.
\end{enumerate}
\end{exercise}

\begin{exercise}
Definisikan $\tan(x) \coloneqq \frac{\sin(x)}{\cos(x)}$ untuk $x$ dengan $\cos(x) \neq 0$, seperti biasa.
\begin{enumerate}[a)]
\item
Tunjukkan bahwa untuk $x \in (\nicefrac{-\pi}{2},\nicefrac{\pi}{2})$,
baik $\sin$ maupun $\tan$ naik tegas, sehingga
fungsi invers $\arcsin$
dan $\arctan$ ada apabila kedua fungsi tersebut dibatasi pada interval tersebut.
\item
Tunjukkan bahwa $\arcsin$ diferensiabel pada $(-1,1)$ dan $\arctan$
diferensiabel pada $\R$, serta bahwa
$\frac{d}{dx} \arcsin(x) = \frac{1}{\sqrt{1-x^2}} \quad (-1<x<1)$ dan
$\frac{d}{dx} \arctan(x) = \frac{1}{1+x^2} \quad (x \in \R)$.
\item
Dengan menggunakan rumus jumlah geometri berhingga, tunjukkan bahwa deret pada
ruas kanan persamaan berikut konvergen untuk setiap $-1 \leq x \leq 1$
(termasuk titik-titik ujung), dan bahwa
\begin{equation*}
\arctan(x) = \int_0^x \frac{1}{1+t^2} dt
=
\sum_{k=0}^\infty \frac{{(-1)}^k}{2k+1} x^{2k+1}
\end{equation*}
Petunjuk: Integralkan jumlah berhingga tersebut, bukan deretnya.
\item
Gunakan hasil ini untuk menunjukkan bahwa
\begin{equation*}
1 - \frac{1}{3} + \frac{1}{5} - \cdots
=
\sum_{k=0}^\infty \frac{{(-1)}^k}{2k+1}
=
\frac{\pi}{4} .
\end{equation*}
\end{enumerate}
\end{exercise}

%%%%%%%%%%%%%%%%%%%%%%%%%%%%%%%%%%%%%%%%%%%%%%%%%%%%%%%%%%%%%%%%%%%%%%%%%%%%%%

\sectionnewpage
\section[Maximum principle and fundamental theorem of algebra]{Maximum principle and the fundamental theorem of algebra}
\label{sec:fundalgeb}

\sectionnotes{half a lecture, optional}

In this section we study the local behavior of polynomials,
and analytic functions in general,
and the growth of polynomials as $z$ goes to infinity.
As an application, we prove the fundamental theorem of algebra:
Any nonconstant polynomial has a complex root.
We consider polynomials in a complex variable $z$
and so we may emphasize their domain by calling them complex polynomials.

\begin{lemma} \label{lemma:polyalwaysgetssmaller}
Let $\epsilon > 0$,
let $p(z)$ be a nonconstant complex polynomial, or more generally
a nonconstant power series converging in $B(z_0,\epsilon)$,
and suppose $p(z_0) \neq 0$.
Then there exists a $w \in B(z_0,\epsilon)$ such that $\babs{p(w)} < \babs{p(z_0)}$.
\end{lemma}

\begin{proof}
We prove this lemma for a polynomial and leave the general case as
\exerciseref{exercise:minprinciple}.
Without loss of generality assume that $z_0 = 0$ and $p(0) = 1$.  Write
\begin{equation*}
p(z) = 1+a_kz^k + a_{k+1}z^{k+1} + \cdots + a_d z^d ,
\end{equation*}
where $a_k \neq 0$.  Pick $t$ such that $a_k e^{ikt} = -\sabs{a_k}$, which
we can do by the discussion on trigonometric functions.  Suppose
$r > 0$ is small enough such that
$1-r^k \sabs{a_k} > 0$.  We have
\begin{equation*}
p(r e^{it}) =
1-r^k \sabs{a_k} + r^{k+1}a_{k+1}e^{i(k+1)t} + \cdots + r^{d}a_{d}e^{idt} .
\end{equation*}
So
\begin{equation*}
\begin{split}
%mbxlatex &
\babs{
p(r e^{it}) } - \abs{
r^{k+1}a_{k+1}e^{i(k+1)t} + \cdots + r^{d}a_{d}e^{idt}
}
%mbxlatex \\
&
%mbxlatex \qquad \qquad \qquad \qquad
\leq
\babs{
p(r e^{it}) 
- r^{k+1}a_{k+1}e^{i(k+1)t} - \cdots - r^{d}a_{d}e^{idt}
}
\\
&
%mbxlatex \qquad \qquad \qquad \qquad
=
\babs{
1-r^k \sabs{a_k}
}
=
1-r^k \sabs{a_k} .
\end{split}
\end{equation*}
In other words,
\begin{equation*}
\babs{
p(r e^{it}) }
\leq
1-r^k \left( \sabs{a_k}
-
r
\abs{
a_{k+1}e^{i(k+1)t} + \cdots + r^{d-k-1}a_{d}e^{idt}
}
\right) .
\end{equation*}
For small enough $r$, the expression in the parentheses is positive
as $\sabs{a_k} > 0$.  Hence, $\babs{p(re^{it})} < 1 = p(0)$.
\end{proof}

What the lemma
says is that the only minima the modulus of analytic functions
has are precisely at the zeros.
It is sometimes called the \emph{\myindex{minimum modulus principle}}.
If $f$ is analytic and nonzero at a point,
then $\nicefrac{1}{f}$ is analytic near that point.  Applying the lemma and
the identity theorem, one obtains the \emph{maximum modulus principle}, or sometimes
just the \emph{maximum principle}.

\begin{thm}[Maximum modulus principle]\index{maximum modulus principle}%
\index{maximum principle!analytic functions}%
\label{thm:maxprinciple}
If $U \subset \C$ is open and connected,
$f \colon U \to \C$ is analytic, and $\babs{f(z)}$ attains a relative
maximum at $z_0 \in U$, then $f$ is constant.
\end{thm}

The details of the proof are left as \exerciseref{exercise:maxprinciple}.

\begin{remark}
The lemma (and the maximum principle) does not hold if we restrict to the real numbers.  For
example, $x^2+1$ has a minimum at $x=0$, but no zero there.
There is a $w$ arbitrarily close to $0$ such that $\sabs{w^2+1} < 1$, but this
$w$ is necessarily not real.  Letting $w = i\epsilon$ for small
$\epsilon > 0$ works.
\end{remark}

The moral of the story is that if $p(0) = 1$, then very close to 0, the
series (or polynomial)
looks like $1+az^k$, and $1+az^k$ has no minimum at the origin.  All the higher
powers of $z$ are too small to make a difference.  For polynomials, we find similar behavior
at infinity.

\begin{lemma}
Let $p(z)$ be a nonconstant complex polynomial.  Then for any $M > 0$, there exists
an $R > 0$ such that
$\babs{p(z)} \geq M$ whenever $\sabs{z} \geq R$.
\end{lemma}

\begin{proof}
Write $p(z) = a_0 + a_1 z + \cdots + a_d z^d$ and suppose that $d \geq 1$
and $a_d \neq 0$.
Suppose $\sabs{z} \geq R$ (so also $\sabs{z}^{-1} \leq R^{-1}$).
We estimate:
\begin{equation*}
\begin{split}
\babs{p(z)}
& \geq
\sabs{a_d z^d} -
\sabs{a_0} - \sabs{a_1 z} - \cdots - \sabs{a_{d-1} z^{d-1} }
\\
& =
\sabs{z}^d \bigl(
\sabs{a_d} -
\sabs{a_0} \, \sabs{z}^{-d} -
\sabs{a_1} \, \sabs{z}^{-d+1} - \cdots - \sabs{a_{d-1}} \, \sabs{z}^{-1}
\bigr)
\\
& \geq
R^d \bigl(\sabs{a_d} -
\sabs{a_0}R^{-d} - \sabs{a_1}R^{1-d} - \cdots - \sabs{a_{d-1}}R^{-1} \bigr)
.
\end{split}
\end{equation*}
Then the expression in parentheses is eventually positive for large enough
$R$.  In particular, for large enough $R$, this expression
is greater than
$\frac{\sabs{a_d}}{2}$, and so
\begin{equation*}
\babs{p(z)}
\geq
R^d \frac{\sabs{a_d}}{2} .
\end{equation*}
Therefore,
we can pick $R$ large enough to be bigger than a given $M$.
\end{proof}

This second lemma does \emph{not} generalize to analytic
functions, even those defined on the entire plane $\C$.  The function
$\cos(z)$ is a counterexample.
We had to look
at the term with the largest degree, and we only have such a term for
a polynomial.  In fact, something that we will not prove is that
an analytic function defined on all of $\C$ satisfying the conclusion
of the lemma must be a polynomial.

The moral of the story here is that for very large $\sabs{z}$ (far away from
the origin) a polynomial of degree $d$ really looks like a constant multiple
of $z^d$.

\begin{thm}[Fundamental theorem of algebra]
\index{fundamental theorem of algebra}%
Let $p(z)$ be a nonconstant complex polynomial, then there exists a $z_0 \in \C$
such that $p(z_0) = 0$.
\end{thm}

\begin{proof}
Let $\mu \coloneqq \inf \bigl\{ \babs{p(z)} : z \in \C \bigr\}$.  Find an $R$ such that
for all $z$ with $\sabs{z} \geq R$, we have $\babs{p(z)} \geq \mu+1$.
Therefore, every $z$ with $\babs{p(z)}$ close to $\mu$ must be in the
closed ball $C(0,R) = \bigl\{ z \in \C : \sabs{z} \leq R \bigr\}$.  As $\babs{p(z)}$
is a continuous real-valued function, it achieves its minimum
on the compact set $C(0,R)$ (closed and bounded) and this minimum must
be $\mu$.  So there is a $z_0 \in C(0,R)$ such that $\babs{p(z_0)} = \mu$.
As that is a minimum of $\babs{p(z)}$ on $\C$, by the first lemma
above, we have $\babs{p(z_0)} = 0$.
\end{proof}

The fundamental theorem also does not generalize to analytic functions.  The
exponential $e^{z}$ is an analytic function on $\C$ with no zeros.

\subsection{Exercises}

\begin{exercise} \label{exercise:minprinciple}
Prove \lemmaref{lemma:polyalwaysgetssmaller} for an analytic function.  That
is, suppose that $p(z)$ is a nonconstant power series converging in
$B(z_0,\epsilon)$.
\end{exercise}

\begin{exercise} \label{exercise:maxprinciple}
Use \lemmaref{lemma:polyalwaysgetssmaller} for analytic functions to prove \thmref{thm:maxprinciple}.
\end{exercise}

\begin{exercise}
Let $U \subset \C$ be open and $z_0 \in U$.
Suppose $f \colon U \to \C$ is analytic and $f(z_0) = 0$.  Show that
there exists an $\epsilon > 0$ such that either
$f(z) \neq 0$ for all $z$ with $0 < \sabs{z} < \epsilon$
or $f(z) = 0$ for all $z \in B(z_0,\epsilon)$.
In other words, zeros of analytic functions are isolated.
Of course, same holds for polynomials.
\end{exercise}

\begin{exnote}
\pagebreak[1]
A \emph{\myindex{rational function}} is a function
$f(z) \coloneqq \frac{p(z)}{q(z)}$
where $p$ and $q$ are polynomials and $q$ is not identically zero.
A point $z_0 \in \C$ where $f(z_0) = 0$ (and therefore $p(z_0) = 0$)
is called a \emph{zero}\index{zero of a function}.
A point $z_0 \in \C$ is called a \emph{\myindex{singularity}} of
$f$ if $q(z_0) = 0$.  As zeros of polynomials are isolated,
singularities of rational functions are isolated
and so are called
\emph{isolated singularities}\index{isolated singularity}.
An isolated singularity is called
\emph{removable}\index{removable singularity}
if $\lim_{z \to z_0} f(z)$ exists.
An isolated singularity is called a \emph{\myindex{pole}} if 
$\lim_{z \to z_0} \babs{f(z)} = \infty$.
We say $f$ has a pole at $\infty$ if
\begin{equation*}
\lim_{z \to \infty} \babs{f(z)} = \infty ,
\end{equation*}
that is, if for every $M > 0$ there exists an $R > 0$ such that
$\babs{f(z)} > M$ for all $z$ with $\sabs{z} > R$.
\end{exnote}

\begin{exercise}
Show that a rational function that is not identically
zero has at most finitely many zeros and
singularities.  In fact, show that if $p$ is a polynomial of 
degree $n > 0$ it has at most $n$ zeros.
\\
Hint: If $z_0$ is a zero of $p$, then without loss of generality assume $z_0 =
0$.  Then use induction.
\end{exercise}

\begin{exercise}
Prove that if $z_0$ is a removable singularity of a rational
function $f(z) \coloneqq \frac{p(z)}{q(z)}$, then there exist
polynomials $\widetilde{p}$ and $\widetilde{q}$ such that
$\widetilde{q}(z_0) \neq 0$ and $f(z) =
\frac{\widetilde{p}(z)}{\widetilde{q}(z)}$.
\\
Hint: Without loss of generality assume $z_0 = 0$.
\end{exercise}

\begin{exercise}
Given a rational function $f$ with an isolated singularity at $z_0$,
show that $z_0$ is either removable or a pole.
\\
Hint: See the previous exercise.
\end{exercise}

\begin{exercise}
Let $f$ be a rational function and $S \subset \C$ be the 
set of the singularities of $f$.
Prove that $f$ is equal to a nonconstant polynomial on $\C \setminus S$
if and only if
$f$ has a pole at infinity and all the singularities are removable.
\\
Hint: See exercises above.
\end{exercise}



%%%%%%%%%%%%%%%%%%%%%%%%%%%%%%%%%%%%%%%%%%%%%%%%%%%%%%%%%%%%%%%%%%%%%%%%%%%%%%

\sectionnewpage
\section{Equicontinuity and the Arzel\`a--Ascoli theorem}
\label{sec:arzelaascoli}

\sectionnotes{2 lectures}

We would like an analogue of Bolzano--Weierstrass.  Something to the tune of
\myquote{every bounded
sequence of functions (with some property) has a convergent subsequence.}
Matters are not
as simple even for continuous functions. 
Not every bounded sequence in the metric space $C\bigl([0,1],\R\bigr)$ has
a convergent subsequence.

\begin{defn}
Let $X$ be a set and
$f_n \colon X \to \C$ be a sequence of functions.  We say that
$\{ f_n \}_{n=1}^\infty$
is \emph{\myindex{pointwise bounded}} if for every $x \in X$, there is an $M_x \in \R$
such that
\begin{equation*}
\babs{f_n(x)} \leq M_x \qquad \text{for all } n \in \N .
\end{equation*}
We say that
$\{ f_n \}_{n=1}^\infty$
is \emph{\myindex{uniformly bounded}} if there is an $M \in \R$
such that
\begin{equation*}
\babs{f_n(x)} \leq M \qquad \text{for all } n \in \N \text{ and all } x \in X.
\end{equation*}
\end{defn}

If $X$ is a compact metric space, then a sequence in $C(X,\C)$
is uniformly bounded if it is bounded as a set in the metric space
$C(X,\C)$ using the uniform norm.

\begin{example}
There exist sequences of 
continuous functions
on $[0,1]$ that are uniformly bounded but contain no subsequence converging
even pointwise.
Let us state without proof that $f_n(x) \coloneqq \sin (2\pi n x)$ is one
such sequence.
Below we will show that there must always exist
a subsequence converging at countably
many points, but $[0,1]$ is uncountable.
\end{example}

\begin{example}
The sequence $f_n(x) \coloneqq x^n$ of continuous functions on $[0,1]$
is uniformly bounded, but contains no subsequence that converges
uniformly,
although the sequence converges pointwise (to a discontinuous function).
\end{example}

\begin{example}
The sequence $\{ f_n \}_{n=1}^\infty$ of functions in $C\bigl([0,1],\R\bigr)$ given by
$f_n(x) \coloneqq \frac{n^3x}{1+n^4x^2}$
converges pointwise to the zero function (obvious at $x=0$, and for $x > 0$,
we have $\frac{n^3x}{1+n^4x^2} \leq \frac{1}{nx}$).
For each $x$, $\{f_n(x)\}_{n=1}^\infty$ is bounded as it converges to 0.
So $\{ f_n \}_{n=1}^\infty$ is pointwise bounded.

Via calculus, we find that the maximum of
$f_n$ on
$[0,1]$ occurs at the critical point
$x=\nicefrac{1}{n^2}$:
\begin{equation*}
\snorm{f_n}_{[0,1]}
=
f_n\left(\nicefrac{1}{n^2}\right)
= \nicefrac{n}{2} .
\end{equation*}
So $\lim_{n\to\infty} \snorm{f_n}_{[0,1]} = \infty$, and
this sequence is not uniformly bounded.
\end{example}

When the domain is countable, we can locate a subsequence
converging at least pointwise.
The proof uses a very common and useful diagonal argument.

\begin{prop} \label{prop:subsequenceoncountableX}
Let $X$ be a countable set and $f_n \colon X \to \C$ give a pointwise bounded
sequence of functions.  Then $\{ f_n \}_{n=1}^\infty$ has a subsequence that converges
pointwise.
\end{prop}

\begin{proof}
Let $x_1,x_2,x_3,\ldots$ be an enumeration of the elements of $X$.
The sequence $\{ f_n(x_1) \}_{n=1}^\infty$ is bounded and hence
we have a subsequence of $\{ f_n \}_{n=1}^{\infty}$, which we denote by
$\{ f_{1,k} \}_{k=1}^\infty$,
such that
$\{ f_{1,k}(x_1) \}_{k=1}^\infty$ converges.
Next $\{ f_{1,k}(x_2) \}_{k=1}^\infty$ is bounded and so 
$\{ f_{1,k} \}_{k=1}^\infty$ has a subsequence
$\{ f_{2,k} \}_{k=1}^\infty$ such that
$\{ f_{2,k}(x_2) \}_{k=1}^\infty$ converges.  Note that
$\{ f_{2,k}(x_1) \}_{k=1}^\infty$ is still convergent.

In general, we have a sequence $\{ f_{m,k} \}_{k=1}^\infty$,
which is a subsequence of $\{ f_{m-1,k} \}_{k=1}^\infty$,
such that $\{ f_{m,k}(x_j) \}_{k=1}^\infty$ converges for $j=1,2,\ldots, m$.
We let $\{ f_{m+1,k} \}_{k=1}^\infty$ be a subsequence of
$\{ f_{m,k} \}_{k=1}^\infty$
such that
$\{ f_{m+1,k}(x_{m+1}) \}_{k=1}^\infty$ converges (and hence it converges for all
$x_j$ for $j=1,2,\ldots,m+1$).  Rinse and repeat.

If $X$ is finite, we are done as the process stops at some point.
If $X$ is countably infinite,
we pick the sequence
$\{ f_{k,k} \}_{k=1}^\infty$.
This is a subsequence of the original sequence $\{ f_n \}_{n=1}^\infty$.
For every $m$, the tail $\{ f_{k,k} \}_{k=m}^\infty$ is a subsequence of
$\{ f_{m,k} \}_{k=1}^\infty$
and hence for any $m$ the sequence $\{ f_{k,k}(x_m) \}_{k=1}^\infty$ converges.
\end{proof}

For larger than countable sets,
we need the functions of the sequence to be related.  When we look at
continuous functions, the concept we need is equicontinuity.

\begin{defn}
Let $(X,d)$ be a metric space.
A set $S$ of functions
$f \colon X \to \C$ is 
\emph{\myindex{uniformly equicontinuous}}
if for every $\epsilon > 0$, there is a $\delta > 0$
such that if $x,y \in X$ with $d(x,y) < \delta$, we have
\begin{equation*}
\babs{f(x)-f(y)} < \epsilon \qquad \text{for all } f \in S .
\end{equation*}
\end{defn}

Notice that functions in a uniformly equicontinuous sequence are
all uniformly continuous.  It is 
not hard to show that a finite set of uniformly continuous functions
is uniformly equicontinuous.  The definition is really interesting
if $S$ is infinite.

Just as for continuity, one can define equicontinuity at a point.
That is, $S$ is \emph{\myindex{equicontinuous}} at $x \in X$
if for every $\epsilon > 0$, there is a $\delta > 0$
such that for $y \in X$ with $d(x,y) < \delta$, we have
$\babs{f(x)-f(y)} < \epsilon$ for all $f \in S$.
We will only deal with compact $X$ here, and
one can prove (exercise) that for a compact metric space $X$,
if $S$ is equicontinuous at every $x \in X$,
then it is uniformly equicontinuous.  For simplicity
we stick to uniform equicontinuity.

\begin{prop}
Suppose $(X,d)$ is a compact metric space,
$f_n \in C(X,\C)$, and $\{ f_n \}_{n=1}^\infty$
converges uniformly.  Then $\{ f_n \}_{n=1}^\infty$ is uniformly equicontinuous.
\end{prop}

\begin{proof}
Let $\epsilon > 0$ be given.
As $\{ f_n \}_{n=1}^\infty$ converges uniformly, there is an $N \in \N$ such that for
all $n \geq N$
\begin{equation*}
\babs{f_n(x)-f_N(x)} < \nicefrac{\epsilon}{3} \qquad \text{for all } x \in X.
\end{equation*}
As $X$ is compact, every continuous function is uniformly continuous.
So $\{ f_1,f_2,\ldots,f_N \}$ is a finite set of uniformly continuous
functions.  And so, as we mentioned above, the set is uniformly equicontinuous.
Hence there is a $\delta > 0$ such that
\begin{equation*}
\babs{f_j(x)-f_j(y)} < \nicefrac{\epsilon}{3} < \epsilon
\end{equation*}
whenever $d(x,y) < \delta$ and $1 \leq j \leq N$.

Take $n > N$.  For $d(x,y) < \delta$, we have
\begin{equation*}
%mbxlatex \begin{aligned}
\babs{f_n(x)-f_n(y)}
%mbxlatex &
\leq
\babs{f_n(x)-f_N(x)}
+
\babs{f_N(x)-f_N(y)}
+
\babs{f_N(y)-f_n(y)}
%mbxlatex \\
%mbxlatex &
<
\nicefrac{\epsilon}{3}
+
\nicefrac{\epsilon}{3}
+
\nicefrac{\epsilon}{3}
=\epsilon . \qedhere
%mbxlatex \end{aligned}
\end{equation*}
\end{proof}

\begin{prop}
A compact metric space $(X,d)$ contains a countable dense subset,
that is, there exists a countable $D \subset X$ such that $\widebar{D} = X$.
\end{prop}

\begin{proof}
For each $n \in \N$ there are finitely many
balls of radius $\nicefrac{1}{n}$ that cover $X$ (as $X$ is compact). That is,
for every $n$, there exists
a finite set of points $x_{n,1},x_{n,2},\ldots,x_{n,k_n}$ such that
\begin{equation*}
X= \bigcup_{j=1}^{k_n} B(x_{n,j},\nicefrac{1}{n}) .
\end{equation*}
Let $D \coloneqq \bigcup_{n=1}^\infty \{ x_{n,1},x_{n,2},\ldots,x_{n,k_n} \}$.
The set $D$ is countable as it is a countable union of finite sets.
For every $x \in X$
and every $\epsilon > 0$, there exists an $n$ such that
$\nicefrac{1}{n} < \epsilon$ and an $x_{n,j} \in D$ such that
\begin{equation*}
x \in B(x_{n,j},\nicefrac{1}{n}) \subset B(x_{n,j},\epsilon) .
\end{equation*}
Hence $x \in \widebar{D}$, so $\widebar{D} = X$, and $D$ is dense.
\end{proof}

We are now ready for the main result of this section,
the Arzel\`a--Ascoli theorem\footnote{%
Named after the Italian mathematicians
\href{https://en.wikipedia.org/wiki/Cesare_Arzel\%C3\%A0}{Cesare Arzel\`a}
(1847--1912), and
\href{https://en.wikipedia.org/wiki/Giulio_Ascoli}{Giulio Ascoli}
(1843--1896).} about existence of convergent subsequences.

\begin{thm}[Arzel\`a--Ascoli]\index{Arzel\`a--Ascoli theorem}
\label{thm:arzelaascoli}
Let $(X,d)$ be a compact metric space, and let $\{ f_n \}_{n=1}^\infty$
be pointwise bounded and uniformly equicontinuous sequence
of functions $f_n \in C(X,\C)$.  Then
$\{f_n\}_{n=1}^\infty$ is uniformly bounded and $\{ f_n \}_{n=1}^\infty$ contains a uniformly
convergent subsequence.
\end{thm}

Basically, a uniformly equicontinuous sequence in the metric space
$C(X,\C)$ that is pointwise bounded
is bounded (in $C(X,\C)$) and furthermore contains a convergent
subsequence in $C(X,\C)$.

As we mentioned before, because $X$ is compact, it is enough
to assume that $\{ f_n \}_{n=1}^\infty$ is equicontinuous;
uniform equicontinuity is automatic via an exercise.

\begin{proof}
We first show that the sequence is uniformly bounded.
By uniform equicontinuity,
there is a $\delta > 0$
such that
for all $x \in X$ and all $n \in \N$,
\begin{equation*}
B(x,\delta) \subset f_n^{-1}\bigl(B(f_n(x),1)\bigr) .
\end{equation*}
The space $X$ is compact, so there exist $x_1,x_2,\ldots,x_k$
such that
\begin{equation*}
X = \bigcup_{j=1}^k B(x_j,\delta) .
\end{equation*}
As $\{ f_n \}_{n=1}^\infty$ is pointwise bounded there exist $M_1,M_2,\ldots,M_k$
such that for $j=1,2,\ldots,k$,
\begin{equation*}
\babs{f_n(x_j)} \leq M_j \qquad \text{for all } n.
\end{equation*}
Let $M \coloneqq 1+ \max \{ M_1,M_2,\ldots,M_k \}$.  Given any
$x \in X$, there is a $j$ such that $x \in B(x_j,\delta)$.  Therefore,
for all $n$, we have
$x \in f_n^{-1}\bigl(B(f_n(x_j),1)\bigr)$, or in other words
\begin{equation*}
\babs{f_n(x)-f_n(x_j)} < 1 .
\end{equation*}
By the reverse triangle inequality,
\begin{equation*}
\babs{f_n(x)} < 1+ \babs{f_n(x_j)} \leq 1+M_j \leq M .
\end{equation*}
As $x$ was arbitrary, $\{f_n\}_{n=1}^\infty$ is uniformly bounded.


Next, pick a countable dense subset $D \subset X$.
By \propref{prop:subsequenceoncountableX}, we find
a subsequence $\{ f_{n_j} \}_{j=1}^\infty$ that converges pointwise on $D$.
Write $g_j \coloneqq f_{n_j}$ for simplicity.
The sequence $\{ g_n \}_{n=1}^\infty$ is 
uniformly equicontinuous.
Let $\epsilon > 0$ be given, then there exists a $\delta > 0$
such that for all $x \in X$ and all $n \in \N$,
\begin{equation*}
B(x,\delta) \subset g_n^{-1}\bigl(B(g_n(x),\nicefrac{\epsilon}{3})\bigr).
\end{equation*}
By density of $D$ and because $\delta$ is fixed, every $x \in X$ is in $B(y,\delta)$
for some $y \in D$.  By compactness of $X$,
there is a finite subset $\{ x_1,x_2,\ldots,x_k \} \subset D$
such that
\begin{equation*}
X = \bigcup_{j=1}^k B(x_j,\delta) .
\end{equation*}
As $\{ x_1,x_2,\ldots,x_k \}$
is a finite set and $\{ g_n \}_{n=1}^\infty$
converges pointwise on $D$, there exists a single $N$ such that for 
all $n,m \geq N$,
\begin{equation*}
\babs{g_n(x_j)-g_m(x_j)} < \nicefrac{\epsilon}{3}
 \qquad \text{for all } j=1,2,\ldots,k.
\end{equation*}

Let $x \in X$ be arbitrary.  There is some $j$ such that
$x \in B(x_j,\delta)$ and so for all $\ell \in \N$,
\begin{equation*}
\babs{g_\ell(x)-g_\ell(x_j)} < \nicefrac{\epsilon}{3}.
\end{equation*}
So for $n,m \geq N$,
\begin{equation*}
\begin{split}
\babs{g_n(x)-g_m(x)} & \leq
\babs{g_n(x)-g_n(x_j)} +
\babs{g_n(x_j)-g_m(x_j)} +
\babs{g_m(x_j)-g_m(x)}
\\
& <
\nicefrac{\epsilon}{3} +
\nicefrac{\epsilon}{3} +
\nicefrac{\epsilon}{3} = \epsilon .
\end{split}
\end{equation*}
Hence, $\{ g_n \}_{n=1}^\infty$ is uniformly Cauchy.  By completeness of $\C$,
it is uniformly convergent.
%FIXME: reference?
\end{proof}

\begin{cor}
Let $(X,d)$ be a compact metric space.
Let $S \subset C(X,\C)$ be a closed, bounded and uniformly equicontinuous set.
Then $S$ is compact.
\end{cor}

The theorem says that $S$
is sequentially compact and that means
compact in a metric space.
Recall that the closed unit ball in $C\bigl([0,1],\R\bigr)$, and therefore also in
$C\bigl([0,1],\C\bigr)$, is not compact.
Hence it cannot be a uniformly equicontinuous set.

\begin{cor}
Suppose $\{ f_n \}_{n=1}^\infty$ is a sequence of differentiable functions on $[a,b]$,
$\{ f_n' \}_{n=1}^\infty$ is uniformly bounded, and there is an
$x_0 \in [a,b]$ such that $\{ f_n(x_0) \}_{n=1}^\infty$ is bounded.
Then there exists a uniformly convergent
subsequence $\{ f_{n_j} \}_{j=1}^\infty$.
\end{cor}

\begin{proof}
The trick is to use the mean value theorem.  If $M$ is the uniform bound on
$\{ f_n' \}_{n=1}^\infty$, then by the mean value theorem for every $n$
\begin{equation*}
\babs{f_n(x)-f_n(y)} \leq M \sabs{x-y} \qquad \text{for all } x,y \in X.
\end{equation*}
All the $f_n$ are Lipschitz with the same constant and hence
the sequence is
uniformly equicontinuous.

Suppose $\babs{f_n(x_0)} \leq M_0$ for all $n$.
For all $x \in [a,b]$,
\begin{equation*}
\babs{f_n(x)} \leq \babs{f_n(x_0)}+ \babs{f_n(x)-f_n(x_0)} \leq M_0+ M \sabs{x-x_0}
\leq M_0 + M(b-a) .
\end{equation*}
So $\{ f_n \}_{n=1}^\infty$ is uniformly bounded.
We apply \hyperref[thm:arzelaascoli]{Arzel\`a--Ascoli} to find the subsequence.
\end{proof}

A classic application of the corollary above to Arzel\`a--Ascoli
in the theory of differential
equations is to prove the Peano existence
theorem, that is, the existence of solutions to ordinary differential
equations.  See \exerciseref{exercise:peanoexistence} below.

\medskip

Another application of Arzel\`a--Ascoli using the same idea as the
corollary above is the following.
Take a continuous $k \colon [0,1] \times [0,1] \to \C$.
For every $f \in C\bigl([0,1],\C\bigr)$ define
\begin{equation*}
T\bigl(f\bigr)(x) \coloneqq  \int_0^1 f(t) \, k(x,t)\,dt .
\end{equation*}
In exercises for earlier sections you have shown that 
$T$ is a linear operator on $C\bigl([0,1],\C\bigr)$.
Via Arzel\`a--Ascoli, we also find (exercise) that
the image of the unit ball of functions
\begin{equation*}
T\bigl( B(0,1) \bigr) = 
\bigl\{
Tf \in C\bigl([0,1],\C\bigr) :  
\snorm{f}_{[0,1]} < 1
\bigr\}
\end{equation*}
has compact closure, usually called
\emph{\myindex{relatively compact}}.
Such an operator is called a \emph{\myindex{compact operator}}.
And they are very useful.  Generally operators defined by
integration tend to be compact.

\subsection{Exercises}

\begin{exercise}
Let $f_n \colon [-1,1] \to \R$ be given by $f_n(x) \coloneqq \frac{nx}{1+{(nx)}^2}$.
Prove that the sequence is uniformly bounded, converges pointwise to 0,
yet there is no
subsequence that converges uniformly.
Which hypothesis of Arzel\`a--Ascoli
is not satisfied?  Prove your assertion.
\end{exercise}

\begin{exercise}
Define $f_n \colon \R \to \R$ by $f_n(x) \coloneqq \frac{1}{{(x-n)}^2+1}$.  Prove that
this sequence is uniformly bounded, uniformly equicontinuous, the sequence
converges pointwise to zero, yet there is no
subsequence that converges uniformly.
Which hypothesis of Arzel\`a--Ascoli
is not satisfied?  Prove your assertion.
\end{exercise}

\begin{exercise}
Let $(X,d)$ be a compact metric space, $C > 0$, $0 < \alpha \leq 1$, and
suppose $f_n \colon X \to \C$ are functions such that
$\babs{f_n(x)-f_n(y)} \leq C {d(x,y)}^\alpha$ for all $x,y \in X$ and
$n \in \N$.  Suppose also that there is a point $p \in X$ such that
$f_n(p) = 0$ for all $n$.
Show that there exists a uniformly convergent subsequence converging to
an $f \colon X \to \C$ that also satisfies $f(p) = 0$ and
$\babs{f(x)-f(y)} \leq C {d(x,y)}^\alpha$.
\end{exercise}

\begin{exercise}
Let $T \colon C\bigl([0,1],\C\bigr) \to C\bigl([0,1],\C\bigr)$ be the operator
given by
\begin{equation*}
T\bigl(f\bigr) (x) \coloneqq \int_0^x f(t)\, dt .
\end{equation*}
(That $T$ is linear and that $Tf$ is continuous follows from
linearity of the integral and
the fundamental theorem of calculus.)
\begin{enumerate}[a)]
\item
Show that $T$ takes the unit ball centered at 0 in $C\bigl([0,1],\C\bigr)$ into a
relatively compact set (a set with compact closure).
That is, $T$ is a compact operator.\\
Hint: See \volIref{\exerciseref*{vI-exercise:relativelycompactseq} in volume I}{\exerciseref{exercise:relativelycompactseq}}.
\item
Let $C \subset C\bigl([0,1],\C\bigr)$ the closed unit ball,
prove that
the image $T(C)$ is not closed (though it is relatively compact).
\end{enumerate}
\end{exercise}

\begin{samepage}
\begin{exercise}
Given $k \in C\bigl([0,1]\times [0,1],\C\bigr)$,
define the operator
$T \colon C\bigl([0,1],\C\bigr) \to C\bigl([0,1],\C\bigr)$ by
\begin{equation*}
T\bigl(f\bigr)(x) \coloneqq  \int_0^1 f(t) \, k(x,t)\,dt .
\end{equation*}
Show that $T$ takes the unit ball centered at 0 in $C\bigl([0,1],\C\bigr)$ into a
relatively compact set (a set with compact closure).
That is, $T$ is a compact operator.\\
Hint: See \volIref{\exerciseref*{vI-exercise:relativelycompactseq} in volume I}{\exerciseref{exercise:relativelycompactseq}}.
\\
Note: That $T$ is a well-defined linear operator was proved in
\exerciseref{exercise:continuouskernel}.
\end{exercise}
\end{samepage}

\begin{exercise}
Suppose $S^1 \subset \C$ is the unit circle, that is, the set where
$\sabs{z}=1$.  Suppose the continuous functions
$f_n \colon S^1 \to \C$ are uniformly bounded.
Let $\gamma \colon [0,1] \to S^1$ be a parametrization of $S^1$,
and $g(z,w)$ a continuous function on $C(0,1) \times S^1$
(here $C(0,1) \subset \C$ is the closed unit ball).  Define
the functions $F_n \colon C(0,1) \to \C$ by
the path integral (see \sectionref{sec:pathintegral})
\begin{equation*}
F_n(z) \coloneqq \int_\gamma f_n(w)\, g(z,w) \, ds(w) . 
\end{equation*}
Show that $\{ F_n \}_{n=1}^\infty$ has a uniformly convergent subsequence.
\end{exercise}

\begin{exercise}
Suppose $(X,d)$ is a compact metric space, $\{ f_n \}_{n=1}^\infty$ a uniformly equicontinuous
sequence of functions in $C(X,\C)$.  Suppose $\{ f_n \}_{n=1}^\infty$ converges
pointwise.  Show that it converges uniformly.
\end{exercise}

\begin{exercise}
Suppose that $\{ f_n \}_{n=1}^\infty$ is a uniformly equicontinuous uniformly bounded sequence of
$2\pi$-periodic functions $f_n \colon \R \to \R$.  Show that there is a
uniformly convergent subsequence.
\end{exercise}

\begin{exercise}
Show that for a compact metric space $X$,
a sequence $\{ f_n \}_{n=1}^\infty$ that is equicontinuous at every $x \in X$
is uniformly equicontinuous.
\end{exercise}

\begin{exercise}
Define $f_n \colon [0,1] \to \C$ by $f_n(t) \coloneqq e^{i(2\pi t + n)}$,
which gives a uniformly equicontinuous uniformly bounded sequence.
Prove a stronger conclusion than that of Arzel\`a--Ascoli for this sequence.
Let $\gamma \in \R$ be given,
and define $g(t) \coloneqq e^{i(2\pi t + \gamma)}$.  Show that there exists 
a subsequence of $\{ f_n \}_{n=1}^\infty$ converging uniformly to~$g$.
\\
Hint: Feel free to use the \emph{\myindex{Kronecker density theorem}}\footnote{%
Named after the German mathematician
\href{https://en.wikipedia.org/wiki/Leopold_Kronecker}{Leopold Kronecker}
(1823--1891).}:
The sequence $\{ e^{in} \}_{n=1}^\infty$ is dense in the unit circle.
\end{exercise}

\begin{exercise} \label{exercise:peanoexistence}
Prove the \emph{\myindex{Peano existence theorem}} (note the weaker
hypotheses than Picard, but also the lack of
uniqueness in this theorem):

\textbf{Theorem:} \emph{Suppose $F \colon I \times J \to \R$ is 
a continuous function where
$I, J \subset \R$ are closed bounded intervals, 
let $I^\circ$ and $J^\circ$ be their interiors,
and
let $(x_0,y_0) \in I^\circ \times J^\circ$.
Then there exists an $h > 0$ and a differentiable
function $f \colon [x_0 - h, x_0 + h] \to J \subset \R$, such that}
\begin{equation*}
f'(x) = F\bigl(x,f(x)\bigr) \qquad \text{and} \qquad f(x_0) = y_0.
\end{equation*}

We can prove this theorem by applying Euler's method
for numerical solutions to ODE and showing that there is
a subsequence of these approximations that converges.
\begin{enumerate}[a)]
\item
Show
that there exists an $h > 0$ so that the following
functions $f_n$ and $s_n$ are well-defined on $[x_0,x_0+h]$
for all $n \in \N$.
Define $f_n(x_0) \coloneqq y_0$ and $s_n(x_0) \coloneqq F(x_0,y_0)$
and for each $\ell=0,1,\ldots,n-1$ in turn
for $x \in \bigl(x_0+\frac{h\ell}{n},x_0+\frac{h(\ell+1)}{n}\bigr]$
let $s_n(x) \coloneqq F\bigl(x_0+\tfrac{h\ell}{n},
f_n(x_0+\tfrac{h\ell}{n})\bigr)$ and 
\begin{equation*}
\begin{split}
f_n(x) & \coloneqq f_n\bigl(x_0+\tfrac{h\ell}{n}\bigr)
+ \bigl(x-x_0-\tfrac{h\ell}{n}\bigr) F\bigl(x_0+\tfrac{h\ell}{n},
f_n(x_0+\tfrac{h\ell}{n}) \bigr) \\
& = f_n\bigl(x_0+\tfrac{h\ell}{n}\bigr)
+ \int_{x_0+\tfrac{h\ell}{n}}^x s_n(t) \, dt \\
& = y_0
+ \int_{x_0}^x s_n(t) \, dt .
\end{split}
\end{equation*}
\item
Show that $\{ f_n \}_{n=1}^\infty$ is equicontinuous
(in fact, Lipschitz with the same constant) and uniformly bounded.
Arzel\`a--Ascoli then says that there exists a
uniformly convergent subsequence $\{ f_{n_k} \}_{k=1}^\infty$
converging to a continuous $f$ defined on $[x_0,x_0+h]$.
\item
Prove that for all $x \in [x_0,x_0+h]$,
\begin{equation*}
\int_{x_0}^x s_{n_k}(t) \, dt
\quad
\text{converges to}
\quad
\int_{x_0}^x F\bigl(t,f(t)\bigr) \, dt ,
\end{equation*}
and therefore, by applying the fundamental theorem of calculus,
$f$ is differentiable on $[x_0,x_0+h]$ and $f'(x) = F\bigl(x,f(x)\bigr)$.
Hint 1: $F$ is uniformly continuous.
Hint 2: Show that $s_{n_k}$ converges uniformly.
\item
Use the solution above to show that $f$ can be defined in $[x_0-h,x_0]$ as
well,
possibly for a smaller $h > 0$.
\end{enumerate}
\end{exercise}

%%%%%%%%%%%%%%%%%%%%%%%%%%%%%%%%%%%%%%%%%%%%%%%%%%%%%%%%%%%%%%%%%%%%%%%%%%%%%%

\sectionnewpage
\section{The Stone--Weierstrass theorem}
\label{sec:stoneweier}

%mbxINTROSUBSECTION

\sectionnotes{3 lectures}

\subsection{Weierstrass approximation}

Perhaps surprisingly, even a very badly behaved continuous function is
a uniform limit of polynomials.
We cannot really get any \myquote{nicer}
functions than polynomials.
The idea of the proof is a very common approximation or \myquote{smoothing} idea
(convolution with an approximate delta function) that has applications
far beyond pure mathematics.

\begin{thm}[Weierstrass approximation theorem]
\index{Weierstrass approximation theorem}
If $f \colon [a,b] \to \C$ is continuous, then there exists a sequence
$\{ p_n \}_{n=1}^\infty$ of polynomials converging to $f$ uniformly on $[a,b]$.
Furthermore, if $f$ is real-valued, we can find $p_n$ with real coefficients.
\end{thm}

\begin{proof}
For $x \in [0,1]$,
define
\begin{equation*}
g(x) \coloneqq f\bigl((b-a)x+a\bigr)-f(a) - x\bigl(f(b)-f(a)\bigr) .
\end{equation*}
If we prove the theorem for $g$ and find the sequence
$\{ p_n \}_{n=1}^\infty$ for $g$,
it is proved for $f$ as we simply
composed with an invertible affine function and added an affine
function to $f$.  We reverse the process and apply that to our
$p_n$, to obtain polynomials approximating $f$.
The function $g$ is defined on $[0,1]$ and $g(0)=g(1)=0$.
For simplicity, assume that
$g$ is defined on $\R$ by letting
$g(x) \coloneqq 0$ if $x < 0$ or $x > 1$.  This extended $g$ is continuous.

Define
\begin{equation*}
c_n \coloneqq {\left( \int_{-1}^1 {(1-x^2)}^n\,dx \right)}^{-1} ,
\qquad
q_n(x) \coloneqq c_n (1-x^2)^n .
\end{equation*}
The choice of $c_n$ is
so that $\int_{-1}^1 q_n(x)\,dx = 1$.
See \figureref{fig:weierqn}.

\begin{myfigureht}
\myincludegraphics{weierqn}{%
Graphs of many functions on the interval minus 1 to 1.
They are all a symmetric "bump" above the x-axis.  The lowest
and widest bump is in dark color, and then taller and thinner peaks
are in increasingly lighter shade of gray.}
\caption{Plot of the approximate delta functions $q_n$ on $[-1,1]$ for
$n=5,10,15,20,\ldots,100$ with higher $n$ in lighter shade.\label{fig:weierqn}}
\end{myfigureht}

The functions $q_n$ are peaks around 0 (ignoring what happens outside
of $[-1,1]$) that get narrower and taller as $n$ increases,
while the area underneath is always 1.
A classic approximation idea
is to do a \emph{\myindex{convolution}} integral with peaks like this:
For $x \in [0,1]$, let
\begin{equation*}
p_n(x) \coloneqq \int_{0}^1 g(t)q_n(x-t) \,dt \quad \left( = \int_{-\infty}^\infty
g(t)q_n(x-t) \,dt \right) .
\end{equation*}
The idea of this convolution is that we do a \myquote{weighted average} of the
function $g$ around the point $x$ using $q_n$ as the weight.
See \figureref{fig:approxdeltaconv}.

\begin{myfigureht}
\myincludegraphics{approxdeltaconv}{%
Graphs of three functions on the interval minus 1 to 1.
There is a point on the x-axis marked simply x.  Symmetric about the
point x, there is a light gray positive function whose graph is a peak
centered at x and decaying rapidly as we get further away from x.
In gray, there is a graph of a function that is simply a jagged line.
In black is the product of the two graphs which is mostly almost zero where
the light gray line is zero, and near the peak the shape of the jagged
gray line is magnified.}
\caption{For $x=0.3$, the plot of $q_{100}(x-t)$ (light gray peak centered
at $x$), some continuous function
$g(t)$ (the jagged line) and the product $g(t)q_{100}(x-t)$ (the bold line).\label{fig:approxdeltaconv}}
\end{myfigureht}

As $q_n$ is a narrow peak, the integral
mostly sees the values of $g$ that are
close to $x$ and it does the weighted average of them.
When the peak gets narrower, we compute this average closer to $x$
and we expect the result to get closer to the value of $g(x)$.  Really, we are
approximating what is called a delta function\footnote{The delta function
is not actually a function,
it is a \myquote{thing} that should give
\myquote{$\int_{-\infty}^\infty g(t) \delta(x-t) \, dt = g(x)$.}}
(don't worry if you have not
heard of this concept),
and functions like $q_n$ are often called
\emph{approximate delta functions}\index{approximate delta function}.
We could do this with any set of polynomials that look like narrower
and narrower peaks near zero.  These just happen to be the simplest ones.
We only need this behavior on $[-1,1]$ as the convolution sees nothing
further than this as $g$ is zero outside $[0,1]$.

Because $q_n$ is a polynomial, we write
\begin{equation*}
q_n(x-t) = a_0(t) + a_1(t)\,x + \cdots + a_{2n}(t)\, x^{2n} ,
\end{equation*}
where $a_k(t)$ are polynomials in $t$, and hence
integrable functions.
So
\begin{equation*}
\begin{split}
p_n(x) & =
\int_{0}^1 g(t)q_n(x-t) \,dt
\\
&=
\left(
\int_0^1
g(t)
a_0(t)\,dt
\right)
+
\left(
\int_0^1
g(t)
a_1(t)\,dt
\right)
\,
x
+
\cdots
+
\left(
\int_0^1
g(t)
a_{2n}(t)\,dt
\right)
\,
x^{2n} .
\end{split}
\end{equation*}
In other words, $p_n$ is a polynomial%
\footnote{%
Do note that the functions $a_j$ depend on $n$, so the coefficients of $p_n$
change as $n$ changes.}
 in $x$.
If $g(t)$ is real-valued, then the functions $g(t)a_j(t)$ are
real-valued and $p_n$ has real coefficients,
proving the \myquote{furthermore} part of the theorem.

We still need to prove that $\{ p_n \}_{n=1}^\infty$ converges to $g$.
We start with estimating the size of~$c_n$.
For $x \in [0,1]$, we have that $1-x \leq 1-x^2$.  We estimate
\begin{equation*}
\begin{split}
c_n^{-1}   = \int_{-1}^1 {(1-x^2)}^n \, dx
& = 2\int_0^1 {(1-x^2)}^n \, dx \\
& \geq 2\int_0^{1} {(1-x)}^n \, dx
= \frac{2}{n+1} .
\end{split}
\end{equation*}
So $c_n \leq \frac{n+1}{2} \leq n$.
Let us see how small $q_n$ is if we ignore some small interval around the origin,
where the peak is.
Given any $\delta > 0$, $\delta < 1$,
we have that for all
$x$ such that $\delta \leq \sabs{x} \leq 1$,
\begin{equation*}
q_n(x) \leq c_n {(1-\delta^2)}^n \leq  n{(1-\delta^2)}^n ,
\end{equation*}
because $q_n$ is increasing on $[-1,0]$ and decreasing on $[0,1]$.
By the ratio test, 
$n{(1-\delta^2)}^n$ goes to 0 as $n$ goes to infinity.

The function $q_n$ is even, $q_n(t) = q_n(-t)$, and $g$
is zero outside of $[0,1]$.
So for $x \in [0,1]$,
\begin{equation*}
p_n(x) = 
\int_{0}^1 g(t)q_n(x-t) \, dt
=
\int_{-x}^{1-x} g(x+t)q_n(-t) \, dt
=
\int_{-1}^{1} g(x+t)q_n(t) \, dt .
\end{equation*}
Let $\epsilon > 0$ be given.
As $[-1,2]$ is compact and $g$ is continuous on $[-1,2]$, we have that $g$ is uniformly continuous.
Pick $0 < \delta < 1$ such that if
$\sabs{x-y} < \delta$ (and $x,y \in [-1,2]$), then
\begin{equation*}
\babs{g(x)-g(y)} < \frac{\epsilon}{2} .
\end{equation*}
Let $M$ be such that $\babs{g(x)} \leq M$ for all $x$.  Let $N$ be
such that for all $n \geq N$,
\begin{equation*}
4M n{(1-\delta^2)}^n < \frac{\epsilon}{2} .
\end{equation*}
Note that 
$\int_{-1}^1 q_n(t) \, dt = 1$ and $q_n(t) \geq 0$ on $[-1,1]$.  So for $n
\geq N$ and every $x \in [0,1]$,
\begin{align*}
\babs{p_n(x)-g(x)} & =
\abs{\int_{-1}^1 g(x+t)q_n(t) \, dt
-g(x)\int_{-1}^1 q_n(t) \, dt} \\
& =
\abs{\int_{-1}^1 \bigl(g(x+t)-g(x)\bigr)q_n(t) \, dt}
\displaybreak[0]\\
& \leq
\int_{-1}^1 \babs{g(x+t)-g(x)} q_n(t) \, dt
\displaybreak[0]\\
& =
\int_{-1}^{-\delta} \babs{g(x+t)-g(x)} q_n(t) \, dt
\quad +
\int_{-\delta}^{\delta} \babs{g(x+t)-g(x)} q_n(t) \, dt
\\
& \phantom{\leq} +
\int_{\delta}^1 \babs{g(x+t)-g(x)} q_n(t) \, dt
\displaybreak[0]\\
& \leq
2M
\int_{-1}^{-\delta} q_n(t) \, dt
\quad
+
\quad
\frac{\epsilon}{2}
\int_{-\delta}^{\delta} q_n(t) \, dt
\quad
+
\quad
2M
\int_{\delta}^1 q_n(t) \, dt
\\
& \leq
2M n{(1-\delta^2)}^n(1-\delta)
\quad
+
\quad
\frac{\epsilon}{2}
\quad
+
\quad
2M n{(1-\delta^2)}^n(1-\delta) \\
& <
4M n{(1-\delta^2)}^n
+
\frac{\epsilon}{2}
< \epsilon . \qedhere
\end{align*}
\end{proof}

A convolution often inherits some property of the functions we are convolving.
In our case the convolution $p_n$ inherited the property of being a
polynomial from $q_n$.  The same idea of the proof is often used 
to get other properties.  If $q_n$ or $g$ is infinitely differentiable, so is $p_n$.
If $q_n$ or $g$ is a solution to a linear differential equation, so is $p_n$.
Etc.

Let us note an immediate application of the Weierstrass theorem.  We have
already seen that countable dense subsets can be very useful.

\begin{cor}
The metric spaces $C\bigl([a,b],\R\bigr)$ and $C\bigl([a,b],\C\bigr)$ each contain a countable dense subset.
\end{cor}

\begin{proof}
Without loss of generality, consider only
$C\bigl([a,b],\R\bigr)$
(why?).
Real polynomials are dense in $C\bigl([a,b],\R\bigr)$ by Weierstrass.
If we show that
every real polynomial can be approximated by polynomials with rational
coefficients, we are done.
Indeed, there are only countably many
rational numbers and so there are only countably many polynomials with
rational coefficients (a countable union of countable sets is countable).

Further without loss of generality, suppose $[a,b]=[0,1]$.  Let
\begin{equation*}
p(x) \coloneqq \sum_{k=0}^n a_k\,  x^k
\end{equation*}
be a polynomial of degree $n$ where $a_k \in \R$.  Given $\epsilon > 0$, pick $b_k \in \Q$
such that $\sabs{a_k-b_k} < \frac{\epsilon}{n+1}$.  Then
if we let
\begin{equation*}
q(x) \coloneqq \sum_{k=0}^n b_k \, x^k ,
\end{equation*}
we have
\begin{equation*}
\babs{p(x)-q(x)}
=
\abs{\sum_{k=0}^n (a_k-b_k) x^k}
\leq
\sum_{k=0}^n \sabs{a_k-b_k} x^k
\leq
\sum_{k=0}^n \sabs{a_k-b_k}
<
\sum_{k=0}^n \frac{\epsilon}{n+1} = \epsilon . \qedhere
\end{equation*}
\end{proof}

\begin{remark}
While we will not prove so, the corollary above implies that
$C\bigl([a,b],\C\bigr)$ has the same cardinality as $\R$, which may be a
bit surprising.  The set of all functions $[a,b] \to \C$ has
cardinality strictly greater than the cardinality of $\R$, it has the
cardinality of the power set of $\R$.  So the
set of continuous functions is a very tiny subset of the set of all
functions.
\end{remark}

\textbf{Warning!}
The fact that every continuous function $f \colon [-1,1] \to \C$ (or any
interval $[a,b]$) can be uniformly
approximated by polynomials
\begin{equation*}
\sum_{k=0}^n a_k\,  x^k
\end{equation*}
does not mean that every continuous $f$ is analytic, that is, equal to a
power series
\begin{equation*}
\sum_{k=0}^\infty c_k\,  x^k .
\end{equation*}
An analytic function is infinitely differentiable,
however,
the function $\sabs{x}$ is continuous and near the origin
approximable by polynomials,
and so provides a counterexample.

The key distinction is that
the polynomials coming from the Weierstrass theorem are not the partial
sums of a power series.  For each one, the coefficients $a_k$ above can be
completely different---they do not need to come from a single sequence
$\{ c_k \}_{k=1}^\infty$.

\medskip

Interestingly, to generalize Weierstrass, we will only need to use it to
approximate the absolute value function by polynomials without
a constant term.

\begin{cor}
Let $[-a,a]$ be an interval.  Then there is a sequence of real polynomials
$\{ p_n \}_{n=1}^\infty$ that converges uniformly to $\sabs{x}$ on $[-a,a]$ and such that
$p_n(0) = 0$ for all $n$.
\end{cor}

\begin{proof}
As $f(x) \coloneqq \sabs{x}$ is continuous and real-valued
on $[-a,a]$, the Weierstrass theorem gives a sequence of
real polynomials $\{ \widetilde{p}_n \}_{n=1}^\infty$ that converges to $f$
uniformly on $[-a,a]$.
Let
\begin{equation*}
p_n(x) \coloneqq \widetilde{p}_n(x) - \widetilde{p}_n(0) .
\end{equation*}
Obviously $p_n(0) = 0$.

Given $\epsilon > 0$, let $N$ be such that
for $n \geq N$, we have
$\bigl\lvert\widetilde{p}_n(x)-\sabs{x}\big\rvert <
\nicefrac{\epsilon}{2}$ for all $x \in [-a,a]$.
In particular,
$\babs{\widetilde{p}_n(0)} < \nicefrac{\epsilon}{2}$.
Then for $n \geq N$,
\begin{equation*}
\babs{ p_n(x)-\sabs{x} }
=
\babs{ \widetilde{p}_n(x) - \widetilde{p}_n(0) - \sabs{x} }
\leq
\babs{ \widetilde{p}_n(x) - \sabs{x} } + \babs{\widetilde{p}_n(0)} < 
\nicefrac{\epsilon}{2} + \nicefrac{\epsilon}{2} = \epsilon . \qedhere
\end{equation*}
\end{proof}

Generalizing the corollary,
we can make the polynomials from the Weierstrass theorem
be equal to our target function at one point, not just for $\sabs{x}$, but
that's the one we will need.  It is also possible (see
\exerciseref{exercise:finitelymanyweierequal})
to make the polynomials
equal at finitely many points by subtracting not a constant but a properly
crafted polynomial.

\subsection{Stone--Weierstrass approximation}

We want to abstract away what is not really
necessary and prove a general version of the Weierstrass theorem,
the Stone--Weierstrass theorem\footnote{%
Named after the American mathematician
\href{https://en.wikipedia.org/wiki/Marshall_Harvey_Stone}{Marshall Harvey Stone}
(1903--1989), and the German mathematician
\href{https://en.wikipedia.org/wiki/Karl_Weierstrass}{Karl Theodor Wilhelm Weierstrass}
(1815--1897).}.
Polynomials are dense in the space of continuous
functions on a compact interval.  What other kinds of families of
functions are also dense?  And if the domain is an
arbitrary metric space, then we no longer have polynomials
to begin with.

\begin{defn}
\pagebreak[2]
A set $\sA$ of complex-valued functions $f \colon X \to \C$ is said to be an 
\emph{\myindex{algebra}} (sometimes
\emph{\myindex{complex algebra}} or \emph{algebra over $\C$}) if for all $f,
g \in \sA$ and $c \in \C$, we have
\begin{enumerate}[(i)]
\item $f+g \in \sA$.
\item $fg \in \sA$.
\item $cg \in \sA$.
\end{enumerate}
A \emph{\myindex{real algebra}} or an
\emph{algebra over $\R$} is a set of real-valued
functions that satisfies the three properties above for $c \in \R$.
\end{defn}

We are interested in the case when
$X$ is a compact metric space.  Then
$C(X,\C)$ and $C(X,\R)$ are metric spaces.
Given a set $\sA \subset C(X,\C)$, the set of all uniform
limits is the metric space closure $\widebar{\sA}$.
When we talk about closure of an algebra
from now on we mean the closure in $C(X,\C)$
as a metric space.  Same for $C(X,\R)$.

The set $\sP$ of all polynomials is an algebra in
$C\bigl([a,b],\C\bigr)$, and we
have shown that its closure $\widebar{\sP} = C\bigl([a,b],\C\bigr)$.
That is, it is dense.  That is the sort of result that we wish to prove.

We leave the following proposition as an exercise.

\begin{prop} \label{prop:closureofalgebra}
Suppose $X$ is a compact metric space.
If $\sA \subset C(X,\C)$ is an algebra, then the closure $\widebar{\sA}$ is also an algebra.
Similarly for a real algebra in $C(X,\R)$.
\end{prop}

We distill the properties of polynomials that are sufficient
for an approximation theorem.

\begin{defn}
Let $\sA$ be a set of complex-valued functions defined on a set $X$.
\begin{enumerate}[(i)]
\item $\sA$ \emph{\myindex{separates points}}
if for every $x,y \in X$ with $x \neq y$, there is an $f \in \sA$ such that
$f(x) \neq f(y)$.
\item 
$\sA$ \emph{\myindex{vanishes at no point}} if for every $x \in X$
there is an $f \in \sA$ such that $f(x) \neq 0$.
\end{enumerate}
\end{defn}

\begin{example}
Given any $X \subset \R$ (or $X \subset \C$),
the set $\sP$ of polynomials in one variable separates points and vanishes at no point
on $X$.  That is, $1 \in \sP$, so it vanishes at no point.
For $x,y \in X$, $x \neq y$, take $f(t) \coloneqq t$, to get $f(x) = x \neq y = f(y)$.
So $\sP$ separates points.
\end{example}

\begin{example}
The set of functions of the form
\begin{equation*}
f(t) = a_0 + \sum_{n=1}^k a_n \cos(nt)
\end{equation*}
is an algebra,
which follows by the identity $\cos(mt)\cos(nt) = 
\frac{\cos((n+m) t)}{2}+
\frac{\cos((n-m) t)}{2}$.
The algebra vanishes at no point as it contains a constant function.
It does not separate points if the domain is an interval
$[-a,a]$, as $f(-t) = f(t)$ for all~$t$.
It does separate points if the domain is $[0,\pi]$; $\cos(t)$
is one-to-one on $[0,\pi]$.
\end{example}

\begin{example}
The set $\sP$ of real polynomials with no constant term is an algebra that vanishes at the origin.
Clearly, any function in the closure of $\sP$ also vanishes at the origin,
so the closure of $\sP$ cannot be $C\bigl([0,1],\R\bigr)$.

Similarly, the set of constant functions is an algebra that does not
separate points.  Uniform limits of constants are constants, so we also do
not obtain all continuous functions.
\end{example}

It is interesting that these two properties,
\myquote{vanishes at no point} and
\myquote{separates points,} are sufficient to obtain approximation of any
real-valued continuous function.
Before we prove this theorem, we note that such an algebra can interpolate
a finite number of values exactly.  We will state this result
only for two points as that is all that we will require.

\begin{prop} \label{prop:SWinterpolate}
Suppose $\sA$ is an algebra of complex-valued functions on a set $X$ that separates points
and vanishes at no point.  Suppose $x,y$ are distinct points of $X$, and
$c,d \in \C$.  Then there is an $f \in \sA$ such that
\begin{equation*}
f(x) = c, \qquad f(y) = d .
\avoidbreak
\end{equation*}
If $\sA$ is a real algebra, the conclusion holds for $c,d \in \R$.
\end{prop}

\begin{proof}
There must exist an $g,h,k \in \sA$
such that 
$g(x) \neq g(y)$, $h(x) \neq 0$, $k(y) \neq 0$.
Let
\begin{equation*}
\begin{split}
f & \coloneqq 
c
\frac{\bigl(g - g(y)\bigr)h}{\bigl(g(x)-g(y)\bigr)h(x) } + 
d
\frac{\bigl(g - g(x)\bigr)k}{\bigl(g(y)-g(x)\bigr)k(y)}
\\
& =
c
\frac{gh - g(y)h}{g(x)h(x)-g(y)h(x) } + 
d
\frac{gk - g(x)k}{g(y)k(y)-g(x)k(y)} .
\end{split}
\end{equation*}
We are not dividing by zero (clear from the first formula).
Also by the first formula, $f(x) = c$ and $f(y) = d$.
By the second formula, $f \in \sA$ (as $\sA$ is an algebra).
\end{proof}

\begin{thm}[Stone--Weierstrass, real version]
\label{thm:SWreal}%
\index{Stone--Weierstrass!real version}%
Let $X$ be a compact metric space and $\sA$ a real algebra of real-valued
continuous functions on $X$, such that $\sA$ separates points and vanishes at
no point.  Then the closure $\widebar{\sA} = C(X,\R)$.
\end{thm}

The proof is divided into several claims.

\medskip

\noindent
\textbf{Claim 1:} \emph{If $f \in \widebar{\sA}$, then $\sabs{f} \in
\widebar{\sA}$.}

\begin{proof}
The function $f$ is bounded (continuous on a compact set), so there is an $M$
such that $\babs{f(x)} \leq M$ for all $x \in X$.
Let $\epsilon > 0$ be given.  By the corollary to the Weierstrass theorem,
there exists a real polynomial $c_1 y + c_2 y^2 + \cdots+ c_n y^n$
(vanishing at $y=0$) such that
\begin{equation*}
\abs{\sabs{y} - \sum_{k=1}^n c_k y^k} < \epsilon
\end{equation*}
for all $y \in [-M,M]$.
Because $\widebar{\sA}$ is an algebra and because there is no constant term in the
polynomial,
\begin{equation*}
\sum_{k=1}^n c_k f^k \in \widebar{\sA} .
\end{equation*}
As $\babs{f(x)} \leq M$, we have that for all $x \in X$,
\begin{equation*}
\abs{\babs{f(x)} - \sum_{k=1}^n c_k {\bigl(f(x)\bigr)}^k}
< \epsilon .
\end{equation*}
So $\sabs{f}$ is in the closure of $\widebar{\sA}$, which is itself closed.
In other words, $\sabs{f} \in \widebar{\sA}$.
\end{proof}

\medskip

\noindent
\textbf{Claim 2:} \emph{If $f \in \widebar{\sA}$ and $g \in \widebar{\sA}$, then
$\max(f,g) \in \widebar{\sA}$ and
$\min(f,g) \in \widebar{\sA}$, where
}
\begin{equation*}
\bigl(\max(f,g)\bigr) (x) \coloneqq \max \bigl\{ f(x), g(x) \bigr\} , \qquad
\text{and} \qquad
\bigl(\min(f,g)\bigr) (x) \coloneqq \min \bigl\{ f(x), g(x) \bigr\} .
\end{equation*}

\begin{proof}
Write:
\begin{equation*}
\max(f,g) = \frac{f+g}{2} + \frac{\sabs{f-g}}{2} ,
\qquad\text{and}\qquad
\min(f,g) = \frac{f+g}{2} - \frac{\sabs{f-g}}{2} .
\end{equation*}
As $\widebar{\sA}$ is an algebra we are done.
\end{proof}

By induction, the claim is also true for the minimum or maximum of a finite
collection of functions.

\medskip

\noindent
\textbf{Claim 3:} \emph{Given $f \in C(X,\R)$, $x \in X$, and $\epsilon > 0$,
there exists a $g_x \in \widebar{\sA}$ with $g_x(x) = f(x)$ and
}
\begin{equation*}
g_x(t) > f(t)-\epsilon \qquad \text{for all } t \in X.
\end{equation*}

\begin{proof}
Fix $f$, $x$, and $\epsilon$.
By \propref{prop:SWinterpolate}, for every $y \in X$, find an $h_y \in
\sA$ such that
\begin{equation*}
h_y(x) = f(x), \qquad h_y(y)=f(y) .
\end{equation*}
As $h_y$ and $f$ are continuous, the function $h_y-f$ is continuous,
and the set
\begin{equation*}
U_y \coloneqq
\bigl\{ t \in X : h_y(t) > f(t) -\epsilon \bigr\}
=
{(h_y-f)}^{-1} \bigl( (-\epsilon,\infty) \bigr)
\end{equation*}
is open (it is the inverse image of an open set by a continuous function).
Furthermore $y \in U_y$.  So the sets $U_y$ cover $X$.
The space $X$ is compact,
so there exist finitely many
points $y_1,y_2,\ldots,y_n$ in $X$ such
that
\begin{equation*}
X = \bigcup_{k=1}^n U_{y_k}  .
\end{equation*}
Let 
\begin{equation*}
g_x \coloneqq \max(h_{y_1},h_{y_2},\ldots,h_{y_n}) .
\end{equation*}
By Claim 2, $g_x \in \widebar{\sA}$.  See \figureref{fig:stonegx}.
Moreover,
\begin{equation*}
g_x(t) > f(t) -\epsilon
\end{equation*}
for all $t \in X$, since for every $t$, there is a $y_k$ such that 
$t \in U_{y_k}$, and so 
$h_{y_k}(t) > f(t) -\epsilon$.
Finally, $h_y(x) = f(x)$ for all $y \in X$, so
$g_x(x) = f(x)$.
\end{proof}

\begin{myfigureht}
\myincludepdft{stonegx}{%
A graph of several functions in two dimensions.  On the horizontal
axis three points from left to right are marked, y sub 1, y sub 2,
and x.  A graph of a function f is given in dark bold line.
A graph of f minus epsilon is given underneath in a dotted line.
A graph of h sub y sub 1 is given in long dashes and starts above
f on the left, crosses the graph of f at y sub 1, then goes below
even f minus epsilon for a while in particular around y sub 2,
then goes up through the graph of f again at x.
A graph of h sub y sub 2 starts below f minus epsilon, then goes up,
crosses the graph of f at y sub 2, then crosses it
back down at x.  Notably at every point either the graph of h sub y sub 1
or h sub y sub 2 is above f minus epsilon, and the higher value of these two
is highlighted in a thick gray line and marked g sub x, again
in particular the graph of g sub x notably lies above f minus epsilon
and agrees with f at x.}
\caption{Construction of $g_x$ out of two $h_{y_1}$ (longer dashes) and
$h_{y_2}$ (shorter dashes).\label{fig:stonegx}}
\end{myfigureht}

What we have now is for each $x$ a function $g_x \in \widebar{\sA}$
that is within $\epsilon$ of $f$ near $x$ (being continuous),
but also $g_x$ is within $\epsilon$ of $f$ from at least one side at
all points.
If we cover $X$ with neighborhoods where $g_x$ is a good approximation,
we can repeat the idea of the argument with a minimum to get
a function that is within $\epsilon$ from both sides.

\medskip
\pagebreak[2]

\noindent
\textbf{Claim 4:} \emph{If $f \in C(X,\R)$ and $\epsilon > 0$ is given, then there
exists an $\varphi \in \widebar{\sA}$ such that}
\begin{equation*}
\babs{f(x) - \varphi(x)} < \epsilon .
\end{equation*}

\begin{proof}
For every $x \in X$, find the function $g_x$ as in Claim 3.
Let
\begin{equation*}
V_x \coloneqq \bigl\{ t \in X : g_x(t) < f(t) + \epsilon \bigr\}.
\end{equation*}
The sets $V_x$ are open as $g_x$ and $f$ are continuous.
As $g_x(x) = f(x)$, we have $x \in V_x$.  So the sets $V_x$ cover $X$.
By compactness of $X$,
there
are finitely many points $x_1,x_2,\ldots,x_n$ such that
\begin{equation*}
X = \bigcup_{k=1}^n V_{x_k} .
\end{equation*}
Let
\begin{equation*}
\varphi \coloneqq \min(g_{x_1},g_{x_2},\ldots,g_{x_n}) .
\end{equation*}
By Claim 2, $\varphi \in \widebar{\sA}$.  Similarly as before (same argument as in
Claim 3), for all $t \in X$,
\begin{equation*}
\varphi(t) < f(t) + \epsilon .
\end{equation*}
Since all the $g_x$ satisfy $g_x(t) > f(t) - \epsilon$ for all $t \in X$,
$\varphi(t) > f(t) - \epsilon$ as well.
Hence, for all $t$,
\begin{equation*}
-\epsilon < \varphi(t) - f(t) < \epsilon ,
\end{equation*}
which is the desired conclusion.
\end{proof}

The proof of the theorem follows from Claim 4.  The claim states that an
arbitrary continuous function is in the closure of $\widebar{\sA}$,
which is already closed.  The theorem is proved.

\begin{example}
The functions of the form
\begin{equation*}
f(t) = \sum_{k=1}^n c_k \, e^{kt},
\end{equation*}
for $c_k \in \R$,
are dense in $C\bigl([a,b],\R\bigr)$.  Such functions are a real
algebra, which follows from $e^{kt} e^{\ell t} = e^{(k+\ell)t}$.  They separate
points as $e^t$ is one-to-one.
As $e^t > 0$ for all $t$, the algebra
does not vanish at any point.
\end{example}

In general, given a set of functions that separates points and does
not vanish at any point, we let these functions
\emph{generate}\index{generate an algebra}
an algebra
by considering all the linear combinations of arbitrary multiples of such
functions.  That is, we consider all real polynomials without constant term
of such functions.  In the example above,
the algebra is generated by $e^t$.  We 
consider polynomials in $e^t$ without constant term.

\begin{example}
We mentioned that the set of all functions of the form
\begin{equation*}
a_0 +
\sum_{n=1}^N a_n \cos(nt)
\end{equation*}
is an algebra.
When considered on $[0,\pi]$, 
it separates points and vanishes nowhere so
\hyperref[thm:SWreal]{Stone--Weierstrass} applies.
As for polynomials, you \emph{do not} want to conclude that every continuous
function on $[0,\pi]$ has a uniformly convergent
Fourier cosine series, that is, that every continuous
function can be written as
\begin{equation*}
a_0 +
\sum_{n=1}^\infty a_n \cos(nt) .
\end{equation*}
That is \emph{not true}!
There exist continuous functions
whose Fourier series does not converge even pointwise
let alone uniformly.  See \sectionref{sec:fourier}.
\end{example}

To obtain Stone--Weierstrass for complex algebras, we must
make an extra assumption.

\begin{defn}
An algebra $\sA$ is \emph{\myindex{self-adjoint}} if for all $f \in \sA$, the function
$\bar{f}$ defined by $\bar{f}(x) \coloneqq \overline{f(x)}$ is in $\sA$, where by the
bar we mean the complex conjugate.
\end{defn}

\begin{thm}[Stone--Weierstrass, complex version]
\label{thm:SWcomplex}%
\index{Stone--Weierstrass!complex version}%
Let $X$ be a compact metric space and $\sA$ an algebra of complex-valued
continuous functions on $X$, such that $\sA$ separates points, vanishes at
no point, and is self-adjoint.  Then the closure $\widebar{\sA} = C(X,\C)$.
\end{thm}

\begin{proof}
Suppose $\sA_\R \subset \sA$ is the set of the real-valued elements of
$\sA$.
For $f \in \sA$, write $f = u+iv$ where $u$ and $v$ are real-valued.
Then
\begin{equation*}
u = \frac{f+\bar{f}}{2}, \qquad
v = \frac{f-\bar{f}}{2i} .
\end{equation*}
So $u, v \in \sA$ as $\sA$ is a self-adjoint algebra, and since they are
real-valued $u, v \in \sA_\R$.

If $x \neq y$, then find an $f \in \sA$ such that $f(x) \neq f(y)$.  If $f
= u+iv$, then it is obvious that either $u(x) \neq u(y)$ or
$v(x) \neq v(y)$.  So $\sA_\R$ separates points.
%
Similarly, for every $x$ find $f \in \sA$ such that $f(x) \neq 0$.  If $f
= u+iv$, then either $u(x) \neq 0$ or $v(x) \neq 0$.
So $\sA_\R$ vanishes at no point.
%
The set $\sA_\R$ is a real algebra, and satisfies the hypotheses of the
\hyperref[thm:SWreal]{real Stone--Weierstrass theorem}.
Given any $f = u+iv \in C(X,\C)$,
we find $g,h \in \sA_\R$ such that
$\babs{u(t)-g(t)} < \nicefrac{\epsilon}{2}$ and
$\babs{v(t)-h(t)} < \nicefrac{\epsilon}{2}$ for all $t \in X$.
Next, $g+i h \in \sA$, and
\begin{multline*}
\abs{f(t) - \bigl(g(t)+ih(t)\bigr)} = 
\abs{u(t)+iv(t) - \bigl(g(t)+ih(t)\bigr)} \\
\leq
\babs{u(t)-g(t)}+\babs{v(t)-h(t)} < \nicefrac{\epsilon}{2} +
\nicefrac{\epsilon}{2} = \epsilon
\end{multline*}
for all $t \in X$.
So $\widebar{\sA} = C(X,\C)$.
\end{proof}

The self-adjoint requirement is necessary, although it is not so obvious to
see it.  For an example, see \exerciseref{exercise:selfadjointSW}.

We give an interesting application.
When working
with functions of two variables, it may be useful to work with functions
of the form $f(x)g(y)$ rather than $F(x,y)$.  For example, they are easier
to integrate.  We have the following.

\begin{example}
Any continuous $F \colon [0,1] \times [0,1] \to \C$ can be
approximated uniformly by functions of the form
\begin{equation*}
\sum_{j=1}^n f_j(x) g_j(y) ,
\end{equation*}
where $f_j \colon [0,1] \to \C$ and $g_j \colon [0,1] \to \C$ are continuous.

Proof:
It is not hard to see that the functions of the above form are a complex
algebra.  It is equally easy to show that they vanish nowhere, separate
points, and the algebra is self-adjoint.  As $[0,1] \times [0,1]$ is
compact,
\hyperref[thm:SWcomplex]{Stone--Weierstrass} obtains the result.
\end{example}

\subsection{Exercises}

\begin{exercise}
Prove \propref{prop:closureofalgebra}.
Hint: If $\{ f_n \}_{n=1}^\infty$ is a sequence in $C(X,\R)$
converging to $f$, then as $f$ is bounded, show
that $f_n$ is uniformly bounded, that is, there exists a
single bound for all $f_n$ (and $f$).
\end{exercise}

\begin{exercise}
Suppose $X \coloneqq \R$ (not compact in particular).
Show that $f(t) \coloneqq e^t$ is not possible to uniformly approximate
by polynomials on $X$.  Hint: Consider $\babs{\frac{e^t}{t^n}}$
as $t \to \infty$.
\end{exercise}

\begin{exercise}
Suppose $f \colon [0,1] \to \C$ is a uniform limit of a sequence of polynomials
of degree at most $d$.  Show that the limit is a polynomial of degree at most $d$.
Conclude that to approximate a function which is not a polynomial, we
need the degree of the approximations to go to infinity.\\
Hint: First prove that if a sequence of polynomials of degree $d$
converges uniformly to the zero function, then 
the coefficients converge to zero.
One way to do this is linear algebra: Consider a polynomial $p$
evaluated at $d+1$ points
to be a linear operator taking the
coefficients of $p$ to the values of $p$
(an operator in $L(\R^{d+1})$).
\end{exercise}

\begin{exercise}
Suppose $f \colon [0,1] \to \R$ is continuous and
$\int_0^1 f(x) x^n \, dx = 0$ for all $n = 0,1,2,\ldots$.
Show that $f(x) = 0$ for all $x \in [0,1]$.
Hint: Approximate by polynomials to show that $\int_0^1 {\bigl( f(x)
\bigr)}^2 \, dx = 0$.
\end{exercise}

\begin{exercise}
Suppose $I \colon C\bigl([0,1],\R\bigr) \to \R$ is 
a linear continuous function such that
$I(x^n) = \frac{1}{n+1}$
for all $n=0,1,2,3,\ldots$.
Prove that $I(f) = \int_0^1 f$ for all $f \in C\bigl([0,1],\R\bigr)$.
\end{exercise}

\begin{exercise}
Let $\sA$ be the collection of real polynomials in $x^2$,
that is, polynomials of the form
$c_0 + c_1 x^2 + c_2 x^4 + \cdots + c_d x^{2d}$.
\begin{enumerate}[a)]
\item
Show that every $f \in C\bigl([0,1],\R\bigr)$ is a uniform limit of
polynomials from $\sA$.
\item
Find an $f \in  C\bigl([-1,1],\R\bigr)$ that is not a
uniform limit of
polynomials from $\sA$.
\item
Which hypothesis of the real Stone--Weierstrass is not satisfied
for the domain $[-1,1]$?
\end{enumerate}
\end{exercise}

\begin{exercise}
\pagebreak[3]
Let $\sabs{z}=1$ define the unit circle $S^1 \subset \C$.
\begin{enumerate}[a)]
\item
Show that functions of the form
\begin{equation*}
\sum_{k=-n}^n c_k\, z^k
\end{equation*}
are dense in $C(S^1,\C)$.  Notice the negative powers.
\item 
Show that functions of the form
\begin{equation*}
c_0
+
\sum_{k=1}^n c_k \, z^k
+
\sum_{k=1}^n c_{-k}\, \bar{z}^k
\end{equation*}
are dense in $C(S^1,\C)$.
These functions are called \emph{\myindex{harmonic polynomials}},
and this approximation leads to, for example, the solution of the
steady state heat problem.
\end{enumerate}
Hint: A good way to write the equation for $S^1$ is $z \bar{z} = 1$.
\end{exercise}

\begin{exercise} \label{exercise:trigpolydense}
Show that any continuous function $f \colon [-\pi,\pi] \to \C$
with $f(-\pi)=f(\pi)$ can be approximated uniformly by
functions of the form
\begin{equation*}
\sum_{k=-n}^n c_k \, e^{ik x}
\end{equation*}
for complex numbers $c_k$.
\end{exercise}

\begin{exercise} \label{exercise:selfadjointSW}
Let $S^1 \subset \C$ be the unit circle, that is, the set where
$\sabs{z} = 1$.  Orient this set counterclockwise.
Let $\gamma(t) \coloneqq e^{it}$.
For the one-form $f(z)\,dz$ we write\footnote{%
Alternatively, one could define $dz \coloneqq dx + i \, dy$ and
extend the path integral from \chapterref{path:chapter} to complex-valued
one-forms.}
\begin{equation*}
\int_{S^1} f(z) \,dz \coloneqq \int_0^{2\pi} f(e^{it}) \, i e^{it} \, dt . 
\end{equation*}
\begin{enumerate}[a)]
\item
Prove that for all nonnegative integers $k = 0,1,2,3,\ldots$, we have
$\int_{S^1} z^k \, dz = 0$.
\item
Prove that if
$P(z) = \sum_{k=0}^n c_k z^k$ is a
polynomial in $z$, then
$\int_{S^1} P(z) \, dz = 0$.
\item
Prove
$\int_{S^1} \bar{z} \, dz \neq 0$.
\item
Conclude that polynomials in $z$ (this algebra of functions is
not self-adjoint) are not dense in $C(S^1,\C)$.
\end{enumerate}
\end{exercise}

\begin{exercise}
Let $(X,d)$ be a compact metric space and
suppose $\sA \subset C(X,\R)$ is a real algebra that separates points, but
vanishes at exactly one point $x_0 \in X$.
That is, $f(x_0) = 0$ for all $f \in \sA$,
but for every $y \in X \setminus \{ x_0 \}$ there is
a $\varphi \in \sA$ such that $\varphi(y) \neq 0$.  Prove that
every function $g \in C(X,\R)$ such that $g(x_0) = 0$ is a uniform limit
of functions from $\sA$.
\end{exercise}

\begin{exercise}
Let $(X,d)$ be a compact metric space with at least two points and
suppose $\sA \subset C(X,\R)$ is a real algebra.
Suppose that for each $y \in X$ the closure $\widebar{\sA}$
contains the function $\varphi_y(x) \coloneqq d(y,x)$.
Show that $\widebar{\sA} = C(X,\R)$.
\end{exercise}

\begin{exercise}
\pagebreak[2]
\leavevmode
\begin{enumerate}[a)]
\item
Suppose $f \colon [a,b] \to \C$ is continuously 
differentiable.  Show that there exists a sequence of polynomials
$\{ p_n \}_{n=1}^\infty$
that converges in the $C^1$ norm to $f$, that is,
$\snorm{f - p_n}_{[a,b]} + \snorm{f'-p_n'}_{[a,b]} \to 0$ as $n \to \infty$.
\item
Suppose $f \colon [a,b] \to \C$ is $k$ times continuously 
differentiable.  Show that there exists a sequence of polynomials
$\{ p_n \}_{n=1}^\infty$
that converges in the $C^k$ norm to $f$, that is,
\begin{equation*}
\sum_{j=0}^k \norm{f^{(j)} - p_n^{(j)}}_{[a,b]} \to 0 \qquad \text{as} \qquad
n \to \infty.
\end{equation*}
\end{enumerate}
\end{exercise}

\begin{exercise}
\pagebreak[2]
\leavevmode
\begin{enumerate}[a)]
\item
Show that an even function $f \colon [-1,1] \to \R$ is a uniform
limit of polynomials with even powers only, that is, polynomials
of the form $a_0 + a_1 x^2 + a_2 x^4 + \cdots + a_k x^{2k}$.
\item
Show that an odd function $f \colon [-1,1] \to \R$ is a uniform
limit of polynomials with odd powers only, that is, polynomials
of the form $b_1 x + b_2 x^3 + b_3 x^5 + \cdots + b_k x^{2k-1}$.
\end{enumerate}
\end{exercise}

\begin{exercise} \label{exercise:finitelymanyweierequal}
\pagebreak[2]
Let $f \colon [a,b] \to \R$ be continuous.
\begin{enumerate}[a)]
\item
Given two points $x_1,x_2 \in [a,b]$, show that
there exists a sequence of real polynomials $\{ p_n \}_{n=1}^\infty$
such that $p_n(x_1) = f(x_1)$ and $p_n(x_2) = f(x_2)$ for all $n$.
\item
Generalize the previous part to $k$ points:
Given the points $x_1,x_2,\ldots,x_k \in [a,b]$, show that
there exists a sequence of real polynomials $\{ p_n \}_{n=1}^\infty$
so that for all $n$, $p_n(x_j) = f(x_j)$ for $j=1,2,\ldots,k$.
\\
Hint: The polynomial $(x-x_1)(x-x_2)\cdots(x-x_{\ell-1})(x-x_{\ell+1})
\cdots(x-x_k)$ is zero at $x_j$ for $j \neq \ell$ but nonzero at
$x_\ell$.  Use it to construct a polynomial that takes prescribed values
at $x_1,x_2,\ldots,x_k$.
\end{enumerate}
\end{exercise}

%%%%%%%%%%%%%%%%%%%%%%%%%%%%%%%%%%%%%%%%%%%%%%%%%%%%%%%%%%%%%%%%%%%%%%%%%%%%%%

\sectionnewpage
\section{Fourier series}
\label{sec:fourier}

%mbxINTROSUBSECTION

\sectionnotes{3--4 lectures}

Fourier series\footnote{%
Named after the French mathematician
\href{https://en.wikipedia.org/wiki/Joseph_Fourier}{Jean-Baptiste Joseph Fourier}
(1768--1830).} is perhaps the most important (and the most difficult)
of the series that we cover in
this book.  We saw a few examples already,
but let us start
at the beginning.

\subsection{Trigonometric polynomials}

A \emph{\myindex{trigonometric polynomial}} is an expression of the form
\begin{equation*}
a_0 + \sum_{n=1}^N \bigl(a_n \cos(nx) + b_n \sin(nx) \bigr),
\end{equation*}
or equivalently, thanks to Euler's formula ($e^{i\theta} = \cos(\theta) + i
\sin(\theta)$):
\begin{equation*}
\sum_{n=-N}^N c_n e^{inx} .
\end{equation*}
The second form is usually more convenient.  If
$z \in \C$ with $\sabs{z}=1,$ we write $z = e^{ix}$, and so
\begin{equation*}
\sum_{n=-N}^N c_n e^{inx} = 
\sum_{n=-N}^N c_n z^n .
\end{equation*}
So a trigonometric polynomial is really a rational function
of the complex variable $z$
(we are allowing negative powers) evaluated on the unit circle.  There is
a wonderful connection between power series
(actually Laurent series due to the negative powers) and
Fourier series because of this observation,
but we will not investigate this further.

\medskip

Another reason why Fourier series is important and comes up in so many
applications is that the functions $e^{inx}$ are eigenfunctions%
\footnote{An eigenfunction is like an eigenvector for a matrix, but for an
operator on a vector space of functions.} of various
differential operators.  For example,
\begin{equation*}
\frac{d}{dx} \bigl[ e^{inx} \bigr] = (in) e^{inx}, \qquad
\frac{d^2}{dx^2} \bigl[ e^{inx} \bigr] = (-n^2) e^{inx} .
\end{equation*}
That is, they are the functions whose derivative is a scalar (the
eigenvalue) times themselves.
Just as eigenvalues and eigenvectors are important in studying matrices,
eigenvalues and eigenfunctions are important when studying linear
differential equations.

\medskip

The functions $\cos (nx)$, $\sin (nx)$, and $e^{inx}$ are $2\pi$-periodic,
hence
trigonometric
polynomials are also $2\pi$-periodic.
We could rescale $x$ to make the period different, but the theory is the
same, so we stick with the period $2\pi$.
The antiderivative of $e^{inx}$ is $\frac{e^{inx}}{in}$ and
so
\begin{equation*}
\int_{-\pi}^\pi e^{inx} \, dx =
\begin{cases}
2\pi & \text{if } n=0, \\
0    & \text{otherwise.}
\end{cases}
\end{equation*}
Consider
\begin{equation*}
f(x) \coloneqq \sum_{n=-N}^N c_n e^{inx} ,
\end{equation*}
and for $m=-N,\ldots,N$ compute
\begin{equation*}
%mbxlatex \begin{aligned}
\frac{1}{2\pi} \int_{-\pi}^\pi
f(x) e^{-imx} \, dx
%mbxlatex &
=
\frac{1}{2\pi} \int_{-\pi}^\pi
\left(\sum_{n=-N}^N c_n e^{i(n-m)x}\right) \, dx
%mbxlatex \\
%mbxlatex &
=
\sum_{n=-N}^N
c_n
\frac{1}{2\pi}
\int_{-\pi}^\pi
e^{i(n-m)x}
 \, dx
=
c_m .
%mbxlatex \end{aligned}
\end{equation*}
We just found a way of computing the coefficients $c_m$ using an integral
of $f$.  If $\sabs{m} > N$, the integral is 0, so we might as
well have included enough zero coefficients to make $\sabs{m} \leq N$.

\begin{prop}
A trigonometric polynomial
$f(x) = \sum_{n=-N}^N c_n\, e^{inx}$
is real-valued for real $x$ if
and only if $c_{-m} = \overline{c_m}$ for all $m=-N,\ldots,N$.
\end{prop}

\begin{proof}
If $f(x)$ is real-valued, that is, $\overline{f(x)} = f(x)$, then
\begin{equation*}
\overline{c_m}
=
\overline{
\frac{1}{2\pi} \int_{-\pi}^\pi
f(x) e^{-imx} \, dx
}
=
\frac{1}{2\pi} \int_{-\pi}^\pi
\overline{
f(x) e^{-imx} } \, dx
=
\frac{1}{2\pi} \int_{-\pi}^\pi
f(x) e^{imx} \, dx
= c_{-m} .
\end{equation*}
The complex conjugate goes inside the integral because the integral is
done on real and imaginary parts separately.

On the other hand, if 
$c_{-m} = \overline{c_m}$, then
\begin{equation*}
\overline{c_{-m}\, e^{-imx}+ c_{m}\, e^{imx}}
=
\overline{c_{-m}}\, e^{imx}+ \overline{c_{m}}\, e^{-imx}
=
c_{m}\, e^{imx}+ c_{-m}\, e^{-imx} ,
\end{equation*}
which is real-valued.  Also $c_0 = \overline{c_0}$, so
$c_0$ is real.
By pairing up the terms, we obtain that $f$ has to be real-valued.
\end{proof}

The functions $e^{inx}$ are also linearly independent.

\begin{prop}
If
\begin{equation*}
\sum_{n=-N}^N c_n \, e^{inx} = 0
\end{equation*}
for all $x \in [-\pi,\pi]$, then $c_n = 0$ for all $n$.
\end{prop}

\begin{proof}
The result follows immediately from the integral formula for $c_n$.
\end{proof}

\subsection{Fourier series}

We now take limits.  The series
\begin{equation*}
\sum_{n=-\infty}^\infty c_n \, e^{inx}
\end{equation*}
is called
the \emph{\myindex{Fourier series}} and the numbers $c_n$
the \emph{\myindex{Fourier coefficients}}.  Using
Euler's formula $e^{i\theta} = \cos(\theta) + i \sin (\theta)$,
we could also develop everything with
sines and cosines, that is, as the series
$a_0 + \sum_{n=1}^\infty a_n \cos(nx) + b_n \sin(nx)$.
It is equivalent, but slightly messier.

Several questions arise.  What functions are expressible as 
Fourier series?  Obviously, they have to be $2\pi$-periodic, but not every
periodic function is expressible with the series.  Furthermore, if we do have
a Fourier series, where does it converge (where and if at all)?  Does it converge
absolutely?  Uniformly?  Also note that the series has two
limits.  When talking about Fourier series convergence, we often
talk about the following limit:
\begin{equation*}
\lim_{N\to\infty} 
\sum_{n=-N}^N c_n e^{inx} .
\end{equation*}
There are other ways we can sum the series to get convergence in more
situations, but we refrain from discussing those.  In light of this,
define the \emph{\myindex{symmetric partial sums}}
\glsadd{not:FSsympartsum}
\begin{equation*}
s_N(f;x) \coloneqq 
\sum_{n=-N}^N c_n \,e^{inx} .
\end{equation*}

\medskip

Conversely, for an integrable function $f \colon [-\pi,\pi] \to
\C$, call the numbers
\begin{equation*}
c_n \coloneqq 
\frac{1}{2\pi} \int_{-\pi}^\pi
f(x) e^{-inx} \, dx
\end{equation*}
its \emph{Fourier coefficients}.  To emphasize the function
the coefficients belong to,
we write $\hat{f}(n)$.\footnote{The notation should seem similar
to Fourier transform to those readers that have seen it.
The similarity is not just
coincidental, we are taking a type of Fourier transform here.}
We then formally write down a Fourier series:\glsadd{not:FS}
\begin{equation*}
f(x) \sim
\sum_{n=-\infty}^\infty c_n \, e^{inx} .
\end{equation*}
As you might imagine such a series might not even converge.
The $\sim$ doesn't imply anything about the two sides being equal
in any way.  It is simply that we created a formal series using the formula
for the coefficients.  We will see that when the functions are
\myquote{nice enough,} we do get convergence.

\begin{example}
Consider the step function $h(x)$ so that $h(x) \coloneqq 1$
on $[0,\pi]$ and $h(x) \coloneqq -1$ on $(-\pi,0)$, extended periodically to a
$2\pi$-periodic function.  With a little bit of calculus,
we compute the coefficients:
\begin{equation*}
%mbxlatex \begin{aligned}
\hat{h}(0)
%mbxlatex &
=
\frac{1}{2\pi} \int_{-\pi}^\pi h(x) \, dx = 0,
%mbxSTARTIGNORE
\qquad
%mbxENDIGNORE
%mbxlatex \\
\hat{h}(n)
%mbxlatex &
=
\frac{1}{2\pi} \int_{-\pi}^\pi h(x) e^{-inx} \, dx
=
\frac{i\bigl( (-1)^n-1 \bigr)}{\pi n} \quad \text{for } n \neq 0 .
%mbxlatex \end{aligned}
\end{equation*}
A little bit of simplification leads to
\begin{equation*}
s_N(h;x) =
\sum_{n=-N}^N \hat{h}(n) \,e^{inx} 
=
\sum_{n=1}^N \frac{2\bigl(1-(-1)^n\bigr)}{\pi n} \sin(n x) .
\avoidbreak
\end{equation*}
See the left hand graph in \figureref{fig:fourierheavicsaw}
for a graph of $h$ and several symmetric partial sums.

For a second example, consider the function $g(x) \coloneqq \sabs{x}$
on $[-\pi,\pi]$ and then extended to a $2\pi$-periodic function.
Computing the coefficients, we find
\begin{equation*}
%mbxlatex \begin{aligned}
\hat{g}(0)
%mbxlatex &
=
\frac{1}{2\pi} \int_{-\pi}^\pi g(x) \, dx = \frac{\pi}{2},
%mbxSTARTIGNORE
\qquad
%mbxENDIGNORE
%mbxlatex \\
\hat{g}(n)
%mbxlatex &
=
\frac{1}{2\pi} \int_{-\pi}^\pi g(x) e^{-inx} \, dx
=
\frac{(-1)^n-1}{\pi n^2} \quad \text{for } n \neq 0 .
%mbxlatex \end{aligned}
\end{equation*}
A little simplification yields
\begin{equation*}
s_N(g;x) =
\sum_{n=-N}^N \hat{g}(n) \,e^{inx} 
=
\frac{\pi}{2} + 
\sum_{n=1}^N \frac{2\bigl((-1)^n-1\bigr)}{\pi n^2} \cos(n x) .
\end{equation*}
See the right hand graph in \figureref{fig:fourierheavicsaw}.

\begin{myfigureht}
\myincludepdft{fourierheavi_csaw}{%
Two graphs of 2 pi periodic functions shown over an interval
slightly larger than from minus pi to pi.  On the left,
the rectangular wave is shown which starts at value 1 when x is
less than minus pi, then becomes minus 1 from minus pi to 0, then
becomes 1 from 0 to pi, then becomes minus 1 again.
On the right, a triangular (in particular continuous) wave is shown
that is an upward sloping straight line hitting pi at minus pi, then a
downward sloping straight line hitting 0 at 0, then going
up again to pi at pi, then down again.
Several partial sum approximations of the Fourier series are graphed
in lighter shades of gray
approximating the two functions.  Notably on the left the approximations
are visibly worse with lots of wiggles around the straight segments, while
on the right the approximations quickly seem to approximate only
approaching slowly near the corners of the graph.}
\caption{The functions $h$ and $g$ in bold, with several
symmetric partial sums in gray.\label{fig:fourierheavicsaw}}
\end{myfigureht}

Note that for both $h$ and $g$, the even coefficients (except $\hat{g}(0)$)
happen to vanish, but that is not really important.
What is important is convergence.
First, at the discontinuity at $x=0$, we find $s_N(h;0) = 0$
for all $N$, so $s_N(h;0)$ converges
to a different number from $h(0)$ (at a nice enough jump discontinuity, the
limit is the average of the two-sided limits, see the exercises).
That should not be surprising;
the coefficients are computed by an integral, and integration
does not notice if the value of a function changes at a single point.
We should remark, however, that we are not guaranteed that in general
the Fourier series converges to the function even at a point where
the function is continuous.
We will prove convergence if the function is at least Lipschitz.

What is really important is how fast the coefficients go to zero.  For the
discontinuous $h$, the coefficients $\hat{h}(n)$ go to zero approximately like
$\nicefrac{1}{n}$.  On the other hand, for the continuous $g$,
the coefficients $\hat{g}(n)$ go to zero approximately like
$\nicefrac{1}{n^2}$.
The Fourier coefficients \myquote{see}
the discontinuity in some sense.

Do note that continuity in this setting is the continuity of the
periodic extension, that is, we include the endpoints $\pm \pi$.
So the function $f(x) = x$ defined on $(-\pi,\pi]$
and extended periodically would be discontinuous at the endpoints $\pm\pi$.
\end{example}

In general, the relationship between regularity of the function and
the rate of decay of the coefficients is somewhat more complicated
than the example above
might make it seem, but there are some quick conclusions we can make.
We forget about finding a series for a function for a moment,
and we consider simply the limit of some given series.
A few sections ago, we proved that the Fourier series 
\begin{equation*}
\sum_{n=1}^\infty \frac{\sin(nx)}{n^2}
\end{equation*}
converges uniformly and hence converges to a continuous function.  This 
example and its proof can be extended to a more general criterion.

\begin{prop}
Let $\sum_{n=-\infty}^\infty c_n\, e^{inx}$ be a Fourier series,
and $C$, $\alpha > 1$ constants such that
\begin{equation*}
\sabs{c_n} \leq \frac{C}{\sabs{n}^\alpha}
\qquad \text{for all } n \in \Z \setminus \{ 0 \}.
\avoidbreak
\end{equation*}
Then the series converges (absolutely and uniformly) to a continuous function on $\R$.
\end{prop}

The proof is to apply the Weierstrass $M$-test (\thmref{thm:weiermtest}) and
the $p$-series test to find that the series converges uniformly and hence
to a continuous function (\corref{cor:metricuniformcontinuous}).
We can also take derivatives.

\begin{prop}
Let $\sum_{n=-\infty}^\infty c_n\, e^{inx}$ be a Fourier series,
and $C$, $\alpha > 2$ constants such that
\begin{equation*}
\sabs{c_n} \leq \frac{C}{\sabs{n}^\alpha}
\qquad \text{for all } n \in \Z \setminus \{ 0 \}.
\avoidbreak
\end{equation*}
Then the series converges to a continuously differentiable function on $\R$.
\end{prop}

The proof is to note that the series converges to a continuous
function by the previous proposition.
In particular, it converges at some point.
Then differentiate the partial sums,
\begin{equation*}
\sum_{n=-N}^{N}
i n c_n \,e^{inx}
\end{equation*}
and notice that for all nonzero $n$
\begin{equation*}
\sabs{i n c_n} \leq \frac{C}{\sabs{n}^{\alpha-1}} .
\end{equation*}
The differentiated series converges uniformly by the $M$-test again.  Since
the differentiated series
converges uniformly, we find that the original series
$\sum_{n=-\infty}^\infty c_n\,e^{inx}$
converges 
to a continuously differentiable function, whose derivative is
the differentiated series (see \thmref{thm:dersconvergecomplex}).

We can iterate this reasoning.
Suppose there is
some $C$ and $\alpha > k+1$ ($k \in \N$) such that
for all nonzero integers $n$,
\begin{equation*}
\sabs{c_n} 
\leq \frac{C}{\sabs{n}^\alpha} .
\end{equation*}
Then 
the Fourier series converges to a $k$-times continuously differentiable
function.  Therefore, the faster the coefficients go to zero, the more
regular the function is.

\subsection{Orthonormal systems}

Let us abstract away the exponentials and
study a more general series for a function.
One fundamental property of the exponentials that makes Fourier series work
is that the exponentials are what we call an \emph{orthonormal system}.
Fix an interval $[a,b]$.  We define an
\emph{\myindex{inner product}} for the space of functions.  We restrict our attention
to Riemann integrable functions as we do not have the Lebesgue
integral, which
would be the natural choice.  Let $f$ and~$g$ be complex-valued 
Riemann integrable functions on $[a,b]$ and define the inner product
\glsadd{not:L2innprod}
\begin{equation*}
\langle f , g \rangle \coloneqq
\int_a^b f(x) \overline{g(x)} \, dx .
\end{equation*}
If you have seen Hermitian inner products in linear algebra, this
is precisely such a product.  We must include the conjugate as we are
working with complex numbers.  We then have the \myquote{size} of $f$, that
is, the $L^2$ norm $\snorm{f}_2$, by (defining the square)
\glsadd{not:L2norm}
\begin{equation*}
\snorm{f}_2^2 \coloneqq
\langle f , f \rangle =
\int_a^b \babs{f(x)}^2 \, dx .
\end{equation*}

\begin{remark}
Note the similarity to finite dimensions.
For $z = (z_1,z_2,\ldots,z_d) \in \C^d$, one defines
\begin{equation*}
\langle z , w \rangle \coloneqq
\sum_{n=1}^d z_n \overline{w_n} .
\end{equation*}
Then the norm is (usually denoted simply by $\snorm{z}$ in $\C^d$
rather than by $\snorm{z}_2$)
\begin{equation*}
\snorm{z}^2 = 
\langle z , z \rangle =
\sum_{n=1}^d \sabs{z_n}^2 .
\avoidbreak
\end{equation*}
This is just the euclidean distance to the origin in $\C^d$ (same as
$\R^{2d}$).
\end{remark}

%Let us get back to function spaces.
In what follows, we will
assume all functions are Riemann integrable.

\begin{defn}
Let $\{ \varphi_n \}_{n=1}^\infty$ be a sequence of integrable complex-valued
functions on $[a,b]$.  We say that this is an
\emph{\myindex{orthonormal system}} if
\begin{equation*}
\langle \varphi_n , \varphi_m \rangle
=
\int_a^b \varphi_n(x) \, \overline{\varphi_m(x)} \, dx
= 
\begin{cases}
1 & \text{if } n=m, \\
0 & \text{otherwise.}
\end{cases}
\end{equation*}
In particular, $\snorm{\varphi_n}_2 = 1$ for all $n$.  If we
only require that 
$\langle \varphi_n , \varphi_m \rangle = 0$ for $m \neq n$, then
the system would be called an \emph{\myindex{orthogonal system}}.
\end{defn}

We noticed above that
\begin{equation*}
{\left\{ \frac{1}{\sqrt{2\pi}} \, e^{inx} \right\}}_{n=1}^\infty
\avoidbreak
\end{equation*}
is an orthonormal system on $[-\pi,\pi]$.
The factor out in front is to make the norm be 1.

Having an orthonormal system $\{ \varphi_n \}_{n=1}^\infty$ on $[a,b]$ and an integrable function $f$
on $[a,b]$, we can write
a Fourier series relative to $\{ \varphi_n \}_{n=1}^\infty$.  Let
\begin{equation*}
c_n \coloneqq
\langle f , \varphi_n \rangle
=
\int_a^b f(x) \overline{\varphi_n(x)} \, dx ,
\end{equation*}
and write
\begin{equation*}
f(x) \sim \sum_{n=1}^\infty c_n \varphi_n .
\end{equation*}

In other words, the series is
\begin{equation*}
\sum_{n=1}^\infty \langle f , \varphi_n \rangle \varphi_n(x) .
\end{equation*}
Notice the similarity to the expression for the orthogonal
projection of a vector onto a subspace from linear algebra.  We are
in fact doing just that, but in a space of functions.

\begin{thm} \label{thm:l2bestapprox}
Suppose $f$ is a Riemann integrable function on $[a,b]$.
Let $\{ \varphi_n \}_{n=1}^\infty$ be an orthonormal system on $[a,b]$ and
suppose
\begin{equation*}
f(x) \sim \sum_{n=1}^\infty c_n \varphi_n(x) .
\end{equation*}
If
\begin{equation*}
s_k (x) \coloneqq \sum_{n=1}^k c_n \varphi_n(x)
\quad\text{and}\quad
p_k (x) \coloneqq \sum_{n=1}^k d_n \varphi_n(x)
\end{equation*}
for some other sequence $\{ d_n \}_{n=1}^\infty$, then
\begin{equation*}
\int_a^b \babs{f(x)-s_k(x)}^2 \, dx = \snorm{f-s_k}_2^2 \leq
\snorm{f-p_k}_2^2 = \int_a^b \sabs{f(x)-p_k(x)}^2 \, dx
\end{equation*}
with equality only if $d_n = c_n$ for all $n=1,2,\ldots,k$.
\end{thm}

In other words, the partial sums of the Fourier series are the best approximation with respect to the
$L^2$ norm.

\begin{proof}
Let us write
\begin{equation*}
\int_a^b \sabs{f-p_k}^2
=
\int_a^b \sabs{f}^2
-
\int_a^b f \widebar{p_k}
-
\int_a^b \widebar{f} p_k
+
\int_a^b \sabs{p_k}^2 .
\end{equation*}
Now
\begin{equation*}
\int_a^b f \widebar{p_k}
=
\int_a^b f \sum_{n=1}^k \overline{d_n} \overline{\varphi_n}
=
 \sum_{n=1}^k \overline{d_n} \int_a^b f \, \overline{\varphi_n}
=
 \sum_{n=1}^k \overline{d_n} c_n ,
\end{equation*}
and
\begin{equation*}
\int_a^b \sabs{p_k}^2
=
\int_a^b
\sum_{n=1}^k d_n \varphi_n
\sum_{m=1}^k \overline{d_m} \overline{\varphi_m}
=
\sum_{n=1}^k
\sum_{m=1}^k 
d_n
\overline{d_m} 
\int_a^b
\varphi_n
\overline{\varphi_m}
=
\sum_{n=1}^k
\sabs{d_n}^2 .
\end{equation*}
So
\begin{equation*}
\begin{split}
\int_a^b \sabs{f-p_k}^2
& =
\int_a^b \sabs{f}^2
-
\sum_{n=1}^k \overline{d_n} c_n
-
\sum_{n=1}^k d_n \overline{c_n}
+
\sum_{n=1}^k
\sabs{d_n}^2
\\
& =
\int_a^b \sabs{f}^2
-
\sum_{n=1}^k \sabs{c_n}^2
+
\sum_{n=1}^k
\sabs{d_n-c_n}^2 .
\end{split}
\end{equation*}
This is minimized precisely when $d_n = c_n$.
\end{proof}

When we do plug in $d_n = c_n$, then
\begin{equation*}
\int_a^b \sabs{f-s_k}^2
=
\int_a^b \sabs{f}^2
-
\sum_{n=1}^k \sabs{c_n}^2 ,
\end{equation*}
and so for all $k$,
\begin{equation*}
\sum_{n=1}^k \sabs{c_n}^2
\leq
\int_a^b \sabs{f}^2 .
\end{equation*}
Note that
\begin{equation*}
\sum_{n=1}^k \sabs{c_n}^2 = \snorm{s_k}_2^2
\end{equation*}
by the calculation above.
We take a limit to obtain
\emph{\myindex{Bessel's inequality}}.

\begin{thm}[Bessel's inequality\footnote{%
Named after the German astronomer, mathematician, physicist, and geodesist
\href{https://en.wikipedia.org/wiki/Friedrich_Bessel}{Friedrich Wilhelm Bessel}
(1784--1846).}] \label{thm:bessels}
Suppose $f$ is a Riemann integrable function on $[a,b]$.
Let $\{ \varphi_n \}_{n=1}^\infty$ be an orthonormal system on $[a,b]$ and
suppose
\begin{equation*}
f(x) \sim \sum_{n=1}^\infty c_n \varphi_n(x) .
\end{equation*}
Then
\begin{equation*}
\sum_{n=1}^\infty \sabs{c_n}^2
\leq
\int_a^b \sabs{f}^2
= \snorm{f}_2^2 .
\end{equation*}
\end{thm}

In particular, $\int_a^b \sabs{f}^2 < \infty$ implies the series
converges and hence
\begin{equation*}
\lim_{k \to \infty} c_k = 0 .
\end{equation*}

\subsection{The Dirichlet kernel and approximate delta functions}

We return to the trigonometric Fourier series.
The system $\{ e^{inx} \}_{n=1}^\infty$ is orthogonal, but not orthonormal if we simply
integrate over $[-\pi,\pi]$.  We can rescale the integral
and hence the inner product to make 
$\{ e^{inx} \}_{n=1}^\infty$ orthonormal.  That is, if we replace
\begin{equation*}
\int_a^b \qquad \text{with} \qquad
\frac{1}{2\pi} \int_{-\pi}^\pi,
\end{equation*}
(we are just rescaling the $dx$ really)\footnote{%
Mathematicians in this field sometimes simplify matters with
a tongue-in-cheek definition that $1=2\pi$.},
then everything works and we obtain that the system
$\{ e^{inx} \}_{n=1}^\infty$
is orthonormal with respect to the inner product
\begin{equation*}
\langle f , g \rangle =
\frac{1}{2\pi} \int_{-\pi}^\pi f(x) \, \overline{g(x)} \, dx .
\end{equation*}

Suppose $f \colon \R \to \C$ is $2\pi$-periodic and integrable
on $[-\pi,\pi]$.
Write
\begin{equation*}
f(x) \sim
\sum_{n=-\infty}^\infty c_n \,e^{inx} ,
\qquad \text{where} \quad
c_n \coloneqq 
\frac{1}{2\pi} \int_{-\pi}^\pi
f(x) e^{-inx} \, dx .
\end{equation*}
Recall the notation for the symmetric partial sums,
$s_N(f;x) \coloneqq 
\sum_{n=-N}^N c_n \,e^{inx}$.
The inequality leading up to Bessel now reads:
\begin{equation*}
\frac{1}{2\pi} \int_{-\pi}^\pi
\babs{s_N(f;x)}^2 \, dx =
\sum_{n=-N}^N \sabs{c_n}^2
\leq
\frac{1}{2\pi} \int_{-\pi}^\pi
\babs{f(x)}^2
\, dx .
\end{equation*}

Let the \emph{\myindex{Dirichlet kernel}} be
\begin{equation*}
D_N(x) \coloneqq \sum_{n=-N}^N e^{inx} .
\end{equation*}
We claim that
\begin{equation*}
D_N(x) =
%\sum_{n=-N}^N e^{inx}
%=
\frac{\sin\bigl( (N+\nicefrac{1}{2})x \bigr)}{\sin(\nicefrac{x}{2})} ,
\end{equation*}
for $x$ such that $\sin(\nicefrac{x}{2}) \neq 0$.  The left-hand
side is continuous on $\R$, and hence the right-hand side extends continuously to
all of $\R$.
To show the claim,
we use a familiar trick:
\begin{equation*}
(e^{ix}-1) D_N(x) = e^{i(N+1)x} - e^{-iNx} .
\end{equation*}
Multiply by $e^{-ix/2}$
\begin{equation*}
(e^{ix/2}-e^{-ix/2}) D_N(x) = e^{i(N+\nicefrac{1}{2})x} -
e^{-i(N+\nicefrac{1}{2})x} .
\avoidbreak
\end{equation*}
The claim follows.

Expand the definition of $s_N$
\begin{multline*}
s_N(f;x) = 
\sum_{n=-N}^N \frac{1}{2\pi} \int_{-\pi}^\pi f(t) e^{-int}  \,  dt ~ e^{inx}
\\
=
\frac{1}{2\pi} \int_{-\pi}^\pi f(t) \sum_{n=-N}^N e^{in(x-t)} \, dt
=
\frac{1}{2\pi} \int_{-\pi}^\pi f(t) D_N(x-t) \, dt .
\end{multline*}
Convolution strikes again!
As $D_N$ and $f$ are $2\pi$-periodic, we may also change variables and write 
\begin{equation*}
s_N(f;x) = 
\frac{1}{2\pi} \int_{x-\pi}^{x+\pi} f(x-t) D_N(t) \, dt
=
\frac{1}{2\pi} \int_{-\pi}^\pi f(x-t) D_N(t) \, dt .
\end{equation*}
See \figureref{fig:approxdeltas} for a plot of $D_N$ for $N=5$ and $N=20$.

\begin{myfigureht}
\myincludegraphics{approxdeltas}{%
A plot of two functions for x within mius pi to pi.
In gray there is a function that starts with small-amplitude slow oscillation
on the left, becomes a sizable bump in the middle and then becomes
a small-amplitude slow oscillation on the right.
In black, it is a function with similar pattern
except the frequency is much higher and instead of a bump in the middle it
becomes a tall peak.  The oscillations go both above and below the
x-axis.}
\caption{Plot of $D_N(x)$ for $N=5$ (gray) and $N=20$
(black).\label{fig:approxdeltas}}
\end{myfigureht}

The central peak gets taller and taller as $N$ gets larger,
and the side peaks stay small.
We are convolving (again) with
\emph{approximate delta functions}\index{approximate delta function},
although these functions have
all these oscillations away from zero.  The oscillations on the side do not go away
but they are eventually so fast that we expect the integral to just sort of
cancel itself out there.
Overall, we expect that
$s_N(f)$ goes to $f$.  Things are not always simple,
but under some conditions on $f$, such a conclusion holds.  For this reason
people write
\begin{equation*}
2\pi \, \delta(x) \sim \sum_{n=-\infty}^\infty e^{inx} ,
\end{equation*}
where $\delta$ is the \myquote{delta function} (not really a function),
which is an object that will give something like
\myquote{$\int_{-\pi}^{\pi} f(x-t) \delta(t) \, dt = f(x)$.}
We can think of $D_N(x)$ converging in some sense to $2 \pi\, \delta(x)$.
However, we have not defined (and will not define) what the delta function
is, nor what it means for it to be a limit of $D_N$ or have a Fourier series.

\subsection{Localization}

If $f$ satisfies a Lipschitz condition at a point, then
the Fourier series converges at that point.

\begin{thm} \label{thm:fourierlocalization}
Let $x$ be fixed and let $f$ be a $2\pi$-periodic function
Riemann integrable on $[-\pi,\pi]$.  Suppose
there exist $\delta > 0$ and $M$ such that
\begin{equation*}
\babs{f(x+t)-f(x)} \leq M \sabs{t}
\end{equation*}
for all $t \in (-\delta,\delta)$, then
\begin{equation*}
\lim_{N \to \infty} s_N(f;x) = f(x) .
\end{equation*}
\end{thm}

In particular,
if $f$ is continuously
differentiable at $x$,
then we obtain convergence at $x$ (exercise).
%We state an often-used version of this corollary.
A function $f \colon [a,b] \to \C$ is
\emph{\myindex{continuous piecewise smooth}}\index{piecewise smooth}
if it is continuous and there exist points
$x_0 = a < x_1 < x_2 < \cdots < x_k = b$
such that for every $j$, $f$ restricted to $[x_j,x_{j+1}]$
is continuously differentiable (up to the endpoints).

\begin{cor} \label{cor:fourierpiecewisesmooth}
Let $f$ be a $2\pi$-periodic function
Riemann integrable on $[-\pi,\pi]$.  Suppose
there exist $x\in \R$ and $\delta > 0$ such that $f$ is continuous piecewise
smooth on $[x-\delta,x+\delta]$, then
\begin{equation*}
\lim_{N \to \infty} s_N(f;x) = f(x) .
\end{equation*}
\end{cor}

The proof of the corollary is left as an exercise.  Let us prove the
theorem.

\begin{proof}[Proof of \thmref{thm:fourierlocalization}]
For all $N$,
\begin{equation*}
\frac{1}{2\pi} \int_{-\pi}^\pi D_N = 1 .
\end{equation*}
Write
\begin{equation*}
\begin{split}
s_N(f;x)-f(x) & =
\frac{1}{2\pi} \int_{-\pi}^\pi f(x-t) D_N(t) \, dt 
-
f(x)
\frac{1}{2\pi} \int_{-\pi}^\pi D_N(t) \, dt
\\
& = 
\frac{1}{2\pi} \int_{-\pi}^\pi \bigl( f(x-t) - f(x) \bigr) D_N(t) \, dt 
\\
& = 
\frac{1}{2\pi} \int_{-\pi}^\pi \frac{f(x-t) - f(x)}{\sin(\nicefrac{t}{2})} \sin\bigl(
(N+\nicefrac{1}{2})t \bigr) \, dt .
\end{split}
\end{equation*}
By the hypotheses,
for small nonzero $t$,
\begin{equation*}
\abs{ \frac{f(x-t) - f(x)}{\sin(\nicefrac{t}{2})} }
\leq
\frac{M\sabs{t}}{\babs{\sin(\nicefrac{t}{2})}} .
\end{equation*}
As $\sin(\theta) = \theta + h(\theta)$ where $\frac{h(\theta)}{\theta} \to
0$ as $\theta \to 0$,
we notice that
$\frac{M\sabs{t}}{\sabs{\sin(\nicefrac{t}{2})}}$ is continuous at the
origin.
Hence,
$\frac{f(x-t) - f(x)}{\sin(\nicefrac{t}{2})}$, as a function of $t$,
is bounded near the origin.
As $t=0$ is the only place on $[-\pi,\pi]$ where the denominator vanishes,
it is the only place where there could be a problem.
So, the function is bounded near $t=0$ and clearly
Riemann integrable on any interval not including $0$, and thus
it is Riemann integrable on $[-\pi,\pi]$.
We use the trigonometric identity
\begin{equation*}
\sin\bigl( (N+\nicefrac{1}{2})t \bigr)
=
\cos(\nicefrac{t}{2}) \sin(Nt) + 
\sin(\nicefrac{t}{2}) \cos(Nt) ,
\end{equation*}
to compute
\begin{multline*}
\frac{1}{2\pi} \int_{-\pi}^\pi \frac{f(x-t) - f(x)}{\sin(\nicefrac{t}{2})} \sin\bigl(
(N+\nicefrac{1}{2})t \bigr) \, dt 
= \\
\frac{1}{2\pi} \int_{-\pi}^\pi
\left( \frac{f(x-t) - f(x)}{\sin(\nicefrac{t}{2})}
\cos (\nicefrac{t}{2}) \right) \sin (Nt) \, dt
+
\frac{1}{2\pi} \int_{-\pi}^\pi \bigl( f(x-t) - f(x) \bigr)
\cos (Nt) \, dt .
\end{multline*}
As functions of $t$, 
$\frac{f(x-t) - f(x)}{\sin(\nicefrac{t}{2})} \cos (\nicefrac{t}{2})$
and
$\bigl( f(x-t) - f(x) \bigr)$ are bounded Riemann integrable functions
and so their Fourier coefficients go to zero by \thmref{thm:bessels}.  So the two
integrals on the right-hand side, which compute the Fourier coefficients
for the real version of the Fourier series go to 0 as $N$ goes to infinity.
This is because $\sin(Nt)$ and $\cos(Nt)$ are also orthonormal systems
with respect to the same inner product.
Hence $s_N(f;x)-f(x)$ goes to 0, that is, $s_N(f;x)$ goes to $f(x)$.
\end{proof}

The theorem also says that convergence depends only on local behavior.
That is, to understand convergence of $s_N(f;x)$ we only need to know
$f$ in some neighborhood of $x$.

\begin{cor}
Suppose $f$ is a $2\pi$-periodic function, Riemann integrable on $[-\pi,\pi]$.
If $J$ is an open interval and $f(x) = 0$ for all $x \in J$,
then $\lim\limits_{N\to\infty} s_N(f;x) = 0$ for all $x \in J$.

In particular, if $f$ and $g$ are $2\pi$-periodic functions,
Riemann integrable on $[-\pi,\pi]$, $J$ an open interval, and $f(x) = g(x)$
for all $x \in J$, then for all $x \in J$,
the sequence
$\bigl\{ s_N(f;x) \bigr\}_{N=1}^\infty$ converges if and only if
$\bigl\{ s_N(g;x) \bigr\}_{N=1}^\infty$ converges.
\end{cor}

The first claim follows by taking $M=0$ in the theorem.
The \myquote{In particular} follows by considering $f-g$, which
is zero on $J$ and $s_N(f-g) = s_N(f) - s_N(g)$.
So convergence at $x$ depends only on the values of the function
near $x$.
However, we saw that the rate of convergence,
that is, how fast does $s_N(f)$
converge to $f$, depends on global behavior of $f$.

Note a subtle difference between the results above and what 
\hyperref[thm:SWcomplex]{Stone--Weierstrass theorem} gives.
%By Stone--Weierstrass, 
Any continuous function on $[-\pi,\pi]$ equal at the endpoints
can be uniformly approximated
by trigonometric polynomials, but these trigonometric polynomials may
not be the partial sums $s_N$.

\subsection{Parseval's theorem}

Finally,
convergence always happens in the $L^2$ sense and
operations on the (infinite) vectors of
Fourier coefficients are the same as the operations using the integral
inner product.

\begin{samepage}
\begin{thm}[Parseval\footnote{%
Named after the French mathematician
\href{https://en.wikipedia.org/wiki/Marc-Antoine_Parseval}{Marc-Antoine
Parseval}
(1755--1836).}] \index{Parseval's theorem}
Let $f$ and $g$ be $2\pi$-periodic functions, Riemann integrable 
on $[-\pi,\pi]$
with
\begin{equation*}
f(x) \sim
\sum_{n=-\infty}^\infty c_n \,e^{inx}
\qquad \text{and} \qquad
g(x) \sim
\sum_{n=-\infty}^\infty d_n \,e^{inx} .
\end{equation*}
Then
\begin{equation*}
\lim_{N\to\infty} \bnorm{f-s_N(f)}_2^2 = 
\lim_{N\to\infty}
\frac{1}{2\pi}
\int_{-\pi}^\pi
\babs{f(x)-s_N(f;x)}^2 \, dx
=0 .
\end{equation*}
Also
\begin{equation*}
\langle f , g \rangle =
\frac{1}{2\pi}
\int_{-\pi}^\pi
f(x) \overline{g(x)}\, dx
=
\sum_{n=-\infty}^\infty c_n \overline{d_n} ,
\end{equation*}
and
\begin{equation*}
\snorm{f}_2^2
=
\frac{1}{2\pi}
\int_{-\pi}^\pi
\babs{f(x)}^2 \, dx
=
\sum_{n=-\infty}^\infty \sabs{c_n}^2.
\end{equation*}
\end{thm}
\end{samepage}

\begin{proof}
There exists (\exerciseref{exercise:contL2close})
a continuous $2\pi$-periodic function $h$ such that
\begin{equation*}
\snorm{f-h}_2 < \epsilon .
\end{equation*}
Via \hyperref[thm:SWcomplex]{Stone--Weierstrass}
(namely \exerciseref{exercise:trigpolydense}),
there is a trigonometric polynomial $P(x)$
such that
$\babs{h(x) - P(x)} < \epsilon$ for all $x$.
Hence,
\begin{equation*}
\snorm{h-P}_2
=
\sqrt{
\frac{1}{2\pi}
\int_{-\pi}^{\pi}
\babs{h(x)-P(x)}^2
\,
dx
}
\leq \epsilon.
\end{equation*}
If $P$ is of degree $N_0$, then for all $N \geq N_0$ ,
\begin{equation*}
\bnorm{h-s_N(h)}_2 \leq \snorm{h-P}_2 \leq \epsilon ,
\end{equation*}
as $s_N(h)$ is the best approximation for $h$ in $L^2$ (\thmref{thm:l2bestapprox}).
By the inequality leading up to Bessel,
\begin{equation*}
\bnorm{s_N(h)-s_N(f)}_2
=
\bnorm{s_N(h-f)}_2
\leq
\snorm{h-f}_2 \leq \epsilon .
\end{equation*}
The $L^2$ norm satisfies the triangle inequality
(\exerciseref{exercise:L2triangleineq}).
Thus, for all $N \geq N_0$,
\begin{equation*}
\bnorm{f-s_N(f)}_2
\leq
\snorm{f-h}_2
+
\bnorm{h-s_N(h)}_2
+
\bnorm{s_N(h)-s_N(f)}_2
\leq 3\epsilon .
\end{equation*}
The first claim follows.

Next,
\begin{equation*}
\langle s_N(f) , g \rangle
=
\frac{1}{2\pi}
\int_{-\pi}^\pi
s_N(f;x) \overline{g(x)} \, dx
=
\sum_{n=-N}^N
c_n 
\frac{1}{2\pi}
\int_{-\pi}^\pi
e^{inx}
\overline{g(x)} \, dx
=
\sum_{n=-N}^N
c_n 
\overline{d_n} .
\end{equation*}
We need the Schwarz (or Cauchy--Schwarz or Cauchy--Bunyakovsky--Schwarz)
inequality for $L^2$, that is,
\begin{equation*}
{\abs{\int_a^b f\bar{g}}}^2
\leq
\left( \int_a^b \sabs{f}^2 \right)
\left( \int_a^b \sabs{g}^2 \right) .
\end{equation*}
Its proof is left as \exerciseref{exercise:L2cauchyschwarz};
it is not much
different from the finite-dimensional version.
So
\begin{equation*}
\begin{split}
\abs{\int_{-\pi}^\pi f\bar{g} - \int_{-\pi}^\pi s_N(f)\bar{g}}
& =
\abs{\int_{-\pi}^\pi \bigl(f- s_N(f)\bigr)\bar{g}} \\
%& \leq
%\int_{-\pi}^\pi \sabs{f- s_N(f)}\, \sabs{g} \\
& \leq
{\left(\int_{-\pi}^\pi \babs{f- s_N(f)}^2 \right)}^{1/2}
{\left( \int_{-\pi}^\pi \sabs{g}^2 \right)}^{1/2} .
\end{split}
\end{equation*}
The right-hand side goes to 0 as $N$ goes to infinity by the first
claim of the theorem.
That is, as $N$ goes to infinity, $\langle s_N(f),g \rangle$
goes to $\langle f,g \rangle$, and
the second claim is proved.  The last claim in the theorem follows by using
$g=f$.
\end{proof}

\subsection{Exercises}

\begin{exercise} \label{exercise:fsweierser}
Consider the Fourier series
\begin{equation*}
\sum_{k=1}^\infty \frac{1}{2^k} \sin(2^k x) .
\end{equation*}
Show that the series converges uniformly and absolutely to a continuous
function.  Remark: This is another example of a nowhere differentiable
function (you do not have to prove that)\footnote{%
See
G.\ H.\ Hardy, \emph{Weierstrass's Non-Differentiable Function},
Transactions of the American Mathematical Society,
\textbf{17}, No.\ 3 (Jul., 1916), pp.\ 301--325.
A thing to notice here is the $n$th Fourier coefficient is $\nicefrac{1}{n}$
if $n=2^k$ and zero otherwise, so the coefficients go to zero like
$\nicefrac{1}{n}$.}.
See \figureref{fig:fourierserweier}.
\begin{myfigureht}
\myincludegraphics{fourierserweier}{%
A plot of a very rough periodic function.  It seems to have lots of small
sharp peaks all over.}
\caption{Plot of 
$\sum_{n=1}^\infty \frac{1}{2^n} \sin(2^n x)$.\label{fig:fourierserweier}}
\end{myfigureht}
\end{exercise}

\begin{exercise}
Suppose that a $2\pi$-periodic function that is Riemann integrable
on $[-\pi,\pi]$, and such that $f$ is continuously differentiable
on some open interval $(a,b)$.  Prove that
for every $x \in (a,b)$, we have $\lim\limits_{N\to\infty} s_N(f;x) = f(x)$.
\end{exercise}

\begin{exercise}
Prove \corref{cor:fourierpiecewisesmooth}, that is,
suppose a $2\pi$-periodic function is continuous piecewise
smooth near a point $x$, then $\lim\limits_{N\to\infty} s_N(f;x) = f(x)$.  Hint: See the previous
exercise.
\end{exercise}

\begin{exercise} \label{exercise:contL2close}
Given a $2\pi$-periodic function $f \colon \R \to \C$, Riemann integrable on
$[-\pi,\pi]$,
and $\epsilon > 0$,
show that there exists a continuous $2\pi$-periodic function $g \colon \R
\to \C$ such that $\snorm{f-g}_2 < \epsilon$.
\end{exercise}

\begin{exercise} \label{exercise:L2cauchyschwarz}
Prove the Cauchy--Bunyakovsky--Schwarz inequality
for Riemann integrable functions:
\begin{equation*}
{\abs{\int_a^b f\bar{g}}}^2
\leq
\left( \int_a^b \sabs{f}^2 \right)
\left( \int_a^b \sabs{g}^2 \right) .
\end{equation*}
\end{exercise}

\begin{exercise} \label{exercise:L2triangleineq}
Prove the $L^2$ triangle inequality 
for Riemann integrable functions on $[-\pi,\pi]$:
\begin{equation*}
\snorm{f+g}_2 \leq \snorm{f}_2 + \snorm{g}_2 .
\end{equation*}
\end{exercise}

\begin{exercise}
\pagebreak[3]
Suppose for some $C$ and $\alpha > 1$, we have
a real sequence $\{ a_n \}_{n=1}^\infty$ with
$\sabs{a_n} \leq \frac{C}{n^\alpha}$ for all $n$.
Let
\begin{equation*}
g(x) \coloneqq \sum_{n=1}^\infty a_n \sin(n x) .
\end{equation*}
\begin{enumerate}[a)]
\item
Show that $g$ is continuous.
\item
Formally (that is, suppose you can differentiate under the sum)
find a solution (formal solution, that is, do not yet worry about convergence)
to the differential equation
\begin{equation*}
y''+ 2 y = g(x)
\end{equation*}
of the form
\begin{equation*}
y(x) = \sum_{n=1}^\infty b_n \sin(n x) .
\end{equation*}
\item
Then show that this solution $y$ is twice continuously differentiable,
and in fact solves the equation.
\end{enumerate}
\end{exercise}

\begin{exercise}
Let $f$ be a $2\pi$-periodic  function such
that $f(x) = x$ for $0 < x < 2\pi$.
Use Parseval's theorem to find
\begin{equation*}
\sum_{n=1}^\infty \frac{1}{n^2} = \frac{\pi^2}{6} .
\end{equation*}
\end{exercise}

\begin{exercise}
Suppose that $c_n = 0$ for all $n < 0$ and $\sum_{n=0}^\infty \sabs{c_n}$
converges.  Let $\D \coloneqq B(0,1) \subset \C$ be the unit disc,
and $\overline{\D} = C(0,1)$ be the closed unit disc.
Show that there exists a continuous function
$f \colon \overline{\D} \to \C$ that is analytic on $\D$
and such that on the boundary of $\D$ we have
$f(e^{i\theta}) = \sum_{n=0}^\infty c_n e^{in\theta}$.\\
Hint: If $z=re^{i\theta}$, then $z^n = r^n e^{in\theta}$.
\end{exercise}

\begin{exercise}
Show that
\begin{equation*}
\sum_{n=1}^\infty e^{-n} \sin(n x)
\end{equation*}
converges to an infinitely differentiable function.
\end{exercise}

\begin{exercise} \label{exercise:fsdiffmindecay}
Let $f$ be a $2\pi$-periodic function
such that $f(x) = f(0) + \int_0^x g$
for a function $g$ that is Riemann integrable on every interval.
Suppose
\begin{equation*}
f(x) \sim
\sum_{n=-\infty}^\infty c_n \,e^{inx} .
\end{equation*}
Show that there exists a $C > 0$ such that
$\sabs{c_n} \leq \frac{C}{\sabs{n}}$ for all nonzero $n$.
\end{exercise}

\begin{exercise}
\leavevmode
\begin{enumerate}[a)]
\item
Let $\varphi$ be the $2\pi$-periodic function 
defined by $\varphi(x) \coloneqq 0$ if $x \in (-\pi,0)$,
and $\varphi(x) \coloneqq 1$ if $x \in (0,\pi)$,
letting $\varphi(0)$ and $\varphi(\pi)$ be arbitrary.  Show
that $\lim\limits_{N \to \infty} s_N(\varphi;0) = \nicefrac{1}{2}$.
\item
Let $f$ be a $2\pi$-periodic function
Riemann integrable on $[-\pi,\pi]$, 
$x \in \R$, $\delta > 0$, and
there are continuously differentiable
$g \colon [x-\delta,x] \to \C$
and $h \colon [x,x+\delta] \to \C$
where $f(t) = g(t)$ for all $t \in [x-\delta,x)$
and
where $f(t) = h(t)$ for all $t \in (x,x+\delta]$.
Then
$\lim\limits_{N\to\infty} s_N(f;x) = \frac{g(x)+h(x)}{2}$, or in other words,
\begin{equation*}
\lim_{N \to \infty} s_N(f;x) =
\frac{1}{2}
\left(
\lim_{t \to x^-} f(t) +
\lim_{t \to x^+} f(t)
\right) .
\end{equation*}
\end{enumerate}
\end{exercise}

\begin{exercise}
Let
$\{ a_n \}_{n=1}^\infty$ be such that 
$\lim_{n\to \infty} a_n = 0$.  Show that there is a continuous
$2\pi$-periodic function $f$ whose Fourier coefficients
$c_{n}$
satisfy that for each $N$ there is a $k \geq N$ where $\sabs{c_k} \geq
a_k$.\\
Remark: The exercise says that if $f$ is only continuous, there is
no \myquote{minimum rate of decay} of the coefficients.  Compare with
\exerciseref{exercise:fsdiffmindecay}.
\\
Hint: Look at \exerciseref{exercise:fsweierser} for inspiration.
\end{exercise}
